\documentclass[11pt, a4paper, oneside]{report}
% Style of the A4 course "Architecture des ordinateurs" (Arabic),
% same palette and fonts as the ENSB slides. Compile from tex/ with
% tectonic or XeLaTeX.

\usepackage[a4paper, top=2.2cm, bottom=2.2cm, left=2cm, right=2cm,
    headheight=14pt]{geometry}
\usepackage{amsmath, amssymb}
\usepackage{booktabs, colortbl, tabularx, multirow, array}
\usepackage{tikz}
\usepackage[skins, breakable]{tcolorbox}
\usepackage{fontawesome5}
\usepackage{listings}
\usepackage{enumitem}
\usepackage{titlesec}
\usepackage{fancyhdr}
\usepackage{polyglossia}

\usetikzlibrary{positioning, calc, arrows.meta, shapes.geometric, fit,
    backgrounds, decorations.pathreplacing, matrix}

%------------------------------------------------
% Languages and fonts
%------------------------------------------------
\setdefaultlanguage[numerals=maghrib]{arabic}
\setotherlanguages{english, french}
% "الفصل" instead of polyglossia's default "باب"
\appto\captionsarabic{\renewcommand{\chaptername}{الفصل}\renewcommand{\appendixname}{الملحق}}

\newfontfamily\arabicfont[
    Script=Arabic, Scale=1.0, Path=../fonts/, Extension=.ttf,
    UprightFont=*-Regular, BoldFont=*-Bold,
    ItalicFont=*-Italic, BoldItalicFont=*-BoldItalic]{Amiri}
\newfontfamily\arabicfontsf[
    Script=Arabic, Scale=1.0, Path=../fonts/, Extension=.ttf,
    UprightFont=*-Regular, BoldFont=*-Bold,
    ItalicFont=*-Italic, BoldItalicFont=*-BoldItalic]{Amiri}
\newfontfamily\arabicfonttt[
    Script=Arabic, Scale=1.0, Path=../fonts/, Extension=.ttf,
    UprightFont=*-Regular, BoldFont=*-Bold,
    ItalicFont=*-Italic, BoldItalicFont=*-BoldItalic]{Amiri}
\newfontfamily\khatt[Script=Arabic, Scale=1.05, Path=../fonts/,
    BoldFont=ae_AlBattar.ttf]{ae_AlBattar.ttf}
% Latin monospace for code (selected directly, not through polyglossia)
\newfontfamily\codefont[Scale=0.92, BoldFont=lmmonolt10-bold.otf,
    ItalicFont=lmmono10-italic.otf]{lmmono10-regular.otf}

\linespread{1.18}
\setlength{\parindent}{0pt}
\setlength{\parskip}{5pt}

%------------------------------------------------
% Colors
%------------------------------------------------
\definecolor{brandTeal}{HTML}{5CC8E6}
\definecolor{brandOrange}{HTML}{FFA27F}
\definecolor{brandLight}{HTML}{F8F9FA}
\definecolor{brandDark}{HTML}{263238}
\definecolor{brandMuted}{HTML}{6C757D}
\definecolor{brandGreen}{HTML}{7BC47F}
\definecolor{brandPurple}{HTML}{B39DDB}
\colorlet{tealDeep}{brandTeal!60!black}
\colorlet{orangeDeep}{brandOrange!70!black}
\definecolor{codeGrey}{HTML}{F1F3F5}
\color{brandDark}

%------------------------------------------------
% Chapter and section titles
%------------------------------------------------
\titleformat{\chapter}[display]
    {\khatt\bfseries}
    {\begin{tikzpicture}
        \node[fill=brandTeal, text=white, rounded corners=6pt, inner xsep=12pt,
            inner ysep=6pt, font=\khatt\bfseries\LARGE]
            {\chaptertitlename\ \thechapter};
     \end{tikzpicture}}
    {10pt}
    {\Huge\color{brandDark}}
    [{\color{brandTeal}\titlerule[2pt]}]
\titlespacing*{\chapter}{0pt}{-10pt}{24pt}
\titleformat{\section}
    {\khatt\bfseries\Large\color{tealDeep}}{\thesection}{0.6em}{}
\titleformat{\subsection}
    {\khatt\bfseries\large\color{brandDark}}{\thesubsection}{0.6em}{}
\titleformat{\subsubsection}
    {\bfseries\color{orangeDeep}}{}{0em}{\faAngleLeft\ }

%------------------------------------------------
% Header and footer
%------------------------------------------------
\pagestyle{fancy}
\fancyhf{}
\fancyhead[R]{\small\color{brandMuted}\leftmark}
\fancyhead[L]{\small\color{brandMuted}هندسة الحواسيب}
\fancyfoot[C]{\small\color{brandMuted}\thepage}
\renewcommand{\headrulewidth}{0.4pt}
\renewcommand{\headrule}{{\color{brandTeal!60}\hrule width\headwidth height\headrulewidth}}
\renewcommand{\chaptermark}[1]{\markboth{\chaptertitlename\ \thechapter: #1}{}}
\fancypagestyle{plain}{\fancyhf{}\fancyfoot[C]{\small\color{brandMuted}\thepage}
    \renewcommand{\headrulewidth}{0pt}}

%------------------------------------------------
% Lists
%------------------------------------------------
\setlist[itemize]{label={\color{brandTeal}\small\faChevronLeft}, itemsep=2pt, topsep=3pt}
\setlist[enumerate]{label={\color{tealDeep}\bfseries\arabic*)}, itemsep=2pt, topsep=3pt}

%------------------------------------------------
% Boxes
%------------------------------------------------
\tcbset{
    base/.style={enhanced, breakable, arc=4pt, boxrule=1.2pt,
        coltitle=white, coltext=brandDark, fonttitle=\khatt\bfseries,
        toptitle=1mm, bottomtitle=1mm, top=2mm, bottom=2mm,
        left=3mm, right=3mm, before skip=8pt, after skip=8pt}}
% Definition
\newtcolorbox{tarif}[1][]{base, colback=brandTeal!8, colframe=brandTeal,
    title={\faLightbulb\quad تعريف#1}}
% Example
\newtcolorbox{mithal}[1][]{base, colback=white, colframe=brandTeal!70!black,
    title={\faPencil*\quad مثال#1}}
% Warning
\newtcolorbox{intabih}[1][]{base, colback=brandOrange!10, colframe=brandOrange,
    title={\faExclamationTriangle\quad انتبه#1}}
% Analogy from everyday life
\newtcolorbox{tashbih}[1][]{base, colback=brandPurple!10, colframe=brandPurple,
    title={\faSmile\quad تشبيه من الحياة اليومية#1}}
% Step by step
% (unbreakable: a split step box could leave an empty page)
\newtcolorbox{khatwa}[1][]{base, unbreakable, colback=brandGreen!8, colframe=brandGreen,
    title={\faShoePrints\quad خطوة بخطوة#1}}
% Summary
\newtcolorbox{khulasa}[1][]{base, colback=brandTeal!8, colframe=brandTeal,
    title={\faCheckCircle\quad خلاصة#1}}
% Exercise and solution
\newcounter{tamrin}[chapter]
\newtcolorbox{tamrin}[1][]{base, colback=white, colframe=brandMuted,
    title={\stepcounter{tamrin}\faEdit\quad تمرين \thechapter.\arabic{tamrin}#1}}
\newtcolorbox{hall}[1][]{base, colback=brandGreen!6, colframe=brandGreen!80!black,
    title={\faCheck\quad الحل#1}}
% Did you know?
\newtcolorbox{hal}[1][]{enhanced, breakable, arc=4pt, boxrule=0pt,
    colback=brandOrange!8, left=4mm, right=4mm, top=2mm, bottom=2mm,
    before skip=8pt, after skip=8pt,
    overlay={\node[anchor=north east, fill=brandOrange, text=white,
        rounded corners=3pt, font=\small\khatt\bfseries, inner sep=3pt]
        at ([xshift=-3mm, yshift=2mm]frame.north east) {\faQuestionCircle\ هل تعلم؟};},
    top=4mm}

%------------------------------------------------
% Terms: Arabic (French / English)
%------------------------------------------------
\newcommand{\mustalah}[3]{\textbf{#1} {\small\color{brandMuted}(\textenglish{#2 / #3})}}
\newcommand{\fr}[1]{\textenglish{#1}}
% Register transfer "dest <- src" as ONE left-to-right unit, so that the
% arrow always points to the destination inside Arabic (RTL) text
\newcommand{\nql}[2]{\textenglish{#1\,$\leftarrow$\,#2}}
\newcommand{\code}[1]{\mbox{\textenglish{\codefont\begingroup\color{tealDeep}#1\endgroup}}}
\newcommand{\hex}[1]{\mbox{\textenglish{#1}\textsubscript{\textenglish{H}}}}
\newcommand{\bin}[1]{\mbox{\textenglish{#1}\textsubscript{\textenglish{2}}}}

%------------------------------------------------
% Tables
%------------------------------------------------
\newcolumntype{R}{>{\raggedleft\arraybackslash}X}
\newcolumntype{C}{>{\centering\arraybackslash}X}
\newcommand{\ras}[1]{\cellcolor{brandTeal!25}\textbf{#1}}
\newcommand{\cell}[2]{\leavevmode\textcolor{#1}{#2}}
\renewcommand{\arraystretch}{1.3}
\setlength{\arrayrulewidth}{0.5pt}
\arrayrulecolor{brandMuted!60}

%------------------------------------------------
% Code listings (always inside the english environment)
%------------------------------------------------
\lstdefinelanguage{asm86}{
    morekeywords={mov, add, adc, sub, sbb, inc, dec, neg, mul, imul, div, idiv,
        cmp, and, or, xor, not, test, shl, shr, sal, sar, rol, ror, rcl, rcr,
        jmp, je, jne, jz, jnz, jg, jge, jl, jle, ja, jae, jb, jbe, jc, jnc, jo,
        js, jns, jcxz, loop, loope, loopne, call, ret, iret, int, push, pop,
        pushf, popf, lea, xchg, in, out, nop, hlt, cbw, cwd, movsb, movsw,
        lodsb, stosb, rep, cld, std, cli, sti, clc, stc, xlat, syscall,
        movzx, movsx, cdq, cqo, addps, mulps, movaps, movups, vaddps, vmovups,
        addss, movss},
    morekeywords=[2]{org, db, dw, dd, dq, equ, proc, endp, offset, ptr, byte,
        word, dword, qword, section, global, extern, dup, resb, resq, rel},
    sensitive=false, morecomment=[l]{;}, morestring=[b]', morestring=[b]"}
\lstdefinestyle{base}{
    basicstyle=\codefont\small, columns=fullflexible, keepspaces=true,
    showstringspaces=false, upquote=true,
    keywordstyle=\color{tealDeep}\bfseries,
    keywordstyle=[2]\color{brandPurple!60!black}\bfseries,
    commentstyle=\color{red!75!black},
    stringstyle=\color{orangeDeep},
    numbers=left, numberstyle=\scriptsize\color{brandMuted}, numbersep=8pt,
    xleftmargin=18pt, aboveskip=0pt, belowskip=0pt}
\lstdefinestyle{asm}{style=base, language=asm86}
\lstdefinestyle{c}{style=base, language=C, morekeywords={uint8_t, int64_t}}
\lstdefinestyle{out}{basicstyle=\codefont\small\color{brandDark}, columns=fullflexible,
    keepspaces=true, aboveskip=0pt, belowskip=0pt}

\newtcolorbox{codebox}[1][]{enhanced, colback=codeGrey, colframe=codeGrey,
    boxrule=0pt, arc=2pt, left=4pt, right=4pt, top=3pt, bottom=3pt,
    before skip=6pt, after skip=6pt,
    overlay={\fill[brandTeal!70] (frame.south west) rectangle ([xshift=3pt]frame.north west);}, #1}
\newtcolorbox{outbox}[1][]{enhanced, colback=brandDark, colframe=brandDark, boxrule=0pt,
    arc=2pt, left=6pt, right=6pt, top=3pt, bottom=3pt, before skip=4pt, after skip=8pt,
    colupper=white, #1}
% \asm{ch3/hello.asm}  \cfile{ch4/sum.c}  \screen{ch3/hello.out}
\newcommand{\asm}[2][]{\begin{codebox}\begin{english}\lstinputlisting[style=asm,#1]{../code/#2}\end{english}\end{codebox}}
\newcommand{\cfile}[2][]{\begin{codebox}\begin{english}\lstinputlisting[style=c,#1]{../code/#2}\end{english}\end{codebox}}
\newcommand{\screen}[2][]{\begin{outbox}{\scriptsize\faDesktop\ \textenglish{\codefont screen}}\par\begin{english}\lstinputlisting[style=out, basicstyle=\codefont\small\color{white},#1]{../code/#2}\end{english}\end{outbox}}

% Line-by-line explanation:  \begin{sharh} \s{1}{...} \s{2-3}{...} \end{sharh}
% (tabularx must be called as \tabularx ... \endtabularx inside an environment)
\newenvironment{sharh}
    {\par\small\tabularx{\linewidth}{|>{\centering\arraybackslash}p{1.3cm}|R|}
     \hline \ras{السطر} & \ras{ماذا يفعل؟} \\ \hline}
    {\endtabularx\par}
\newcommand{\s}[2]{\textenglish{#1} & #2 \\ \hline}

%------------------------------------------------
% TikZ: right-to-left text inside nodes
%------------------------------------------------
\newcommand{\nodecolor}{}
\tikzset{
    every node/.append style={
        execute at begin node={\nodecolor\everypar{\beginR}\beginR},
        execute at end node=\endR},
    lawn/.code={\def\nodecolor{\color{#1}}},
    % labels placed on arrow paths are invisible without an explicit text
    % colour (RTL node hook), hence nodes={lawn=brandDark}
    sahm/.style={-{Stealth[length=2.6mm]}, thick, color=brandMuted, nodes={lawn=brandDark}},
    sahmT/.style={-{Stealth[length=2.6mm]}, very thick, color=tealDeep, nodes={lawn=brandDark}},
    sahmO/.style={-{Stealth[length=2.6mm]}, very thick, color=orangeDeep, nodes={lawn=brandDark}},
    unit/.style={draw=#1, fill=#1!15, very thick, rounded corners=4pt,
        align=center, font=\bfseries\small, inner sep=5pt},
    unit/.default=brandTeal,
    reg/.style={draw=brandDark, fill=white, thick, minimum height=0.65cm,
        font=\codefont\small, inner sep=2pt},
}


\begin{document}

%================================================
% Cover
%================================================
\begin{titlepage}
\thispagestyle{empty}
\begin{tikzpicture}[remember picture, overlay]
    \fill[brandTeal!12] (current page.north west) rectangle ([yshift=-8cm]current page.north east);
    \draw[brandTeal, line width=1.5pt, opacity=0.35]
        ([xshift=-3cm, yshift=4cm]current page.south east) circle (6cm);
    \draw[brandOrange, line width=1pt, opacity=0.35]
        ([xshift=-3cm, yshift=4cm]current page.south east) circle (7cm);
    % a small chip drawing
    \begin{scope}[shift={([xshift=4.5cm, yshift=-4.2cm]current page.north west)}, scale=0.9]
        \fill[brandDark] (-1.5,-1.5) rectangle (1.5,1.5);
        \fill[brandTeal] (-1,-1) rectangle (1,1);
        \node[font=\codefont\bfseries\Large, lawn=white] at (0,0) {\textenglish{CPU}};
        \foreach \i in {-1.2,-0.6,0,0.6,1.2} {
            \draw[brandMuted, line width=2pt] (\i,1.5) -- (\i,2);
            \draw[brandMuted, line width=2pt] (\i,-1.5) -- (\i,-2);
            \draw[brandMuted, line width=2pt] (1.5,\i) -- (2,\i);
            \draw[brandMuted, line width=2pt] (-1.5,\i) -- (-2,\i);
        }
    \end{scope}
\end{tikzpicture}

\vspace*{1.2cm}
\begin{flushright}
    {\khatt\bfseries\color{tealDeep}\Large مقياس \textenglish{UEF31}}\\[0.6cm]
    {\khatt\bfseries\color{brandDark}\fontsize{38}{46}\selectfont هندسة الحواسيب}\\[0.4cm]
    {\large\color{brandMuted}\textenglish{Architecture des ordinateurs}}
\end{flushright}

\vspace{4.5cm}
\begin{tcolorbox}[enhanced, colback=white, colframe=brandTeal, arc=6pt, boxrule=1.2pt,
    width=0.85\linewidth]
    \large
    \begin{itemize}
        \item الفصل الأول: التنظيم العام للوحدة المركزية للحاسوب
        \item الفصل الثاني: البنية الداخلية للمعالج
        \item الفصل الثالث: دراسة حالة، المعالج \textenglish{Intel 8086}
        \item الفصل الرابع: المعالجات ذات \textenglish{32} و\textenglish{64} بت
        \item الفصل الخامس: بنية المعالجات الحديثة
    \end{itemize}
\end{tcolorbox}

\vfill
\begin{flushright}
    \begin{tabular}{@{}ll}
        \textbf{الأستاذ:} & د. بن عبد الرزاق يوسف \\
        \textbf{المعامل:} & \textenglish{3} \quad\textbf{الأرصدة:} \textenglish{5} \\
        \textbf{التقييم:} & امتحان \textenglish{60\%}، مراقبة مستمرة \textenglish{40\%} \\
    \end{tabular}
\end{flushright}
\end{titlepage}

%================================================
% How to use this course
%================================================
\chapter*{كيف تستعمل هذا الدرس؟}
\addcontentsline{toc}{chapter}{كيف تستعمل هذا الدرس؟}
\markboth{كيف تستعمل هذا الدرس؟}{}

هذا الدرس مكتوب لك أنت، حتى لو كانت هذه أول مرة تفتح فيها «صندوق» الحاسوب.
نبدأ من الأشياء البسيطة، ونصعد خطوة خطوة حتى نكتب برامج حقيقية بلغة التجميع
\textenglish{(Assembly)}، ثم نفهم كيف تعمل المعالجات الحديثة في هواتفنا وحواسيبنا.

\section*{الصناديق الملونة}
ستجد في كل صفحة صناديق ملونة، لكل لون دور:

\begin{tarif}
    كلمة جديدة مهمة، مع اسمها بالفرنسية والإنجليزية لأنك ستجدها في الكتب والمراجع.
\end{tarif}
\begin{tashbih}
    نشرح الفكرة الصعبة بشيء تعرفه من حياتك اليومية: المطبخ، المدرسة، الحافلة\dots
\end{tashbih}
\begin{khatwa}
    حساب أو تنفيذ نقوم به مرحلة بمرحلة، دون أن نقفز أي خطوة.
\end{khatwa}
\begin{intabih}
    خطأ يقع فيه كثير من الطلبة، أو تفصيل صغير يغيّر النتيجة.
\end{intabih}
\begin{hal}
    معلومة إضافية ممتعة، ليست ضرورية للامتحان لكنها تساعدك على الفهم.
\end{hal}

\section*{البرامج في هذا الدرس}
\begin{itemize}
    \item كل برنامج مكتوب كاملًا، ومرفوق بجدول \textbf{«ماذا يفعل كل سطر؟»}.
    \item برامج المعالج \textenglish{8086} مكتوبة لبرنامج المحاكاة \textenglish{EMU8086}،
          وبرامج \textenglish{64} بت مكتوبة لـ\textenglish{NASM} على نظام \textenglish{Linux}.
    \item كل البرامج \textbf{جُرّبت فعلًا}: ما تراه في الشاشة السوداء \faDesktop{} هو ما يظهر عند التنفيذ،
          وجداول السجلات مأخوذة من التنفيذ الحقيقي.
\end{itemize}

\begin{intabih}
    لا تكتفِ بالقراءة! انسخ كل برنامج في \textenglish{EMU8086}، ونفّذه تعليمة بتعليمة
    (زر \textenglish{Single Step})، وراقب تغيّر السجلات. هكذا يُفهم هذا المقياس.
\end{intabih}

\tableofcontents

\chapter{التنظيم العام للوحدة المركزية للحاسوب}

\begin{khulasa}[: ماذا سنتعلم في هذا الفصل؟]
\begin{itemize}
    \item ما الحاسوب، وكيف تطوّر عبر التاريخ.
    \item المكونات الأساسية: المعالج، الذاكرة، وحدات الإدخال والإخراج.
    \item نموذجا \textenglish{Von Neumann} و\textenglish{Harvard}، والفرق بينهما.
    \item أنواع الذاكرة، وكيف نحسب سعتها.
    \item الحافلات الثلاث: حافلة المعطيات، حافلة العناوين، حافلة التحكم.
\end{itemize}
\end{khulasa}

%================================================
\section{ما الحاسوب؟}
%================================================

\begin{tarif}
    \mustalah{الحاسوب}{Ordinateur}{Computer}: آلة إلكترونية تستقبل \textbf{معطيات}،
    وتعالجها حسب \textbf{برنامج} (قائمة من التعليمات)، ثم تعطي \textbf{نتائج}، ويمكنها أن تحفظ
    المعطيات والنتائج.
\end{tarif}

كل ما يفعله الحاسوب يمكن تلخيصه في أربعة أفعال: \textbf{يستقبل}، \textbf{يعالج}،
\textbf{يحفظ}، \textbf{يُخرج}.

لنأخذ مثالًا بسيطًا: تستعمل تطبيق الآلة الحاسبة في حاسوبك لتحسب $25 \times 4$.
\begin{enumerate}
    \item \textbf{الإدخال}: تكتب \textenglish{25}، ثم \textenglish{×}، ثم \textenglish{4} بلوحة المفاتيح.
          كل ضغطة تتحول إلى رمز رقمي يصل إلى الحاسوب.
    \item \textbf{الحفظ المؤقت}: يحفظ الحاسوب العددين في الذاكرة، لأنه يحتاج إليهما بعد قليل.
    \item \textbf{المعالجة}: يقرأ المعالج العددين، وينفذ تعليمة الضرب، فيحصل على \textenglish{100}.
    \item \textbf{الإخراج}: يرسل الحاسوب النتيجة إلى الشاشة فتراها.
\end{enumerate}
كل برنامج تستعمله، مهما كان معقدًا (لعبة، متصفح، تطبيق صور)، هو في الحقيقة ملايين من هذه
الخطوات الصغيرة، تُنفَّذ بسرعة كبيرة جدًا.

\begin{center}
\begin{tikzpicture}[node distance=1.4cm]
    \node[unit=brandGreen, minimum width=2.8cm, minimum height=1.2cm] (in) {الإدخال\\{\small لوحة المفاتيح، الفأرة}};
    \node[unit, minimum width=2.8cm, minimum height=1.2cm, left=of in] (cpu) {المعالجة\\{\small المعالج}};
    \node[unit=brandOrange, minimum width=2.8cm, minimum height=1.2cm, left=of cpu] (out) {الإخراج\\{\small الشاشة، الطابعة}};
    \node[unit=brandPurple, minimum width=2.8cm, minimum height=1.1cm, below=1cm of cpu] (mem) {الحفظ\\{\small الذاكرة، القرص}};
    \draw[sahmT] (in) -- (cpu);
    \draw[sahmT] (cpu) -- (out);
    \draw[sahmT, <->] (cpu) -- (mem);
\end{tikzpicture}
\end{center}

\begin{tashbih}
    الحاسوب يشبه \textbf{مطبخ مطعم}: المكونات (الخضر، اللحم) هي \textbf{المعطيات}، والوصفة هي
    \textbf{البرنامج}، والطباخ هو \textbf{المعالج}، والثلاجة هي \textbf{الذاكرة}، والطبق الجاهز
    هو \textbf{النتيجة}. الطباخ لا يقرر وحده: هو ينفذ الوصفة خطوة بخطوة.
\end{tashbih}

\subsection{العتاد والبرمجيات}
\begin{itemize}
    \item \mustalah{العتاد}{Matériel}{Hardware}: كل ما يمكن لمسه: المعالج، الذاكرة، الشاشة، الأسلاك.
    \item \mustalah{البرمجيات}{Logiciel}{Software}: البرامج التي تقول للعتاد ماذا يفعل:
          نظام التشغيل، المتصفح، الألعاب.
\end{itemize}
في هذا المقياس ندرس \textbf{العتاد من الداخل}، ونرى كيف ينفذ البرامج.

\subsection{طبقات الحاسوب}
يمكن أن نتخيل الحاسوب طبقات بعضها فوق بعض. كل طبقة تستعمل خدمات الطبقة التي تحتها، ولا تحتاج
أن تعرف تفاصيلها:

\begin{center}
\begin{tikzpicture}[l/.style={draw=#1, fill=#1!15, very thick, minimum width=9cm, minimum height=0.75cm, font=\small}]
    \node[l=brandGreen] at (0,1.8) {التطبيقات: المتصفح، الألعاب، معالج النصوص};
    \node[l=brandPurple] at (0,0.9) {نظام التشغيل: \textenglish{Windows}، \textenglish{Linux}، \textenglish{Android}};
    \node[l=brandOrange] at (0,0) {لغة الآلة ولغة التجميع: التعليمات التي يفهمها المعالج};
    \node[l=brandTeal] at (0,-0.9) {العتاد: المعالج، الذاكرة، الحافلات، الدارات الإلكترونية};
    \draw[sahmT] (5,-1.1) -- (5,2) node[midway, right, font=\small] {من الآلة إلى الإنسان};
\end{tikzpicture}
\end{center}

في هذا الدرس نشتغل على \textbf{الطبقتين السفليين}: العتاد، ولغة التجميع التي تتكلم معه مباشرة.

\subsection{كيف يعمل برنامج؟}
عندما تضغط مرتين على أيقونة برنامج، تحدث هذه المراحل:
\begin{enumerate}
    \item البرنامج محفوظ على \textbf{القرص} (\textenglish{SSD} أو \textenglish{HDD}) في شكل ملف من الأصفار والآحاد.
    \item نظام التشغيل \textbf{ينسخه في الذاكرة المركزية} \textenglish{(RAM)}، لأن القرص بطيء جدًا والمعالج لا يقرأ منه مباشرة.
    \item المعالج يقرأ التعليمات من الذاكرة \textbf{واحدة بعد الأخرى} وينفذها.
    \item عندما تغلق البرنامج، تُحرَّر مساحته في الذاكرة لبرامج أخرى.
\end{enumerate}

%================================================
\section{تاريخ الحواسيب باختصار}
%================================================

لم يظهر الحاسوب فجأة. هذه أهم المحطات:

\begin{center}
\begin{tikzpicture}[
    yr/.style={circle, fill=brandTeal, minimum size=0.35cm, inner sep=0pt},
    txt/.style={text width=2.7cm, align=center, font=\small}]
    \draw[brandTeal, line width=3pt] (0,0) -- (-14,0);
    \foreach \x/\y/\t/\p in {
        0/1642/{آلة \textenglish{Pascal}\\ للجمع والطرح}/above,
        -2/1837/{آلة \textenglish{Babbage}\\ فكرة البرنامج}/below,
        -4/1945/{\textenglish{ENIAC}\\ مصابيح مفرغة}/above,
        -6/1947/{اختراع\\ \textbf{الترانزستور}}/below,
        -8/1958/{\textbf{الدارة المندمجة}\\ \textenglish{(IC)}}/above,
        -10/1971/{أول معالج دقيق\\ \textenglish{Intel 4004}}/below,
        -12/1978/{المعالج\\ \textenglish{Intel 8086}}/above,
        -14/2005/{المعالجات\\ \textbf{متعددة الأنوية}}/below} {
        \node[yr] (n) at (\x, 0) {};
        \node[txt, \p=0.25cm of n] {\textbf{\textenglish{\y}}\\ \t};
    }
\end{tikzpicture}
\end{center}

\begin{table}[h]
\centering
\begin{tabular}{|c|c|c|c|}
    \hline
    \ras{الجيل} & \ras{الفترة} & \ras{التقنية الأساسية} & \ras{مثال} \\ \hline
    الأول   & \textenglish{1945--1955} & المصابيح المفرغة & \textenglish{ENIAC} (\textenglish{30} طنًا!) \\ \hline
    الثاني  & \textenglish{1955--1965} & الترانزستور & \textenglish{IBM 7090} \\ \hline
    الثالث  & \textenglish{1965--1980} & الدارات المندمجة & \textenglish{IBM 360} \\ \hline
    الرابع  & \textenglish{1980--} الآن & المعالج الدقيق (\textenglish{VLSI}) & الحاسوب الشخصي، الهاتف \\ \hline
\end{tabular}
\end{table}

كل جيل جديد جاء بتقنية تجعل المكونات \textbf{أصغر} و\textbf{أسرع} و\textbf{أقل استهلاكًا للطاقة}:
\begin{itemize}
    \item \textbf{المصابيح المفرغة} كانت كبيرة، تسخن كثيرًا، وتحترق باستمرار؛ كان يجب تبديل بعضها كل يوم تقريبًا.
    \item \textbf{الترانزستور} يقوم بالعمل نفسه (مفتاح كهربائي يفتح ويغلق)، لكنه أصغر بمئات المرات ولا يحترق.
    \item \textbf{الدارة المندمجة} تضع آلاف الترانزستورات على قطعة صغيرة واحدة من السيليسيوم، بدل توصيلها بالأسلاك.
    \item \textbf{المعالج الدقيق} يضع المعالج \textbf{كله} في شريحة واحدة. من هنا ولد الحاسوب الشخصي ثم الهاتف الذكي.
\end{itemize}

\begin{tarif}
    \mustalah{قانون مور}{Loi de Moore}{Moore's law}: لاحظ \textenglish{Gordon Moore} (أحد مؤسسي \textenglish{Intel})
    سنة 1965 أن عدد الترانزستورات في الشريحة \textbf{يتضاعف تقريبًا كل سنتين}. بقيت هذه الملاحظة صحيحة
    لأكثر من خمسين سنة: من \textenglish{2300} ترانزستور في \textenglish{4004} إلى عشرات المليارات اليوم.
\end{tarif}

\begin{hal}
    كان \textenglish{ENIAC} يملأ قاعة كاملة ويستهلك \textenglish{150} كيلوواط، بينما هاتفك
    أقوى منه بملايين المرات ويعمل ببطارية صغيرة. السبب: الترانزستورات أصبحت صغيرة جدًا،
    حتى صار المعالج الواحد يحتوي على مليارات منها.
\end{hal}

%================================================
\section{تذكير سريع: كيف يعدّ الحاسوب؟}
%================================================

الحاسوب لا يعرف إلا حالتين كهربائيتين: \textbf{يمر التيار} أو \textbf{لا يمر}. نكتبهما
\textenglish{1} و\textenglish{0}.

\begin{tarif}
    \begin{itemize}
        \item \mustalah{البت}{Bit}{Bit}: أصغر وحدة معلومة، قيمته \textenglish{0} أو \textenglish{1}.
        \item \mustalah{البايت}{Octet}{Byte}: مجموعة من \textbf{8 بتات}.
        \item \mustalah{الكلمة}{Mot}{Word}: عدد البتات التي يعالجها المعالج مرة واحدة
              (\textenglish{16} بتًا في المعالج \textenglish{8086}).
    \end{itemize}
\end{tarif}

نكتب الأعداد في ثلاثة أنظمة: العشري (الأساس 10)، الثنائي (الأساس 2)، والست عشري (الأساس 16،
أرقامه من \textenglish{0} إلى \textenglish{9} ثم \textenglish{A} إلى \textenglish{F}).

\begin{center}
\small
\begin{tabular}{|c|c|c||c|c|c|}
    \hline
    \ras{عشري} & \ras{ثنائي} & \ras{ست عشري} & \ras{عشري} & \ras{ثنائي} & \ras{ست عشري} \\ \hline
    \textenglish{0} & \textenglish{0000} & \textenglish{0} & \textenglish{8}  & \textenglish{1000} & \textenglish{8} \\ \hline
    \textenglish{1} & \textenglish{0001} & \textenglish{1} & \textenglish{9}  & \textenglish{1001} & \textenglish{9} \\ \hline
    \textenglish{2} & \textenglish{0010} & \textenglish{2} & \textenglish{10} & \textenglish{1010} & \textenglish{A} \\ \hline
    \textenglish{3} & \textenglish{0011} & \textenglish{3} & \textenglish{11} & \textenglish{1011} & \textenglish{B} \\ \hline
    \textenglish{4} & \textenglish{0100} & \textenglish{4} & \textenglish{12} & \textenglish{1100} & \textenglish{C} \\ \hline
    \textenglish{5} & \textenglish{0101} & \textenglish{5} & \textenglish{13} & \textenglish{1101} & \textenglish{D} \\ \hline
    \textenglish{6} & \textenglish{0110} & \textenglish{6} & \textenglish{14} & \textenglish{1110} & \textenglish{E} \\ \hline
    \textenglish{7} & \textenglish{0111} & \textenglish{7} & \textenglish{15} & \textenglish{1111} & \textenglish{F} \\ \hline
\end{tabular}
\end{center}

\begin{khatwa}[: تحويل \textenglish{2F}\textsubscript{\textenglish{H}} إلى الثنائي وإلى العشري]
    \begin{enumerate}
        \item كل رقم ست عشري يساوي \textbf{4 بتات}: \textenglish{2 = 0010} و\textenglish{F = 1111}.
        \item نضعهما جنبًا إلى جنب: $\hex{2F} = \bin{0010\,1111}$.
        \item إلى العشري: $2 \times 16 + 15 = 32 + 15 = 47$.
    \end{enumerate}
\end{khatwa}

\begin{khatwa}[: تحويل العدد الثنائي \textenglish{1011} إلى العشري]
    كل بت له «وزن» هو قوة من قوى 2، تبدأ من اليمين بـ$2^0 = 1$:
    \begin{center}
    \begin{tabular}{|c|c|c|c|c|}
        \hline \ras{البت} & 1 & 0 & 1 & 1 \\ \hline
        \ras{الوزن} & $2^3 = 8$ & $2^2 = 4$ & $2^1 = 2$ & $2^0 = 1$ \\ \hline
        \ras{البت × الوزن} & 8 & 0 & 2 & 1 \\ \hline
    \end{tabular}
    \end{center}
    نجمع: $8 + 0 + 2 + 1 = 11$. إذن $\bin{1011} = 11$.
\end{khatwa}

\begin{khatwa}[: تحويل العدد العشري 13 إلى الثنائي (القسمة المتتالية على 2)]
    \begin{enumerate}
        \item $13 \div 2 = 6$، والباقي \textbf{1}.
        \item $6 \div 2 = 3$، والباقي \textbf{0}.
        \item $3 \div 2 = 1$، والباقي \textbf{1}.
        \item $1 \div 2 = 0$، والباقي \textbf{1}. نتوقف لأن الحاصل أصبح 0.
        \item نقرأ البواقي \textbf{من الأخير إلى الأول}: \bin{1101}. تحقق: $8 + 4 + 0 + 1 = 13$.
    \end{enumerate}
\end{khatwa}

\begin{tashbih}
    لماذا نستعمل الست عشري؟ لأن العدد الثنائي طويل جدًا ومتعب للعين: \bin{1010\,1111\,0011\,0110}.
    الست عشري يختصر كل 4 بتات في رقم واحد، فيصبح العدد \hex{AF36}. هو «الاختزال» الذي يستعمله المبرمجون
    لكتابة الثنائي بسرعة، كما نكتب «د.» بدل «دكتور».
\end{tashbih}

\begin{intabih}
    في لغة التجميع نكتب حرف \code{h} بعد العدد الست عشري (\code{2Fh})، وحرف \code{b} بعد العدد
    الثنائي (\code{00101111b}). العدد دون حرف هو عدد عشري.
\end{intabih}

\subsection{وحدات قياس الذاكرة}
\begin{center}
\begin{tabular}{|c|c|c|}
    \hline
    \ras{الوحدة} & \ras{الرمز} & \ras{القيمة} \\ \hline
    كيلوبايت  & \textenglish{KB} & $2^{10} = 1024$ بايت \\ \hline
    ميغابايت  & \textenglish{MB} & $2^{20} = 1\,048\,576$ بايت \\ \hline
    غيغابايت  & \textenglish{GB} & $2^{30}$ بايت (حوالي مليار) \\ \hline
    تيرابايت  & \textenglish{TB} & $2^{40}$ بايت \\ \hline
\end{tabular}
\end{center}

%================================================
\section{المكونات الأساسية للحاسوب}
%================================================

داخل كل حاسوب، من أكبر خادم إلى أصغر ساعة ذكية، نجد ثلاثة أجزاء أساسية:

\begin{center}
\begin{tikzpicture}
    \node[unit, minimum width=4cm, minimum height=2.6cm] (cpu) at (0,0) {};
    \node[font=\bfseries] at (0, 1) {المعالج \textenglish{(CPU)}};
    \node[unit=brandOrange, minimum width=1.6cm, font=\small] at (-0.95, 0) {\textenglish{UC}\\وحدة التحكم};
    \node[unit=brandGreen, minimum width=1.6cm, font=\small] at (0.95, 0) {\textenglish{UAL}\\الحساب};
    \node[draw=brandDark, fill=white, font=\small, minimum width=3.4cm] at (0,-0.9) {السجلات};
    \node[unit=brandPurple, minimum width=3.4cm, minimum height=2.6cm] (mem) at (-6,0) {الذاكرة المركزية\\{\small \textenglish{RAM} و\textenglish{ROM}}};
    \node[unit=brandGreen, minimum width=3.4cm, minimum height=2.6cm] (io) at (6,0) {وحدات\\الإدخال والإخراج};
    \draw[sahmT, <->] (cpu) -- (mem);
    \draw[sahmT, <->] (cpu) -- (io);
\end{tikzpicture}
\end{center}

\begin{enumerate}
    \item \mustalah{المعالج}{Processeur}{Processor, CPU}: «دماغ» الحاسوب. يقرأ التعليمات وينفذها.
    \item \mustalah{الذاكرة المركزية}{Mémoire centrale}{Main memory}: تحفظ البرنامج
          الذي يعمل الآن ومعطياته.
    \item \mustalah{وحدات الإدخال والإخراج}{Périphériques d'E/S}{I/O devices}: تسمح
          للحاسوب بالتواصل مع العالم الخارجي.
\end{enumerate}

\begin{khatwa}[: رحلة حرف من لوحة المفاتيح إلى الشاشة]
    تضغط على المفتاح \textenglish{A} في برنامج كتابة:
    \begin{enumerate}
        \item لوحة المفاتيح (\textbf{إدخال}) ترسل رمز المفتاح إلى الحاسوب.
        \item المعالج يُقاطَع (يتوقف لحظة عن عمله)، فيقرأ الرمز، ويحفظه في \textbf{الذاكرة} داخل النص الذي تكتبه.
        \item برنامج الكتابة يطلب من المعالج «رسم» الحرف: يحسب المعالج النقاط الضوئية التي تشكّل الحرف.
        \item هذه النقاط تُرسل إلى بطاقة الشاشة، ثم إلى الشاشة (\textbf{إخراج}).
    \end{enumerate}
    كل هذا يحدث في جزء صغير جدًا من الثانية، فتشعر أن الحرف ظهر فورًا.
\end{khatwa}

\begin{tarif}
    \mustalah{الوحدة المركزية}{Unité centrale}{Central unit}: المعالج مع الذاكرة المركزية.
    هي الجزء الذي ينفذ البرامج فعلًا.
\end{tarif}

%================================================
\section{النموذجان الأساسيان: فون نيومان وهارفارد}
%================================================

\subsection{نموذج فون نيومان \textenglish{(Von Neumann)}}

في سنة 1945 اقترح العالم \textenglish{John von Neumann} فكرة بسيطة وعبقرية: نحفظ
\textbf{البرنامج} و\textbf{المعطيات} معًا في \textbf{الذاكرة نفسها}. قبل ذلك، كان تغيير البرنامج
يتطلب إعادة توصيل الأسلاك يدويًا!

حسب فون نيومان، يتكون الحاسوب من \textbf{خمس وحدات}: وحدة الإدخال، وحدة الإخراج، الذاكرة، وحدة
الحساب والمنطق، ووحدة التحكم. والفكرة الأهم هي \mustalah{البرنامج المخزَّن}{Programme enregistré}{Stored program}:
التعليمات أعداد مثل المعطيات تمامًا، فنستطيع أن نحفظها في الذاكرة نفسها، ونقرأها ونغيّرها بالطريقة نفسها.

\begin{center}
\begin{tikzpicture}
    \node[unit, minimum width=4.2cm, minimum height=2.4cm] (cpu) at (0,0) {};
    \node[font=\bfseries] at (0,0.85) {المعالج};
    \node[unit=brandOrange, font=\small, minimum width=1.7cm] at (-1,-0.15) {وحدة التحكم};
    \node[unit=brandGreen, font=\small, minimum width=1.7cm] at (1,-0.15) {\textenglish{UAL}};
    \node[unit=brandPurple, minimum width=4.2cm, minimum height=2.4cm, align=center] (mem) at (-7,0) {};
    \node[font=\bfseries] at (-7,0.85) {ذاكرة واحدة};
    \node[draw=brandDark, fill=brandTeal!20, font=\small, minimum width=3.4cm] at (-7,0.1) {البرنامج (التعليمات)};
    \node[draw=brandDark, fill=brandOrange!20, font=\small, minimum width=3.4cm] at (-7,-0.6) {المعطيات};
    \draw[brandDark, line width=5pt, -{Stealth[length=4mm]}] (-4.8,0) -- (-2.2,0);
    \draw[brandDark, line width=5pt, -{Stealth[length=4mm]}] (-2.2,0) -- (-4.8,0);
    \node[font=\small, above] at (-3.5,0.2) {حافلة واحدة};
    \node[unit=brandGreen, minimum width=2.6cm, minimum height=1.2cm] (io) at (5,0) {الإدخال\\والإخراج};
    \draw[sahmT, <->] (cpu) -- (io);
\end{tikzpicture}
\end{center}

\begin{itemize}
    \item \textbf{ميزته}: بسيط ومرن. نغيّر البرنامج بمجرد تحميل برنامج آخر في الذاكرة.
    \item \textbf{عيبه}: حافلة واحدة للتعليمات والمعطيات، فلا يمكن قراءة تعليمة ومعطاة في
          اللحظة نفسها. هذا الاختناق يسمى \mustalah{عنق زجاجة فون نيومان}{Goulot de Von Neumann}{Von Neumann bottleneck}.
\end{itemize}

\begin{tashbih}
    تخيّل مكتبة لها \textbf{باب واحد ضيق}: الطلبة الذين يدخلون لإرجاع الكتب (المعطيات) والذين يدخلون
    لأخذ التعليمات (البرنامج) يمرون كلهم من الباب نفسه، واحدًا واحدًا. مهما كان الموظفون سريعين في
    الداخل، يبقى الباب هو ما يحدد السرعة. هذا هو عنق الزجاجة.
\end{tashbih}

\subsection{نموذج هارفارد \textenglish{(Harvard)}}

في هذا النموذج نستعمل \textbf{ذاكرتين منفصلتين}: واحدة للتعليمات وأخرى للمعطيات،
ولكل واحدة حافلتها الخاصة.

\begin{center}
\begin{tikzpicture}
    \node[unit, minimum width=3.4cm, minimum height=2.2cm] (cpu) at (0,0) {المعالج};
    \node[unit=brandTeal, minimum width=3.6cm, minimum height=1.6cm] (im) at (-6.5,0) {ذاكرة التعليمات\\{\small البرنامج}};
    \node[unit=brandOrange, minimum width=3.6cm, minimum height=1.6cm] (dm) at (6.5,0) {ذاكرة المعطيات\\{\small المتغيرات}};
    \draw[tealDeep, line width=4pt, -{Stealth[length=4mm]}] (im) -- (cpu) node[midway, above, font=\small, lawn=brandDark] {حافلة التعليمات};
    \draw[orangeDeep, line width=4pt, {Stealth[length=4mm]}-{Stealth[length=4mm]}] (cpu) -- (dm) node[midway, above, font=\small, lawn=brandDark] {حافلة المعطيات};
\end{tikzpicture}
\end{center}

\begin{itemize}
    \item \textbf{ميزته}: أسرع، لأن المعالج يقرأ التعليمة التالية ويقرأ المعطاة في الوقت نفسه.
    \item \textbf{عيبه}: أعقد وأغلى، والمساحة المخصصة لكل ذاكرة ثابتة.
\end{itemize}

\begin{tashbih}
    في نموذج هارفارد للمكتبة \textbf{بابان}: باب لأخذ التعليمات، وباب للمعطيات. الطابوران يتقدمان في
    الوقت نفسه، فيمر الناس أسرع. لكن بناء بابين أغلى، وإذا امتلأ أحد الجانبين لا نستطيع استعمال مساحة الجانب الآخر.
\end{tashbih}

\subsection{مقارنة واستعمالات}
\begin{center}
\begin{tabularx}{\linewidth}{|>{\centering\arraybackslash}p{3cm}|C|C|}
    \hline
    \ras{} & \ras{فون نيومان} & \ras{هارفارد} \\ \hline
    الذاكرة   & واحدة للبرنامج والمعطيات & ذاكرتان منفصلتان \\ \hline
    الحافلات  & حافلة واحدة مشتركة & حافلتان مستقلتان \\ \hline
    السرعة    & أبطأ (عنق الزجاجة) & أسرع (قراءتان معًا) \\ \hline
    التعقيد   & بسيط ورخيص & أعقد وأغلى \\ \hline
    الاستعمال & الحواسيب الشخصية، الخوادم & المتحكمات الدقيقة (\textenglish{Arduino})، معالجات الإشارة \textenglish{(DSP)} \\ \hline
\end{tabularx}
\end{center}

\begin{hal}
    المعالجات الحديثة في حاسوبك تجمع بين النموذجين: الذاكرة الرئيسية \textenglish{(RAM)} واحدة
    (فون نيومان)، لكن داخل المعالج توجد ذاكرتان صغيرتان سريعتان منفصلتان: واحدة للتعليمات
    وأخرى للمعطيات (هارفارد). نسمي ذلك \textbf{هارفارد المعدّل}، وسنعود إليه في الفصل الخامس.
\end{hal}

%================================================
\section{المعالج: الدور والمكونات}
%================================================

\begin{tarif}
    \mustalah{المعالج}{Processeur, UCT}{CPU}: الدارة التي تنفذ تعليمات البرنامج واحدة تلو الأخرى.
\end{tarif}

يتكون المعالج من ثلاثة أجزاء رئيسية (سندرسها بالتفصيل في الفصل الثاني):

\begin{center}
\begin{tabularx}{\linewidth}{|>{\centering\arraybackslash}p{4.2cm}|R|R|}
    \hline
    \ras{الجزء} & \ras{دوره} & \ras{التشبيه} \\ \hline
    وحدة التحكم \textenglish{(UC)} & تقرأ التعليمة، تفهمها، وتعطي الأوامر لباقي الأجزاء & قائد الأوركسترا \\ \hline
    وحدة الحساب والمنطق \textenglish{(UAL)} & تنفذ العمليات الحسابية (\textenglish{+ −}) والمنطقية (و، أو) & الآلة الحاسبة \\ \hline
    السجلات \textenglish{(Registres)} & ذاكرة صغيرة جدًا وسريعة جدًا داخل المعالج & ورقة المسودة على المكتب \\ \hline
\end{tabularx}
\end{center}

\subsection{الساعة وسرعة المعالج}
يعمل المعالج على إيقاع \mustalah{ساعة}{Horloge}{Clock} ترسل نبضات منتظمة. في كل نبضة يتقدم
المعالج خطوة صغيرة في عمله.

\begin{tarif}
    \textbf{التردد} $f$: عدد النبضات في الثانية، ويقاس بالهرتز \textenglish{(Hz)}.
    \textbf{الدورة} $T$: مدة نبضة واحدة، و$T = \dfrac{1}{f}$.
\end{tarif}

\begin{khatwa}[: معالج تردده \textenglish{3 GHz}]
    \begin{enumerate}
        \item $f = 3\ \text{GHz} = 3 \times 10^{9}$ نبضة في الثانية.
        \item $T = \dfrac{1}{3 \times 10^{9}} \approx 0.33 \times 10^{-9}$ ثانية، أي حوالي
              \textenglish{0.33} نانوثانية.
        \item في الوقت الذي ترمش فيه عينك (\textenglish{0.1} ثانية تقريبًا) يقوم هذا المعالج بـ
              \textenglish{300} مليون دورة!
    \end{enumerate}
\end{khatwa}

\begin{intabih}
    التردد وحده لا يحدد سرعة المعالج. ما يهمنا هو \textbf{عدد التعليمات المنفذة في الثانية}.
    معالج بـ\textenglish{3 GHz} يحتاج 4 دورات لكل تعليمة ينفذ $\frac{3 \times 10^9}{4} = 750$ مليون تعليمة
    في الثانية، بينما معالج بـ\textenglish{2 GHz} ينفذ تعليمتين في كل دورة ينفذ $4 \times 10^9$ تعليمة، فهو
    أسرع رغم أن تردده أقل.
\end{intabih}

%================================================
\section{الذاكرة}
%================================================

\subsection{كيف تبدو الذاكرة من الداخل؟}
نتخيل الذاكرة مثل \textbf{عمارة كبيرة من الصناديق}. لكل صندوق رقم يسمى \textbf{العنوان}،
وفي كل صندوق نحفظ \textbf{بايتًا واحدًا}.

\begin{center}
\begin{tikzpicture}[cellx/.style={draw=brandDark, fill=white, minimum width=2.4cm, minimum height=0.6cm, font=\codefont},
    addr/.style={font=\codefont\small, lawn=tealDeep}]
    \foreach \a/\v [count=\i from 0] in {0000/3A, 0001/FF, 0002/07, 0003/41, 0004/00} {
        \node[cellx] (c\i) at (0, -0.6*\i) {\textenglish{\v}};
        \node[addr, right=0.2cm of c\i] {\textenglish{\a}};
    }
    \node[font=\small\bfseries] at (2.15, 0.6) {العنوان};
    \node[font=\small\bfseries] at (0, 0.6) {المحتوى};
    \node[font=\small, text width=6.5cm, align=right] at (-5.5,-1.2)
        {العنوان \textenglish{0003} يحتوي على \textenglish{41}\textsubscript{\textenglish{H}}.\\
         العنوان لا يتغير أبدًا، أما المحتوى فيمكن تغييره.};
\end{tikzpicture}
\end{center}

\begin{tashbih}
    العنوان يشبه \textbf{رقم المنزل} في شارع، والمحتوى يشبه \textbf{الشخص الذي يسكن فيه}.
    رقم المنزل ثابت، أما السكان فيمكن أن يتغيروا.
\end{tashbih}

\subsection{عمليتان فقط: القراءة والكتابة}
\begin{itemize}
    \item \textbf{القراءة}: يرسل المعالج عنوانًا، فتعيد الذاكرة المحتوى (المحتوى لا يتغير).
    \item \textbf{الكتابة}: يرسل المعالج عنوانًا ومعطاة، فتحفظ الذاكرة المعطاة مكان القديمة.
\end{itemize}

\begin{mithal}
    الخانة \textenglish{0002} تحتوي على \hex{07}. ينفذ المعالج «اكتب \hex{55} في \textenglish{0002}».
    بعدها تحتوي الخانة على \hex{55}، و\hex{07} ضاعت نهائيًا. لهذا يجب أن ننتبه أين نكتب: الكتابة
    في المكان الخطأ تمحو معطيات مهمة.
\end{mithal}

\begin{tarif}
    \mustalah{ذاكرة الوصول العشوائي}{Mémoire à accès aléatoire}{Random Access Memory}: «عشوائي» هنا تعني
    أننا نستطيع الوصول إلى \textbf{أي خانة مباشرة} وفي الزمن نفسه، دون المرور بالخانات التي قبلها، عكس
    شريط الكاسيت القديم الذي يجب لفّه حتى نصل إلى الأغنية.
\end{tarif}

\subsection{أنواع الذاكرة}
\begin{center}
\begin{tikzpicture}[level distance=1.4cm,
    level 1/.style={sibling distance=7cm},
    level 2/.style={sibling distance=2.3cm},
    every node/.append style={unit, font=\small, minimum width=1.9cm},
    edge from parent/.style={draw, brandMuted, thick}]
    \node[unit=brandDark, fill=brandDark!10] {الذاكرة المركزية}
        child {node[unit=brandOrange] {\textenglish{RAM}\\{\scriptsize قراءة وكتابة}}
            child {node[unit=brandOrange] {\textenglish{SRAM}\\{\scriptsize سريعة}}}
            child {node[unit=brandOrange] {\textenglish{DRAM}\\{\scriptsize رخيصة}}}}
        child {node[unit=brandTeal] {\textenglish{ROM}\\{\scriptsize قراءة فقط}}
            child {node {\textenglish{PROM}}}
            child {node {\textenglish{EPROM}}}
            child {node {\textenglish{EEPROM}\\{\scriptsize و\textenglish{Flash}}}}};
\end{tikzpicture}
\end{center}

\begin{center}
\small
\begin{tabularx}{\linewidth}{|>{\centering\arraybackslash}p{2.2cm}|R|>{\centering\arraybackslash}p{2.4cm}|R|}
    \hline
    \ras{النوع} & \ras{الوصف} & \ras{عند انقطاع التيار} & \ras{الاستعمال} \\ \hline
    \textenglish{SRAM} & ترانزستورات فقط، سريعة جدًا وغالية & تُمحى & الذاكرة المخبئية داخل المعالج \\ \hline
    \textenglish{DRAM} & مكثّف صغير لكل بت، يجب «تنشيطه» باستمرار & تُمحى & الذاكرة الرئيسية (شرائط \textenglish{RAM}) \\ \hline
    \textenglish{ROM}  & تُكتب مرة واحدة في المصنع & تبقى & برامج الإقلاع القديمة \\ \hline
    \textenglish{EPROM} & تُمحى بالأشعة فوق البنفسجية ثم تُعاد كتابتها & تبقى & النماذج الأولية \\ \hline
    \textenglish{EEPROM / Flash} & تُمحى كهربائيًا & تبقى & \textenglish{BIOS/UEFI}، المفاتيح \textenglish{USB}، \textenglish{SSD} \\ \hline
\end{tabularx}
\end{center}

\begin{tashbih}
    \textbf{\textenglish{SRAM}} مثل مفتاح الضوء: يبقى في وضعه (مشتعل أو مطفأ) ما دام التيار موجودًا.
    أما \textbf{\textenglish{DRAM}} فمثل كأس فيه ثقب صغير: الماء (الشحنة الكهربائية) يتسرب ببطء، فيجب إعادة ملئه
    آلاف المرات في الثانية. هذه العملية تسمى \textbf{التنشيط} \textenglish{(refresh)}. لكن الكأس أصغر بكثير من المفتاح،
    لذلك نستطيع وضع مليارات منها في شريط ذاكرة رخيص.
\end{tashbih}

\begin{hal}
    لماذا يحتاج الحاسوب إلى \textenglish{ROM} أو \textenglish{Flash}؟ عند تشغيله تكون \textenglish{RAM} فارغة تمامًا.
    يجب أن يجد المعالج برنامجًا أول يبدأ به: برنامج الإقلاع (\textenglish{BIOS} أو \textenglish{UEFI})، المحفوظ في
    ذاكرة لا تُمحى. هذا البرنامج يفحص العتاد، ثم يبحث عن نظام التشغيل على القرص ويحمّله في \textenglish{RAM}.
\end{hal}

\begin{tarif}
    ذاكرة \mustalah{متطايرة}{Volatile}{Volatile} تفقد محتواها عند انقطاع التيار (\textenglish{RAM})،
    وذاكرة \mustalah{غير متطايرة}{Non volatile}{Non-volatile} تحتفظ به (\textenglish{ROM}، \textenglish{Flash}).
\end{tarif}

\subsection{حساب سعة الذاكرة}
إذا كان للمعالج $n$ خطًا (سلكًا) للعنوان، فيمكنه أن يكتب $2^n$ عنوانًا مختلفًا.

\begin{tarif}
    $\text{عدد الخانات} = 2^{n}$ \qquad و\qquad
    $\text{السعة} = 2^{n} \times \text{حجم الخانة}$
\end{tarif}

\begin{khatwa}[: سعة ذاكرة المعالج \textenglish{8086}]
    للمعالج \textenglish{8086} حافلة عناوين من \textbf{20 بتًا}، وكل خانة تحتوي على بايت واحد.
    \begin{enumerate}
        \item عدد العناوين: $2^{20}$.
        \item $2^{20} = 2^{10} \times 2^{10} = 1024 \times 1024 = 1\,048\,576$ عنوانًا.
        \item السعة $= 2^{20} \times 1$ بايت $= 1$ ميغابايت \textenglish{(1 MB)}.
    \end{enumerate}
\end{khatwa}

\begin{khatwa}[: كم خطًا للعنوان نحتاج لذاكرة \textenglish{64 KB}؟]
    \begin{enumerate}
        \item $64\ \text{KB} = 64 \times 1024 = 2^{6} \times 2^{10} = 2^{16}$ بايت.
        \item إذن $2^{n} = 2^{16}$، أي $n = 16$ خطًا.
    \end{enumerate}
\end{khatwa}

%================================================
\section{الحافلات}
%================================================

\begin{tarif}
    \mustalah{الحافلة}{Bus}{Bus}: مجموعة من الأسلاك المتوازية، تنقل المعلومات بين المعالج
    والذاكرة ووحدات الإدخال والإخراج. كل سلك ينقل بتًا واحدًا في كل مرة.
\end{tarif}

\begin{center}
\begin{tikzpicture}
    \node[unit, minimum width=2.8cm, minimum height=4cm] (cpu) at (5,0) {المعالج};
    \node[unit=brandPurple, minimum width=2.6cm, minimum height=1.2cm] (mem) at (0,2.6) {الذاكرة};
    \node[unit=brandGreen, minimum width=2.6cm, minimum height=1.2cm] (io) at (-4,2.6) {الإدخال\\والإخراج};
    % buses
    \draw[brandOrange, line width=5pt] (3.6,0.9) -- (-6,0.9);
    \draw[brandTeal!80!black, line width=5pt] (3.6,0) -- (-6,0);
    \draw[brandMuted, line width=5pt] (3.6,-0.9) -- (-6,-0.9);
    \node[font=\small\bfseries, lawn=orangeDeep, anchor=east] at (-6.1,0.9) {حافلة العناوين};
    \node[font=\small\bfseries, lawn=tealDeep, anchor=east] at (-6.1,0) {حافلة المعطيات};
    \node[font=\small\bfseries, lawn=brandMuted, anchor=east] at (-6.1,-0.9) {حافلة التحكم};
    \foreach \x/\n in {0/mem, -4/io} {
        \draw[brandOrange, line width=2pt, -{Stealth}] (\x-0.6,0.9) -- (\x-0.6,1.95);
        \draw[brandTeal!80!black, line width=2pt, {Stealth}-{Stealth}] (\x,0) -- (\x,1.95);
        \draw[brandMuted, line width=2pt, {Stealth}-{Stealth}] (\x+0.6,-0.9) -- (\x+0.6,1.95);
    }
\end{tikzpicture}
\end{center}

\begin{center}
\begin{tabularx}{\linewidth}{|>{\centering\arraybackslash}p{3cm}|>{\centering\arraybackslash}p{2.6cm}|R|}
    \hline
    \ras{الحافلة} & \ras{الاتجاه} & \ras{ماذا تنقل؟} \\ \hline
    \textbf{العناوين} \textenglish{(Address bus)} & من المعالج فقط & رقم الخانة التي يريد المعالج قراءتها أو الكتابة فيها \\ \hline
    \textbf{المعطيات} \textenglish{(Data bus)} & في الاتجاهين & القيمة المقروءة أو المكتوبة \\ \hline
    \textbf{التحكم} \textenglish{(Control bus)} & في الاتجاهين & إشارات مثل: قراءة \textenglish{(RD)}، كتابة \textenglish{(WR)}، الساعة، المقاطعة \\ \hline
\end{tabularx}
\end{center}

\begin{tashbih}
    تخيّل أنك تطلب طردًا من مكتب البريد: \textbf{العنوان} يخبر الموظف أي رف يقصد (حافلة العناوين)،
    و\textbf{الطرد} نفسه هو المعطاة (حافلة المعطيات)، وكلمة «أعطني» أو «خذ» هي إشارة التحكم
    (قراءة أو كتابة).
\end{tashbih}

\begin{khatwa}[: المعالج يقرأ الخانة \textenglish{0003}]
    \begin{enumerate}
        \item يضع المعالج العنوان \textenglish{0003} على \textbf{حافلة العناوين}.
        \item يرسل إشارة \textbf{قراءة} \textenglish{(RD)} على \textbf{حافلة التحكم}.
        \item تضع الذاكرة محتوى الخانة (\hex{41}) على \textbf{حافلة المعطيات}.
        \item يأخذ المعالج القيمة \hex{41} ويحفظها في أحد سجلاته.
    \end{enumerate}
\end{khatwa}

\subsection{عرض الحافلات}
\begin{itemize}
    \item \textbf{عرض حافلة العناوين} يحدد \textbf{سعة الذاكرة القصوى}: $n$ خطًا تعطي $2^n$ عنوانًا.
    \item \textbf{عرض حافلة المعطيات} يحدد \textbf{كم بتًا ننقل في كل مرة}: حافلة من 16 بتًا
          تنقل بايتين في المرة الواحدة.
\end{itemize}

\begin{khatwa}[: نقل 1000 بايت بحافلات مختلفة]
    \begin{enumerate}
        \item حافلة معطيات من \textbf{8 بتات} تنقل بايتًا واحدًا في كل مرة: نحتاج إلى $1000$ عملية نقل.
        \item حافلة من \textbf{16 بتًا} تنقل بايتين: $1000 \div 2 = 500$ عملية.
        \item حافلة من \textbf{64 بتًا} تنقل 8 بايتات: $1000 \div 8 = 125$ عملية فقط.
    \end{enumerate}
    كلما كانت الحافلة أعرض، انتقلت المعطيات أسرع، كما يمر عدد أكبر من السيارات في طريق سيّار
    من 8 مسالك مقارنة بطريق من مسلك واحد.
\end{khatwa}

\begin{center}
\begin{tabular}{|c|c|c|c|}
    \hline
    \ras{المعالج} & \ras{حافلة المعطيات} & \ras{حافلة العناوين} & \ras{الذاكرة القصوى} \\ \hline
    \textenglish{Intel 8086} & 16 بتًا & 20 بتًا & \textenglish{$2^{20}$ = 1 MB} \\ \hline
    \textenglish{Intel 80386} & 32 بتًا & 32 بتًا & \textenglish{$2^{32}$ = 4 GB} \\ \hline
    \textenglish{x86-64} الحديث & 64 بتًا & 48 بتًا (عمليًا) & \textenglish{$2^{48}$ = 256 TB} \\ \hline
\end{tabular}
\end{center}

%================================================
\section{تمارين}
%================================================

\begin{tamrin}
    حوّل الأعداد الآتية: (أ) \hex{3C} إلى العشري والثنائي. (ب) العدد العشري 200 إلى الست عشري.
\end{tamrin}
\begin{hall}
    (أ) $\hex{3C} = 3 \times 16 + 12 = 60$، وبالثنائي $\bin{0011\,1100}$.\\
    (ب) $200 = 12 \times 16 + 8$، والرقم 12 يُكتب \textenglish{C}، إذن $200 = $ \hex{C8}.
\end{hall}

\begin{tamrin}
    معالج له حافلة عناوين من 24 بتًا، وكل خانة ذاكرة تحتوي على بايت. ما السعة القصوى للذاكرة؟
\end{tamrin}
\begin{hall}
    $2^{24} = 2^{4} \times 2^{20} = 16 \times 1\ \text{MB} = 16\ \text{MB}$.
\end{hall}

\begin{tamrin}
    نريد ذاكرة سعتها \textenglish{4 GB}. كم خطًا تحتاج حافلة العناوين؟
\end{tamrin}
\begin{hall}
    $4\ \text{GB} = 2^{2} \times 2^{30} = 2^{32}$ بايت، إذن نحتاج إلى \textbf{32 خطًا}.
\end{hall}

\begin{tamrin}
    أجب بصحيح أو خطأ مع التعليل:
    (أ) في نموذج فون نيومان، يُحفظ البرنامج في ذاكرة مختلفة عن المعطيات.
    (ب) حافلة العناوين تنقل المعلومات في الاتجاهين.
    (ج) ذاكرة \textenglish{DRAM} تحتفظ بمحتواها بعد إطفاء الحاسوب.
\end{tamrin}
\begin{hall}
    (أ) \textbf{خطأ}: في نموذج فون نيومان ذاكرة واحدة للبرنامج والمعطيات؛ الفصل خاصية نموذج هارفارد.\\
    (ب) \textbf{خطأ}: المعالج وحده يرسل العناوين، فهي في اتجاه واحد.\\
    (ج) \textbf{خطأ}: \textenglish{DRAM} ذاكرة متطايرة، تُمحى عند انقطاع التيار.
\end{hall}

\begin{tamrin}
    معالج تردده \textenglish{2 GHz}، وتحتاج كل تعليمة في المتوسط إلى 4 دورات. كم تعليمة ينفذ في الثانية؟
\end{tamrin}
\begin{hall}
    عدد الدورات في الثانية $= 2 \times 10^{9}$. عدد التعليمات $= \dfrac{2 \times 10^{9}}{4} = 0.5 \times 10^{9}$،
    أي \textbf{500 مليون تعليمة في الثانية}.
\end{hall}

\chapter{البنية الداخلية للمعالج}

\begin{khulasa}[: ماذا سنتعلم في هذا الفصل؟]
\begin{itemize}
    \item ما الموجود داخل المعالج، ودور كل جزء.
    \item السجلات وأنواعها.
    \item وحدة الحساب والمنطق: كيف يجمع المعالج عددين بالبوابات المنطقية، وما الأعلام.
    \item وحدة التحكم: من يقرر ماذا يحدث في كل لحظة.
    \item مجموعة التعليمات وأنماط العنونة.
    \item المراحل التي تمر بها كل تعليمة: الجلب، فك الترميز، التنفيذ، الكتابة.
\end{itemize}
\end{khulasa}

%================================================
\section{نظرة إلى داخل المعالج}
%================================================

لنفتح المعالج ونرى ما بداخله. الشكل الآتي يمثل معالجًا بسيطًا، لكنه يحتوي على كل الأفكار
الأساسية الموجودة في المعالجات الحقيقية.

\begin{center}
\begin{tikzpicture}[r/.style={reg, minimum width=1.5cm},
    bus/.style={line width=3pt, color=brandTeal!70!black},
    link/.style={thick, color=brandTeal!70!black},
    sig/.style={dashed, thick, color=orangeDeep, -{Stealth[length=2mm]}}]
    % processor frame
    \draw[brandTeal, very thick, rounded corners=8pt, fill=brandTeal!4] (-8.4,-3.3) rectangle (2.2,3.4);
    \node[font=\bfseries, lawn=tealDeep, anchor=north east] at (2.1,3.35) {المعالج};
    % internal bus
    \draw[bus] (-7.6,0) -- (1.8,0);
    \node[font=\scriptsize, lawn=tealDeep] at (0.2,-0.3) {الحافلة الداخلية};
    % control unit
    \draw[orangeDeep, very thick, rounded corners=5pt, fill=brandOrange!12] (-8.1,0.5) rectangle (-4.9,3.0);
    \node[font=\small\bfseries, lawn=orangeDeep] at (-6.5,2.7) {وحدة التحكم \textenglish{(UC)}};
    \node[r] (ri) at (-6.5,2.05) {\textenglish{RI}};
    \node[reg, font=\small, minimum width=1.3cm] (dec) at (-7.25,1.05) {المفكك};
    \node[reg, font=\small, minimum width=1.3cm] (seq) at (-5.8,1.05) {المتتالي};
    \draw[sahm] (ri) -- (dec);
    \draw[sahm] (dec) -- (seq);
    \draw[link] (ri.east) -- ++(0.35,0) |- (-4.55,0.35) -- (-4.55,0);
    % clock
    \node[unit=brandPurple, font=\small] (clk) at (-5.7,3.95) {الساعة};
    \draw[sahm] (clk) -- (-5.7,3.0);
    % registers above the bus
    \node[r] (co) at (-2.6,2.3) {\textenglish{CO}};
    \node[r] (ram) at (0.6,2.3) {\textenglish{RAM}};
    \node[r] (rim) at (0.6,1.1) {\textenglish{RIM}};
    \draw[sahmT] (co) -- node[above, font=\scriptsize] {عنوان التعليمة} (ram);
    \node[font=\scriptsize, lawn=brandMuted] at (-2.6,1.65) {\textenglish{+1} بعد كل جلب};
    \draw[link] (co.south) ++(0,-0.15) -- (-2.6,0);
    \draw[link] (rim.south) -- (0.6,0);
    % ALU with two inputs, ACC under it
    \draw[brandGreen!70!black, very thick, fill=brandGreen!20]
        (-5.4,-0.8) -- (-3.6,-0.8) -- (-4.0,-1.9) -- (-5.0,-1.9) -- cycle;
    \node[font=\bfseries] at (-4.5,-1.35) {\textenglish{UAL}};
    \node[r] (acc) at (-4.5,-2.75) {\textenglish{ACC}};
    \node[r, font=\small] (fl) at (-2.3,-2.75) {الأعلام};
    \draw[sahmT] (-4.5,-1.9) -- (acc.north) node[pos=0.5, right, font=\scriptsize] {النتيجة};
    \draw[sahmT] (-3.8,-1.6) -| (fl.north);
    \draw[sahm] (-3.9,0) -- (-3.9,-0.8) node[pos=0.5, right, font=\scriptsize] {المعامل 2};
    \draw[sahm] (acc.west) -- ++(-0.9,0) |- (-6.0,-0.45) -| (-5.1,-0.8);
    \node[font=\scriptsize, lawn=brandDark] at (-6.75,-1.6) {المعامل 1};
    \draw[link] (acc.south) -- (-4.5,-3.2) -- (-1.0,-3.2) -- (-1.0,0);
    % control signals
    \draw[sig] (seq.south) -- (-5.8,0.15);
    \node[font=\scriptsize, lawn=orangeDeep, align=center] at (-7.3,0.25) {إشارات التحكم};
    % memory
    \node[unit=brandPurple, minimum width=2cm, minimum height=4cm] (mem) at (4.4,1.2) {الذاكرة\\المركزية};
    \draw[sahmO] (ram.east) -- node[above, font=\scriptsize] {حافلة العناوين} ([yshift=1.1cm]mem.west);
    \draw[sahmT, {Stealth[length=2.6mm]}-{Stealth[length=2.6mm]}] (rim.east) -- node[below, font=\scriptsize] {حافلة المعطيات} ([yshift=-0.1cm]mem.west);
\end{tikzpicture}
\end{center}

\begin{center}
\small
\begin{tabularx}{\linewidth}{|>{\centering\arraybackslash}p{2.8cm}|R|}
    \hline
    \ras{العنصر} & \ras{دوره في جملة واحدة} \\ \hline
    \textenglish{CO} (العدّاد البرنامجي) & يحفظ عنوان \textbf{التعليمة القادمة} \\ \hline
    \textenglish{RI} (سجل التعليمة) & يحفظ \textbf{التعليمة الحالية} التي نفذها الآن \\ \hline
    \textenglish{RAM} (سجل عنوان الذاكرة) & يحفظ العنوان الذي نرسله إلى الذاكرة \\ \hline
    \textenglish{RIM} (سجل معطيات الذاكرة) & يحفظ المعطاة القادمة من الذاكرة أو الذاهبة إليها \\ \hline
    \textenglish{ACC} (المركِّم) & يحفظ نتيجة الحساب \\ \hline
    الأعلام & تحفظ معلومات عن آخر نتيجة: هل هي صفر؟ سالبة؟ \dots \\ \hline
    \textenglish{UAL} & تنفذ العمليات الحسابية والمنطقية \\ \hline
    المفكك والمتتالي & يفهمان التعليمة ويرسلان الأوامر بالترتيب الصحيح \\ \hline
\end{tabularx}
\end{center}

\begin{intabih}
    لا تخلط بين \textenglish{RAM} الذاكرة (\textenglish{Random Access Memory}) و\textenglish{RAM}
    السجل (\textenglish{Registre d'Adresse Mémoire}) في الكتب الفرنسية. في الإنجليزية يسمى هذا
    السجل \textenglish{MAR}، ويسمى \textenglish{RIM} سجل \textenglish{MDR}، ويسمى \textenglish{CO}
    سجل \textenglish{PC} \textenglish{(Program Counter)}، وفي المعالج \textenglish{8086} نسميه \textenglish{IP}.
\end{intabih}

%================================================
\section{السجلات}
%================================================

\begin{tarif}
    \mustalah{السجل}{Registre}{Register}: خانة ذاكرة صغيرة جدًا (8 أو 16 أو 32 أو 64 بتًا)
    موجودة \textbf{داخل المعالج}. هي أسرع ذاكرة في الحاسوب كله.
\end{tarif}

\begin{tashbih}
    عندما تحل تمرينًا في الرياضيات، لا تذهب إلى المكتبة في كل مرة لتقرأ عددًا! تكتب الأعداد التي
    تحتاجها على \textbf{ورقة المسودة} أمامك. السجلات هي ورقة المسودة للمعالج، والذاكرة هي المكتبة.
\end{tashbih}

\subsection{لماذا السجلات مهمة للسرعة؟}
\begin{center}
\begin{tabular}{|c|c|c|}
    \hline
    \ras{المكان} & \ras{زمن الوصول التقريبي} & \ras{لو كانت الدورة ثانية واحدة\dots} \\ \hline
    سجل & أقل من \textenglish{1 ns} & ثانية واحدة \\ \hline
    الذاكرة الرئيسية \textenglish{RAM} & حوالي \textenglish{80 ns} & دقيقة ونصف تقريبًا \\ \hline
    قرص \textenglish{SSD} & حوالي \textenglish{100 µs} & أكثر من يوم \\ \hline
\end{tabular}
\end{center}
لهذا نضع في السجلات المعطيات التي نستعملها كثيرًا، ونقلل من الذهاب إلى الذاكرة.

\subsection{أنواع السجلات}
\begin{enumerate}
    \item \mustalah{سجلات المعطيات}{Registres de données}{Data registers}: تحفظ الأعداد التي نحسب بها
          (مثل \textenglish{AX}، \textenglish{BX} في المعالج \textenglish{8086}).
    \item \mustalah{سجلات العناوين}{Registres d'adresse}{Address registers}: تحفظ عناوين في
          الذاكرة (مثل \textenglish{SI}، \textenglish{DI}، \textenglish{SP}).
    \item \mustalah{سجلات التحكم والحالة}{Registres de contrôle}{Control registers}: تساعد المعالج
          على تنظيم عمله: \textenglish{CO}، \textenglish{RI}، وسجل الأعلام.
\end{enumerate}

\begin{mithal}[: السجلات أثناء العمل]
    المعالج ينفذ الآن التعليمة الموجودة في العنوان 11، وهي «اجمع \textenglish{ACC} مع محتوى الخانة 21».
    \begin{itemize}
        \item \textenglish{CO} يحتوي على 12: عنوان التعليمة \textbf{القادمة}، لا الحالية.
        \item \textenglish{RI} يحتوي على التعليمة الحالية نفسها، لكي تقرأها وحدة التحكم طيلة تنفيذها.
        \item \textenglish{RAM} يحتوي على 21: العنوان الذي نرسله الآن إلى الذاكرة.
        \item \textenglish{RIM} يستقبل محتوى الخانة 21 القادم من الذاكرة.
        \item \textenglish{ACC} يحتوي على العدد الأول، وسيستقبل النتيجة.
    \end{itemize}
    كل سجل له دور واضح، ولهذا تسمى سجلات خاصة، عكس السجلات العامة التي يستعملها المبرمج كما يشاء.
\end{mithal}

\subsection{سجل الأعلام}
\begin{tarif}
    \mustalah{العلم}{Indicateur, drapeau}{Flag}: بت واحد تضعه وحدة الحساب في 1 أو 0 حسب
    نتيجة آخر عملية. نستعمل الأعلام بعد ذلك لاتخاذ القرارات (مثل: «إذا كانت النتيجة صفرًا، اقفز»).
\end{tarif}

\begin{center}
\begin{tabularx}{\linewidth}{|>{\centering\arraybackslash}p{2.6cm}|R|R|}
    \hline
    \ras{العلم} & \ras{يصبح 1 عندما\dots} & \ras{مثال على 8 بتات} \\ \hline
    \textenglish{ZF} (الصفر) & تكون النتيجة \textbf{صفرًا} & $5 - 5 = 0$ \\ \hline
    \textenglish{CF} (الحمل) & يخرج حمل من آخر بت (النتيجة أكبر من السعة) & $200 + 100 = 300 > 255$ \\ \hline
    \textenglish{SF} (الإشارة) & يكون البت الأخير للنتيجة 1 (عدد سالب) & $3 - 7 = -4$ \\ \hline
    \textenglish{OF} (الفيضان) & تكون النتيجة خاطئة إذا اعتبرنا الأعداد بإشارة & $100 + 50 = 150 > 127$ \\ \hline
\end{tabularx}
\end{center}

\subsection{الفرق بين \textenglish{CF} و\textenglish{OF}}
البايت نفسه يمكن قراءته بطريقتين، والمعالج لا يعرف أي طريقة يقصدها المبرمج:
\begin{itemize}
    \item \textbf{دون إشارة} \textenglish{(unsigned)}: القيم من 0 إلى 255.
    \item \textbf{بإشارة} \textenglish{(signed)}: القيم من $-128$ إلى $+127$. البت الأخير (البت 7) يدل على الإشارة: 1 = سالب.
\end{itemize}
لذلك يحسب المعالج العلمين معًا: \textenglish{CF} يهم من يعمل بأعداد دون إشارة، و\textenglish{OF} يهم من يعمل
بأعداد بإشارة.

\begin{khatwa}[: $100 + 50$ على 8 بتات]
    \begin{enumerate}
        \item $100 = $ \bin{0110\,0100} و$50 = $ \bin{0011\,0010}.
        \item المجموع: $\bin{1001\,0110} = 150$ (دون إشارة). لا حمل من البت 7، إذن \textenglish{CF = 0}: صحيح لمن يعمل دون إشارة.
        \item لكن البت 7 أصبح 1! من يقرأ بإشارة يرى عددًا سالبًا: $150 - 256 = -106$.
        \item جمعنا عددين موجبين فحصلنا على عدد سالب: هذا مستحيل، إذن النتيجة خاطئة بإشارة، و\textenglish{OF = 1}.
    \end{enumerate}
\end{khatwa}

\begin{intabih}
    قاعدة سهلة لـ\textenglish{OF}: يحدث الفيضان عندما نجمع عددين \textbf{من الإشارة نفسها} فنحصل على نتيجة
    \textbf{من الإشارة الأخرى}. جمع موجب مع سالب لا يسبب فيضانًا أبدًا.
\end{intabih}

%================================================
\section{وحدة الحساب والمنطق \textenglish{(UAL)}}
%================================================

\begin{tarif}
    \mustalah{وحدة الحساب والمنطق}{Unité Arithmétique et Logique, UAL}{Arithmetic Logic Unit, ALU}:
    الجزء الذي ينفذ العمليات الحسابية (الجمع، الطرح، المقارنة\dots) والمنطقية (\textenglish{AND}،
    \textenglish{OR}، \textenglish{XOR}، \textenglish{NOT}، الإزاحة).
\end{tarif}

\begin{center}
\begin{tikzpicture}
    \draw[brandGreen!70!black, very thick, fill=brandGreen!20]
        (-2.2,1) -- (-0.4,1) -- (0,0.4) -- (0.4,1) -- (2.2,1) -- (1.2,-1) -- (-1.2,-1) -- cycle;
    \node[font=\bfseries\large] at (0,-0.2) {\textenglish{UAL}};
    \draw[sahmT] (-1.3,2) node[above, font=\small] {المعامل \textenglish{A}} -- (-1.3,1);
    \draw[sahmT] (1.3,2) node[above, font=\small] {المعامل \textenglish{B}} -- (1.3,1);
    \draw[sahmO] (3.6,0) node[right, font=\small] {العملية (من وحدة التحكم)} -- (1.7,0);
    \draw[sahmT] (0,-1) -- (0,-2) node[below, font=\small] {النتيجة};
    \draw[sahm] (-1.6,-0.2) -- (-3.4,-0.2) node[left, font=\small] {الأعلام \textenglish{ZF, CF, SF, OF}};
\end{tikzpicture}
\end{center}

\subsection{كيف تجمع البوابات المنطقية؟ نصف المجمّع}

لنجمع بتين $A$ و$B$. هناك أربع حالات فقط:

\begin{center}
\begin{tabular}{|c|c||c|c|}
    \hline
    \ras{$A$} & \ras{$B$} & \ras{المجموع $S$} & \ras{الحمل $C$} \\ \hline
    0 & 0 & 0 & 0 \\ \hline
    0 & 1 & 1 & 0 \\ \hline
    1 & 0 & 1 & 0 \\ \hline
    1 & 1 & 0 & \cell{orangeDeep}{\textbf{1}} \\ \hline
\end{tabular}
\end{center}

نلاحظ أن: $S$ يساوي 1 عندما يكون بت واحد فقط يساوي 1، وهذه هي البوابة \textenglish{XOR}.
و$C$ يساوي 1 عندما يكون البتان يساويان 1، وهذه هي البوابة \textenglish{AND}.

\begin{tarif}
    \mustalah{نصف المجمّع}{Demi-additionneur}{Half adder}: \quad
    $S = A \oplus B$ \quad و\quad $C = A \cdot B$
\end{tarif}

\begin{center}
\begin{tikzpicture}[gate/.style={draw=brandDark, fill=white, thick, minimum width=1.5cm,
        minimum height=0.9cm, font=\codefont\bfseries, rounded corners=2pt}]
    \node[gate] (x) at (0,0.8) {\textenglish{XOR}};
    \node[gate] (a) at (0,-0.8) {\textenglish{AND}};
    \node[font=\bfseries] (A) at (-3.2,0.4) {\textenglish{A}};
    \node[font=\bfseries] (B) at (-3.2,-0.4) {\textenglish{B}};
    \fill (-2,0.4) circle (2pt); \fill (-1.5,-0.4) circle (2pt);
    \draw[thick] (A) -- (-2,0.4) |- ([yshift=0.15cm]x.west);
    \draw[thick] (-2,0.4) |- ([yshift=0.15cm]a.west);
    \draw[thick] (B) -- (-1.5,-0.4) |- ([yshift=-0.15cm]x.west);
    \draw[thick] (-1.5,-0.4) |- ([yshift=-0.15cm]a.west);
    \draw[sahmT] (x.east) -- (2,0.8) node[right] {\textenglish{S}};
    \draw[sahmO] (a.east) -- (2,-0.8) node[right] {\textenglish{C}};
\end{tikzpicture}
\end{center}

\subsection{المجمّع الكامل}
لجمع أعداد من عدة بتات، نحتاج أن نضيف \textbf{الحمل القادم} من البت السابق. المجمّع الكامل
له ثلاثة مداخل: $A$ و$B$ والحمل الداخل $C_{in}$.

\begin{tarif}
    \mustalah{المجمّع الكامل}{Additionneur complet}{Full adder}: \quad
    $S = A \oplus B \oplus C_{in}$ \quad و\quad
    $C_{out} = A \cdot B + C_{in} \cdot (A \oplus B)$
\end{tarif}

\begin{center}
\small
\begin{tabular}{|c|c|c||c|c|}
    \hline
    \ras{$A$} & \ras{$B$} & \ras{$C_{in}$} & \ras{$S$} & \ras{$C_{out}$} \\ \hline
    0&0&0&0&0 \\ \hline 0&0&1&1&0 \\ \hline 0&1&0&1&0 \\ \hline 0&1&1&0&1 \\ \hline
    1&0&0&1&0 \\ \hline 1&0&1&0&1 \\ \hline 1&1&0&0&1 \\ \hline 1&1&1&1&1 \\ \hline
\end{tabular}
\end{center}

\subsection{مجمّع 4 بتات}
نضع 4 مجمّعات كاملة جنبًا إلى جنب: الحمل الخارج من كل واحد يدخل في الذي بعده،
تمامًا كما نفعل عند الجمع بالورقة والقلم.

\begin{center}
\begin{tikzpicture}[fa/.style={draw=brandTeal, fill=brandTeal!15, very thick, minimum width=1.6cm,
        minimum height=1.1cm, font=\codefont\bfseries, rounded corners=3pt}]
    \foreach \i/\x in {0/0, 1/-3, 2/-6, 3/-9} {
        \node[fa] (f\i) at (\x,0) {\textenglish{FA}};
        \draw[thick] ([xshift=-0.35cm]f\i.north) -- ++(0,0.6) node[above, font=\small] {\textenglish{A\i}};
        \draw[thick] ([xshift=0.35cm]f\i.north) -- ++(0,0.6) node[above, font=\small] {\textenglish{B\i}};
        \draw[sahmT] (f\i.south) -- ++(0,-0.6) node[below, font=\small] {\textenglish{S\i}};
    }
    \draw[sahmO] (1.7,0) node[right, font=\small] {\textenglish{0}} -- (f0.east);
    \draw[sahmO] (f0.west) -- (f1.east);
    \draw[sahmO] (f1.west) -- (f2.east);
    \draw[sahmO] (f2.west) -- (f3.east);
    \draw[sahmO] (f3.west) -- ++(-0.9,0) node[left, font=\small] {\textenglish{CF}};
\end{tikzpicture}
\end{center}

\begin{khatwa}[: جمع $0111 + 0110$ (أي $7 + 6$) بالمجمّع]
    \begin{center}
    \begin{tabular}{|c|c|c|c|c|c|}
        \hline
        \ras{البت} & \ras{$A$} & \ras{$B$} & \ras{$C_{in}$} & \ras{$S$} & \ras{$C_{out}$} \\ \hline
        0 & 1 & 0 & 0 & 1 & 0 \\ \hline
        1 & 1 & 1 & 0 & 0 & 1 \\ \hline
        2 & 1 & 1 & 1 & 1 & 1 \\ \hline
        3 & 0 & 0 & 1 & 1 & 0 \\ \hline
    \end{tabular}
    \end{center}
    نقرأ $S$ من البت 3 إلى البت 0: $1101 = 13$، وهو فعلًا $7 + 6$. الحمل الأخير 0، إذن \textenglish{CF = 0}.
\end{khatwa}

\subsection{والطرح؟ نستعمل المجمّع نفسه!}
لا يحتاج المعالج إلى دارة خاصة للطرح. يستعمل فكرة \mustalah{المتمم لاثنين}{Complément à deux}{Two's complement}:
$A - B = A + (-B)$، و$-B$ نحصل عليه بقلب كل بتات $B$ ثم إضافة 1.

\begin{khatwa}[: $7 - 5$ على 4 بتات]
    \begin{enumerate}
        \item $5 = 0101$. نقلب البتات: $1010$. نضيف 1: $1011$. إذن $-5 = 1011$.
        \item نجمع: $0111 + 1011 = 1\,0010$. الحمل الخارج نتركه، فيبقى $0010 = 2$.
        \item النتيجة $7 - 5 = 2$. صحيح! والدارة نفسها جمعت وطرحت.
    \end{enumerate}
\end{khatwa}

\subsection{المقارنة بالطرح}
لا يحتاج المعالج إلى دارة خاصة للمقارنة: ليقارن $A$ و$B$ يحسب $A - B$ \textbf{دون أن يحفظ النتيجة}،
ثم ينظر إلى الأعلام:
\begin{itemize}
    \item إذا كان \textenglish{ZF = 1} فإن $A = B$.
    \item إذا كان \textenglish{CF = 1} (حدث «استلاف») فإن $A < B$ (للأعداد دون إشارة).
\end{itemize}

\subsection{تجربة حقيقية: كيف تتغير الأعلام؟}
نفذنا البرنامج الآتي على المعالج \textenglish{8086}. لا تقلق إن لم تفهم كل التعليمات الآن،
سندرسها في الفصل الثالث؛ ركّز فقط على الأعلام.

\asm{ch2/flags.asm}

\begin{center}
\small
\begin{tabular}{|c|l|c|c|c|c|c|}
    \hline
    \ras{السطر} & \ras{التعليمة} & \ras{\textenglish{AL}} & \ras{\textenglish{CF}} & \ras{\textenglish{ZF}} & \ras{\textenglish{SF}} & \ras{\textenglish{OF}} \\ \hline
    4 & \code{add al, 100} ($200 + 100$) & $\hex{2C} = 44$ & \cell{orangeDeep}{\textbf{1}} & 0 & 0 & 0 \\ \hline
    5 & \code{mov al, 100} & $\hex{64} = 100$ & 1 & 0 & 0 & 0 \\ \hline
    6 & \code{add al, 50} ($100 + 50$) & $\hex{96} = 150$ & 0 & 0 & 1 & \cell{orangeDeep}{\textbf{1}} \\ \hline
    8 & \code{sub al, 5} ($5 - 5$) & $\hex{00} = 0$ & 0 & \cell{orangeDeep}{\textbf{1}} & 0 & 0 \\ \hline
    10 & \code{sub al, 7} ($3 - 7$) & $\hex{FC} = 252$ & \cell{orangeDeep}{\textbf{1}} & 0 & \cell{orangeDeep}{\textbf{1}} & 0 \\ \hline
\end{tabular}
\end{center}

\begin{khatwa}[: لماذا $200 + 100$ أعطى 44؟]
    \begin{enumerate}
        \item $200 + 100 = 300$، لكن \textenglish{AL} يحتوي على 8 بتات فقط، أي من 0 إلى 255.
        \item $300 - 256 = 44$: يبقى في \textenglish{AL} العدد 44، و«الـ256» التي لم تجد مكانًا تذهب إلى
              العلم \textenglish{CF}، فيصبح \textenglish{CF = 1}.
    \end{enumerate}
\end{khatwa}

\begin{intabih}
    لاحظ السطر 5: التعليمة \code{mov} \textbf{لا تغيّر الأعلام}، لذلك بقي \textenglish{CF = 1}. فقط
    تعليمات الحساب والمنطق تغيّر الأعلام.
\end{intabih}

%================================================
\section{وحدة التحكم \textenglish{(UC)}}
%================================================

\begin{tarif}
    \mustalah{وحدة التحكم}{Unité de Commande, UC}{Control Unit}: الجزء الذي يقرأ التعليمات
    ويفهمها، ثم يرسل في كل نبضة ساعة \textbf{إشارات} إلى باقي الأجزاء (افتح هذا السجل، اجمع،
    اقرأ من الذاكرة\dots) بالترتيب الصحيح.
\end{tarif}

\begin{tashbih}
    وحدة التحكم هي \textbf{قائد الأوركسترا}: لا يعزف أي آلة بنفسه، لكنه يقول لكل عازف متى يبدأ ومتى
    يتوقف. بدونه يعزف الجميع في الوقت نفسه فتصبح ضوضاء.
\end{tashbih}

مكوناتها الأساسية:
\begin{itemize}
    \item \textbf{سجل التعليمة} \textenglish{(RI)}: يحفظ التعليمة الحالية.
    \item \mustalah{المفكك}{Décodeur}{Decoder}: يقرأ الرمز العملي \textenglish{(opcode)} ويعرف
          ما المطلوب: جمع؟ نقل؟ قفز؟
    \item \mustalah{المتتالي}{Séquenceur}{Sequencer}: يرسل الإشارات المناسبة، واحدة بعد الأخرى،
          على إيقاع الساعة.
\end{itemize}

\subsection{الأوامر الدقيقة نبضة بنبضة}
كل مرحلة من مراحل التعليمة تنقسم إلى خطوات صغيرة جدًا تسمى \mustalah{العمليات الدقيقة}{Micro-opérations}{Micro-operations}.
في كل نبضة ساعة يرسل المتتالي إشارات تفتح «أبوابًا» بين السجلات. هذا مثال لجلب تعليمة:

\begin{center}
\small
\begin{tabularx}{\linewidth}{|c|R|R|}
    \hline
    \ras{النبضة} & \ras{العملية الدقيقة} & \ras{الإشارة التي يرسلها المتتالي} \\ \hline
    \textenglish{T1} & \nql{RAM}{CO} & «انسخ \textenglish{CO} في \textenglish{RAM}» \\ \hline
    \textenglish{T2} & \nql{RIM}{[RAM]}، \nql{CO}{CO + 1} & «اقرأ من الذاكرة» و«زِد \textenglish{CO}» \\ \hline
    \textenglish{T3} & \nql{RI}{RIM} & «انسخ \textenglish{RIM} في \textenglish{RI}» \\ \hline
    \textenglish{T4}\dots & حسب التعليمة & المفكك يختار الإشارات التالية \\ \hline
\end{tabularx}
\end{center}
نلاحظ أن الجلب متشابه في كل التعليمات، أما ما بعده فيختلف حسب الرمز العملي.

\subsection{نوعان من وحدات التحكم}
\begin{center}
\begin{tabularx}{\linewidth}{|>{\centering\arraybackslash}p{3.2cm}|R|R|}
    \hline
    \ras{} & \ras{المبرمجة بالأسلاك \textenglish{(câblée)}} & \ras{المبرمجة دقيقًا \textenglish{(microprogrammée)}} \\ \hline
    المبدأ & دارة منطقية ثابتة تولّد الإشارات & كل تعليمة = برنامج صغير مخزن في ذاكرة داخلية \\ \hline
    السرعة & سريعة جدًا & أبطأ قليلًا \\ \hline
    التعديل & صعب (يجب تغيير الدارة) & سهل (نغيّر البرنامج الدقيق) \\ \hline
    أين نجدها & معالجات \textenglish{RISC} & معالجات \textenglish{CISC} مثل \textenglish{8086} \\ \hline
\end{tabularx}
\end{center}

%================================================
\section{مجموعة التعليمات}
%================================================

\begin{tarif}
    \mustalah{مجموعة التعليمات}{Jeu d'instructions}{Instruction Set}: قائمة كل التعليمات التي
    يعرف المعالج تنفيذها. كل معالج له مجموعته الخاصة.
\end{tarif}

\subsection{شكل التعليمة}
كل تعليمة في لغة الآلة تتكون من جزأين:

\begin{center}
\begin{tikzpicture}
    \node[draw=brandOrange, fill=brandOrange!20, thick, minimum width=3.5cm, minimum height=1cm, font=\bfseries] (op) at (0,0) {الرمز العملي};
    \node[draw=brandTeal, fill=brandTeal!20, thick, minimum width=5cm, minimum height=1cm, font=\bfseries, anchor=east] (oper) at (op.west) {المعاملات};
    \node[below=0.15cm of op, font=\small] {\textenglish{opcode}: ماذا نفعل؟};
    \node[below=0.15cm of oper, font=\small] {\textenglish{operands}: على ماذا نعمل؟};
    \node[above=0.2cm of op, font=\codefont] {\textenglish{ADD}};
    \node[above=0.2cm of oper, font=\codefont] {\textenglish{AX, 5}};
\end{tikzpicture}
\end{center}

\begin{tarif}
    \mustalah{لغة التجميع}{Langage d'assemblage}{Assembly language}: كتابة التعليمات بكلمات قصيرة
    يفهمها الإنسان (\code{ADD}، \code{MOV}) بدل الأعداد الثنائية. برنامج \mustalah{المجمِّع}{Assembleur}{Assembler}
    يحوّلها إلى لغة الآلة (أصفار وآحاد).
\end{tarif}

\begin{mithal}
    في المعالج \textenglish{8086}، التعليمة \code{add al, 5} تُكتب في لغة الآلة بالبايتين \textenglish{04 05}، أي
    \textenglish{0000 0100 0000 0101}: البايت الأول هو الرمز العملي «أضف قيمة فورية إلى \textenglish{AL}»، والثاني هو العدد 5.
    المبرمج يكتب \code{add al, 5}، والمجمِّع يكتب الأصفار والآحاد بدلًا منه (سنرى ذلك بالتفصيل في الفصل الثالث).
\end{mithal}

\subsection{كم تعليمة يمكن أن يملك المعالج؟}
إذا كان الرمز العملي مكتوبًا على $k$ بتات، فهناك $2^k$ رمزًا مختلفًا، أي $2^k$ تعليمة على الأكثر.

\begin{khatwa}
    رمز عملي من 6 بتات: $2^6 = 64$ تعليمة ممكنة. رمز من 8 بتات: $2^8 = 256$ تعليمة.
\end{khatwa}

\subsection{عدد العناوين في التعليمة}
لنحسب \textenglish{A = B + C} في أربعة أنواع من الآلات. كلما قلّ عدد المعاملات في التعليمة، قصرت
التعليمة، لكن احتجنا إلى تعليمات أكثر:

\begin{center}
\small
\begin{tabularx}{\linewidth}{|>{\centering\arraybackslash}p{3cm}|>{\raggedright\arraybackslash}X|>{\centering\arraybackslash}p{3cm}|}
    \hline
    \ras{نوع الآلة} & \ras{البرنامج} & \ras{عدد التعليمات} \\ \hline
    3 عناوين & \code{ADD A, B, C} & 1 \\ \hline
    عنوانان & \code{MOV A, B} ثم \code{ADD A, C} & 2 \\ \hline
    عنوان واحد (بمركِّم) & \code{LOAD B} ثم \code{ADD C} ثم \code{STORE A} & 3 \\ \hline
    دون عنوان (بمكدس) & \code{PUSH B}، \code{PUSH C}، \code{ADD}، \code{POP A} & 4 \\ \hline
\end{tabularx}
\end{center}
المعالج \textenglish{8086} آلة \textbf{بعنوانين}: \code{add ax, bx} تعني $\text{AX} \leftarrow \text{AX} + \text{BX}$.

\subsection{أصناف التعليمات}
\begin{center}
\begin{tabularx}{\linewidth}{|>{\centering\arraybackslash}p{3.5cm}|R|>{\centering\arraybackslash}p{4.2cm}|}
    \hline
    \ras{الصنف} & \ras{الدور} & \ras{أمثلة (\textenglish{8086})} \\ \hline
    نقل المعطيات & نسخ قيمة من مكان إلى آخر & \code{MOV}، \code{PUSH}، \code{POP} \\ \hline
    الحساب & الجمع، الطرح، الضرب، القسمة & \code{ADD}، \code{SUB}، \code{MUL}، \code{INC} \\ \hline
    المنطق & العمليات على البتات & \code{AND}، \code{OR}، \code{XOR}، \code{SHL} \\ \hline
    التحكم في التدفق & القفز وتغيير ترتيب التنفيذ & \code{JMP}، \code{JE}، \code{CALL}، \code{RET} \\ \hline
    الإدخال والإخراج & التواصل مع الأجهزة & \code{IN}، \code{OUT}، \code{INT} \\ \hline
\end{tabularx}
\end{center}

\subsection{\textenglish{CISC} و\textenglish{RISC}}
\begin{itemize}
    \item \textbf{\textenglish{CISC}} \textenglish{(Complex Instruction Set Computer)}: تعليمات كثيرة
          ومعقدة، بعضها يقوم بعمل كبير. مثال: \textenglish{Intel x86}.
    \item \textbf{\textenglish{RISC}} \textenglish{(Reduced Instruction Set Computer)}: تعليمات قليلة
          وبسيطة وسريعة، كلها بالطول نفسه. مثال: \textenglish{ARM} في الهواتف، \textenglish{RISC-V}.
\end{itemize}

%================================================
\section{أنماط العنونة}
%================================================

\begin{tarif}
    \mustalah{نمط العنونة}{Mode d'adressage}{Addressing mode}: الطريقة التي نحدد بها
    \textbf{أين توجد المعطاة} التي تعمل عليها التعليمة.
\end{tarif}

\begin{tashbih}
    تريد أن تعطي صديقك كتابًا:
    \textbf{(1)} تعطيه الكتاب في يده (فوري)،
    \textbf{(2)} تقول له: «الكتاب على الرف رقم 12» (مباشر)،
    \textbf{(3)} تقول له: «انظر في الورقة الموجودة في جيبي، فيها رقم الرف» (غير مباشر)،
    \textbf{(4)} تقول: «الرف الثالث بعد الرف رقم 10» (مفهرس).
\end{tashbih}

\subsection{الأنماط الأساسية}
\begin{center}
\small
\begin{tabularx}{\linewidth}{|>{\centering\arraybackslash}p{2.4cm}|>{\centering\arraybackslash}p{3.3cm}|R|}
    \hline
    \ras{النمط} & \ras{مثال \textenglish{8086}} & \ras{أين المعطاة؟} \\ \hline
    فوري \textenglish{(immédiat)} & \code{mov ax, 5} & داخل التعليمة نفسها \\ \hline
    سجلي \textenglish{(registre)} & \code{mov bx, ax} & في سجل \\ \hline
    مباشر \textenglish{(direct)} & \code{mov cx, [val]} & في الذاكرة، وعنوانها مكتوب في التعليمة \\ \hline
    غير مباشر \textenglish{(indirect)} & \code{mov dx, [si]} & في الذاكرة، وعنوانها موجود في سجل \\ \hline
    مفهرس \textenglish{(indexé)} & \code{mov ax, [tab+di]} & العنوان = عنوان بداية + قيمة سجل الفهرس \\ \hline
    نسبي \textenglish{(relatif)} & \code{jmp suite} & العنوان = العنوان الحالي + إزاحة (للقفز) \\ \hline
\end{tabularx}
\end{center}

\subsection{متى نستعمل كل نمط؟}
\begin{itemize}
    \item \textbf{الفوري}: للثوابت (عدد الدورات، القيمة الأولية). أسرع نمط، لأن القيمة وصلت مع التعليمة،
          فلا نذهب إلى الذاكرة. لكن لا يمكن تغييرها أثناء التنفيذ.
    \item \textbf{السجلي}: للحسابات السريعة بين قيم موجودة داخل المعالج. سريع جدًا أيضًا.
    \item \textbf{المباشر}: للمتغيرات التي لها مكان ثابت في الذاكرة (مثل متغير عام). نحتاج إلى قراءة واحدة من الذاكرة.
    \item \textbf{غير المباشر}: عندما لا نعرف العنوان عند كتابة البرنامج، بل نحسبه أثناء التنفيذ. هو أساس
          \textbf{المؤشرات} في لغة \textenglish{C}.
    \item \textbf{المفهرس}: للجداول: نثبت عنوان البداية، ونغيّر الفهرس في كل دورة من الحلقة لنمر على كل العناصر.
    \item \textbf{النسبي}: للقفزات. بما أن المسافة نسبية، يعمل البرنامج حتى لو حُمّل في مكان آخر من الذاكرة.
\end{itemize}

\begin{khatwa}[: كم مرة نذهب إلى الذاكرة؟]
    \begin{center}
    \begin{tabular}{|l|c|}
        \hline \ras{التعليمة} & \ras{قراءات الذاكرة لجلب المعطاة} \\ \hline
        \code{mov ax, 5} (فوري) & 0 \\ \hline
        \code{mov ax, bx} (سجلي) & 0 \\ \hline
        \code{mov ax, [val]} (مباشر) & 1 \\ \hline
        \code{mov ax, [si]} (غير مباشر بسجل) & 1 \\ \hline
        غير مباشر بالذاكرة (العنوان نفسه في الذاكرة) & 2 \\ \hline
    \end{tabular}
    \end{center}
    كل ذهاب إلى الذاكرة يكلف وقتًا، لذلك يحاول المبرمج إبقاء المعطيات المهمة في السجلات.
\end{khatwa}

\subsection{الأنماط في برنامج حقيقي}
\asm{ch2/modes.asm}

\begin{center}
\begin{tikzpicture}[cellx/.style={draw=brandDark, fill=white, minimum width=2cm, minimum height=0.55cm, font=\codefont\small},
    addr/.style={font=\codefont\small, lawn=tealDeep}]
    \foreach \a/\v/\n [count=\i from 0] in {0117/1234/val, 0119/000A/tab, 011B/0014/{}, 011D/001E/{}} {
        \node[cellx] (c\i) at (0, -0.55*\i) {\textenglish{\v}};
        \node[addr, right=0.15cm of c\i] {\textenglish{\a}};
        \node[font=\codefont\small, left=0.15cm of c\i] {\textenglish{\n}};
    }
    \node[font=\small\bfseries] at (0,0.55) {الذاكرة};
    \draw[sahmO] (2.9,-0.55) -- (2.9,-1.65) node[midway, right, font=\scriptsize] {\textenglish{+4}};
\end{tikzpicture}
\end{center}

\begin{sharh}
    \s{1}{برنامج من النوع \textenglish{.COM}: يبدأ في العنوان \hex{100}.}
    \s{3}{\textbf{فوري}: العدد 5 مكتوب داخل التعليمة. النتيجة: \textenglish{AX = 0005}.}
    \s{4}{\textbf{سجلي}: ننسخ \textenglish{AX} في \textenglish{BX}. النتيجة: \textenglish{BX = 0005}.}
    \s{5}{\textbf{مباشر}: نقرأ الذاكرة في عنوان \code{val} (\hex{0117}). النتيجة: \textenglish{CX = 1234}.}
    \s{6}{\code{lea} تضع \textbf{عنوان} \code{val} (لا قيمته) في \textenglish{SI}: \textenglish{SI = 0117}.}
    \s{7}{\textbf{غير مباشر}: نقرأ الذاكرة في العنوان الموجود في \textenglish{SI}. النتيجة: \textenglish{DX = 1234}.}
    \s{8}{\textenglish{DI = 4}: كل عنصر من \code{tab} كلمة (بايتان)، إذن العنصر الثالث على بعد 4 بايتات.}
    \s{9}{\textbf{مفهرس}: العنوان $\hex{0119} + 4 = \hex{011D}$، وفيه 30. النتيجة: \textenglish{AX = 001E}.}
    \s{10}{نهاية البرنامج والعودة إلى النظام.}
\end{sharh}

\begin{intabih}
    القوسان \code{[ ]} يعنيان دائمًا: «\textbf{محتوى الذاكرة} في هذا العنوان». \code{mov dx, si}
    ينسخ العدد \hex{0117} نفسه، أما \code{mov dx, [si]} فيقرأ ما هو مخزن في العنوان \hex{0117}.
\end{intabih}

%================================================
\section{مراحل تنفيذ تعليمة}
%================================================

كل تعليمة، مهما كانت بسيطة، تمر بالمراحل نفسها. نسمي ذلك \mustalah{دورة التعليمة}{Cycle d'instruction}{Instruction cycle}.

\begin{center}
\begin{tikzpicture}
    \foreach \a/\t/\c/\d in {90/{1. الجلب\\\textenglish{Fetch}}/brandTeal/{نقرأ التعليمة من الذاكرة},
                             0/{2. فك الترميز\\\textenglish{Decode}}/brandOrange/{نفهم ما المطلوب},
                             -90/{3. التنفيذ\\\textenglish{Execute}}/brandGreen/{نقوم بالعملية},
                             180/{4. الكتابة\\\textenglish{Write back}}/brandPurple/{نحفظ النتيجة}} {
        \node[unit=\c, minimum width=2.9cm, minimum height=1.2cm] (n\a) at (\a:2.6) {\t};
    }
    \draw[sahmT] (n90) to[bend left=25] (n0);
    \draw[sahmT] (n0) to[bend left=25] (n-90);
    \draw[sahmT] (n-90) to[bend left=25] (n180);
    \draw[sahmT] (n180) to[bend left=25] (n90);
    \node[font=\small, text width=3cm, align=center] at (0,0) {ثم نعيد\\مع التعليمة التالية};
\end{tikzpicture}
\end{center}

\begin{enumerate}
    \item \textbf{الجلب}: يُرسل \textenglish{CO} إلى \textenglish{RAM}، نقرأ الذاكرة، تصل التعليمة إلى
          \textenglish{RIM} ثم إلى \textenglish{RI}. ثم يزيد \textenglish{CO} ليشير إلى التعليمة التالية.
    \item \textbf{فك الترميز}: يقرأ المفكك الرمز العملي الموجود في \textenglish{RI} ويحدد المعاملات.
    \item \textbf{التنفيذ}: إن احتجنا معطاة من الذاكرة نقرؤها، ثم تحسب \textenglish{UAL} النتيجة.
    \item \textbf{الكتابة}: نحفظ النتيجة في سجل أو في الذاكرة، ونحدّث الأعلام.
\end{enumerate}

\subsection{مثال كامل خطوة بخطوة}
نأخذ آلة بسيطة بمركِّم \textenglish{ACC}، ونريد حساب \textenglish{z = x + y}. المتغير \textenglish{x}
في العنوان 20، و\textenglish{y} في 21، و\textenglish{z} في 22.

\begin{center}
\begin{tabular}{|c|l|l|}
    \hline
    \ras{العنوان} & \ras{المحتوى} & \ras{المعنى} \\ \hline
    10 & \code{LOAD 20}  & $\text{ACC} \leftarrow [20]$ \\ \hline
    11 & \code{ADD 21}   & $\text{ACC} \leftarrow \text{ACC} + [21]$ \\ \hline
    12 & \code{STORE 22} & $[22] \leftarrow \text{ACC}$ \\ \hline
    20 & 7 & \textenglish{x} \\ \hline
    21 & 5 & \textenglish{y} \\ \hline
    22 & 0 & \textenglish{z} (قبل التنفيذ) \\ \hline
\end{tabular}
\end{center}

نتابع التعليمة الثانية \code{ADD 21} خطوة بخطوة (في البداية \textenglish{CO = 11} و\textenglish{ACC = 7}):

\begin{center}
\small
\begin{tabularx}{\linewidth}{|>{\centering\arraybackslash}p{1.6cm}|R|c|c|c|c|c|}
    \hline
    \ras{المرحلة} & \ras{ماذا يحدث؟} & \ras{\textenglish{CO}} & \ras{\textenglish{RAM}} & \ras{\textenglish{RIM}} & \ras{\textenglish{RI}} & \ras{\textenglish{ACC}} \\ \hline
    الجلب & \nql{RAM}{CO} & 11 & \cell{orangeDeep}{11} & -- & -- & 7 \\ \cline{2-7}
     & قراءة الذاكرة: \nql{RIM}{[11]} & 11 & 11 & \cell{orangeDeep}{\code{ADD 21}} & -- & 7 \\ \cline{2-7}
     & \nql{RI}{RIM}، و\nql{CO}{CO + 1} & \cell{orangeDeep}{12} & 11 & \code{ADD 21} & \cell{orangeDeep}{\code{ADD 21}} & 7 \\ \hline
    فك الترميز & المفكك يفهم: «جمع مع محتوى العنوان 21» & 12 & 11 & \code{ADD 21} & \code{ADD 21} & 7 \\ \hline
    التنفيذ & \nql{RAM}{21}، ثم قراءة: \nql{RIM}{[21]} & 12 & \cell{orangeDeep}{21} & \cell{orangeDeep}{5} & \code{ADD 21} & 7 \\ \cline{2-7}
     & \textenglish{UAL}: $7 + 5$ & 12 & 21 & 5 & \code{ADD 21} & 7 \\ \hline
    الكتابة & \nql{ACC}{12}، وتحديث الأعلام & 12 & 21 & 5 & \code{ADD 21} & \cell{orangeDeep}{\textbf{12}} \\ \hline
\end{tabularx}
\end{center}

\begin{khulasa}
    بعد هذه التعليمة: \textenglish{ACC = 12} و\textenglish{CO = 12}، فالتعليمة القادمة هي \code{STORE 22}
    التي تكتب 12 في الذاكرة، فتصبح \textenglish{z = 12}.
\end{khulasa}

\begin{hal}
    المعالجات الحديثة لا تنتظر انتهاء تعليمة لتبدأ الأخرى: بينما تُنفَّذ تعليمة، تُفك ترميز
    التي بعدها وتُجلب الثالثة. هذا هو \textbf{خط الأنابيب} \textenglish{(pipeline)} الذي سندرسه في الفصل الرابع.
\end{hal}

%================================================
\section{تمارين}
%================================================

\begin{tamrin}
    احسب بالمجمّع الكامل (بتًا بتًا) $1011 + 0110$ على 4 بتات. ما النتيجة وما قيمة \textenglish{CF}؟
\end{tamrin}
\begin{hall}
    \begin{tabular}{|c|c|c|c|c|c|}
        \hline \ras{البت} & \ras{$A$} & \ras{$B$} & \ras{$C_{in}$} & \ras{$S$} & \ras{$C_{out}$} \\ \hline
        0 & 1 & 0 & 0 & 1 & 0 \\ \hline 1 & 1 & 1 & 0 & 0 & 1 \\ \hline
        2 & 0 & 1 & 1 & 0 & 1 \\ \hline 3 & 1 & 0 & 1 & 0 & 1 \\ \hline
    \end{tabular}\\
    النتيجة $0001$ و\textenglish{CF = 1}: $11 + 6 = 17 = 16 + 1$، و«الـ16» ذهبت إلى \textenglish{CF}.
\end{hall}

\begin{tamrin}
    على 8 بتات، أعط قيم \textenglish{AL} و\textenglish{CF} و\textenglish{ZF} بعد: (أ) $255 + 1$ (ب) $10 - 10$ (ج) $10 - 20$.
\end{tamrin}
\begin{hall}
    (أ) $256 - 256 = 0$: \textenglish{AL = 0}، \textenglish{CF = 1}، \textenglish{ZF = 1}.\\
    (ب) \textenglish{AL = 0}، \textenglish{CF = 0}، \textenglish{ZF = 1}.\\
    (ج) $10 - 20 = -10$، وعلى 8 بتات $256 - 10 = 246 = \hex{F6}$: \textenglish{CF = 1} (استلاف)، \textenglish{ZF = 0}.
\end{hall}

\begin{tamrin}
    حدد نمط العنونة في كل تعليمة: (أ) \code{mov al, [bx]} (ب) \code{mov al, 10h} (ج) \code{mov al, [10h]} (د) \code{mov al, bl}.
\end{tamrin}
\begin{hall}
    (أ) غير مباشر بالسجل \quad (ب) فوري \quad (ج) مباشر \quad (د) سجلي.
\end{hall}

\begin{tamrin}
    في المثال السابق (\textenglish{z = x + y})، اكتب جدول المراحل للتعليمة الثالثة \code{STORE 22}.
\end{tamrin}
\begin{hall}
    \textbf{الجلب}: \nql{RAM}{12}، \nql{RIM}{STORE 22}، \nql{RI}{RIM}، \textenglish{CO = 13}.
    \textbf{فك الترميز}: «اكتب \textenglish{ACC} في العنوان 22».
    \textbf{التنفيذ والكتابة}: \nql{RAM}{22}، \nql{RIM}{ACC} (12)، ثم إشارة كتابة: $[22] = 12$.
\end{hall}

\chapter{دراسة حالة: المعالج \textenglish{Intel 8086}}

\begin{khulasa}[: ماذا سنتعلم في هذا الفصل؟]
\begin{itemize}
    \item عائلة \textenglish{80x86}، والبنية الداخلية للمعالج \textenglish{8086}.
    \item سجلات المعالج \textenglish{8086}، والعنونة المقطعية (كيف نحسب العنوان الفيزيائي).
    \item كتابة أول برنامج، وكيف تُنفَّذ تعليمة داخل المعالج.
    \item التعليمات: النقل، الحساب، المنطق، القفز، الحلقات، ثم كيف تتحول التعليمة إلى لغة الآلة.
    \item الإدخال والإخراج، المكدس، الإجراءات، والمقاطعات.
\end{itemize}
\end{khulasa}

%================================================
\section{عائلة \textenglish{80x86}}
%================================================

في سنة 1978 أطلقت شركة \textenglish{Intel} المعالج \textenglish{8086}. نجح نجاحًا كبيرًا لأن شركة
\textenglish{IBM} اختارت أخاه \textenglish{8088} لأول حاسوب شخصي \textenglish{(IBM PC)} سنة 1981. ومن
يومها، كل معالج جديد من \textenglish{Intel} و\textenglish{AMD} يبقى قادرًا على تنفيذ برامج \textenglish{8086}.
هذه العائلة تسمى \textenglish{x86}.

\begin{tarif}
    \mustalah{التوافق مع الأجيال السابقة}{Compatibilité ascendante}{Backward compatibility}: المعالج الجديد يفهم
    كل تعليمات المعالجات القديمة من عائلته. برنامج كُتب للمعالج \textenglish{8086} سنة 1980 يمكن أن يعمل
    على معالج \textenglish{Core i7} حديث. هذا ما جعل عائلة \textenglish{x86} تبقى أكثر من 45 سنة.
\end{tarif}

\begin{hal}
    الفرق الوحيد بين \textenglish{8086} و\textenglish{8088} هو \textbf{حافلة المعطيات الخارجية}: 16 بتًا في \textenglish{8086}
    و8 بتات فقط في \textenglish{8088}. من الداخل هما متطابقان. اختارت \textenglish{IBM} المعالج \textenglish{8088} لأن
    الدارات المساعدة ذات 8 بتات كانت أرخص في ذلك الوقت.
\end{hal}

\begin{center}
\small
\begin{tabular}{|c|c|c|c|l|}
    \hline
    \ras{المعالج} & \ras{السنة} & \ras{عدد البتات} & \ras{الذاكرة القصوى} & \ras{الجديد فيه} \\ \hline
    \textenglish{8086 / 8088} & 1978 & 16 & \textenglish{1 MB} & أول معالج من العائلة \\ \hline
    \textenglish{80286} & 1982 & 16 & \textenglish{16 MB} & النمط المحمي \\ \hline
    \textenglish{80386} & 1985 & 32 & \textenglish{4 GB} & أول معالج 32 بتًا \textenglish{(IA-32)} \\ \hline
    \textenglish{80486} & 1989 & 32 & \textenglish{4 GB} & ذاكرة مخبئية ووحدة الأعداد الحقيقية داخله \\ \hline
    \textenglish{Pentium} & 1993 & 32 & \textenglish{4 GB} & تنفيذ تعليمتين معًا \\ \hline
    \textenglish{Athlon 64} \textenglish{(AMD)} & 2003 & 64 & $2^{48}$ عمليًا & المعمارية \textenglish{AMD64 (x86-64)} \\ \hline
    \textenglish{Core i7} وما بعده & 2008 & 64 & ضخمة & عدة أنوية في الشريحة نفسها \\ \hline
\end{tabular}
\end{center}

\begin{tarif}
    \textbf{الخصائص العامة للمعالج \textenglish{8086}}:
    \begin{itemize}
        \item معالج \textbf{16 بتًا}: سجلاته وحافلة معطياته من 16 بتًا.
        \item حافلة عناوين من \textbf{20 بتًا}: يستطيع عنونة \textenglish{1 MB} من الذاكرة.
        \item حوالي \textenglish{29\,000} ترانزستور، وتردد بين \textenglish{5} و\textenglish{10 MHz}.
        \item يملك حوالي \textbf{100 تعليمة} مختلفة (معالج \textenglish{CISC}).
    \end{itemize}
\end{tarif}

%================================================
\section{البنية الداخلية للمعالج \textenglish{8086}}
%================================================

ينقسم المعالج \textenglish{8086} إلى وحدتين تعملان \textbf{في الوقت نفسه}:

\begin{center}
\begin{tikzpicture}[r/.style={reg, minimum width=1.25cm, font=\codefont\small}]
    % BIU
    \draw[brandOrange, very thick, rounded corners=6pt, fill=brandOrange!8] (0.3,-3.2) rectangle (6.3,2.6);
    \node[font=\bfseries, lawn=orangeDeep] at (3.3,2.25) {وحدة الربط بالحافلة \textenglish{(BIU)}};
    \node[r] at (4.6,1.4) {\textenglish{CS}}; \node[r] at (4.6,0.8) {\textenglish{DS}};
    \node[r] at (4.6,0.2) {\textenglish{SS}}; \node[r] at (4.6,-0.4) {\textenglish{ES}};
    \node[r] at (4.6,-1.0) {\textenglish{IP}};
    \node[draw=brandDark, fill=brandOrange!25, thick, minimum width=1.4cm, minimum height=0.6cm, font=\small] at (2,1.1) {مجمّع $\Sigma$};
    \node[font=\scriptsize, align=center] at (2,0.35) {يحسب العنوان\\الفيزيائي};
    \foreach \i in {0,...,5} \node[draw=brandDark, fill=white, minimum width=0.55cm, minimum height=0.5cm] at (1.2+0.6*\i,-2.1) {};
    \node[font=\scriptsize] at (2.7,-2.75) {طابور التعليمات (6 بايتات)};
    % EU
    \draw[brandTeal, very thick, rounded corners=6pt, fill=brandTeal!8] (-6.6,-3.2) rectangle (-0.3,2.6);
    \node[font=\bfseries, lawn=tealDeep] at (-3.45,2.25) {وحدة التنفيذ \textenglish{(EU)}};
    \foreach \n/\y in {AX/1.4, BX/0.8, CX/0.2, DX/-0.4} \node[r] at (-1.3,\y) {\textenglish{\n}};
    \foreach \n/\y in {SI/1.4, DI/0.8, BP/0.2, SP/-0.4} \node[r] at (-2.7,\y) {\textenglish{\n}};
    \draw[brandGreen!70!black, very thick, fill=brandGreen!20] (-5.9,1.2) -- (-4.1,1.2) -- (-4.5,-0.2) -- (-5.5,-0.2) -- cycle;
    \node[font=\bfseries\small] at (-5,0.5) {\textenglish{UAL}};
    \node[r, minimum width=2.5cm, font=\small] at (-5,-1.0) {الأعلام \textenglish{FLAGS}};
    \node[unit=brandOrange, font=\small, minimum width=2.6cm] at (-2.2,-2.1) {وحدة التحكم};
    % links
    \draw[sahmT, <->] (-0.3,-2.1) -- (0.3,-2.1);
    \draw[brandDark, line width=4pt, {Stealth[length=4mm]}-{Stealth[length=4mm]}] (6.3,0.2) -- (7.8,0.2);
    \node[font=\small, align=center] at (7.1,-0.5) {حافلة\\النظام};
\end{tikzpicture}
\end{center}

\begin{itemize}
    \item \mustalah{وحدة الربط بالحافلة}{Bus Interface Unit}{BIU}: تتكلم مع الذاكرة. تحسب العناوين،
          تجلب التعليمات مسبقًا وتضعها في \textbf{طابور} من 6 بايتات، وتقرأ وتكتب المعطيات.
    \item \mustalah{وحدة التنفيذ}{Execution Unit}{EU}: تأخذ التعليمات من الطابور، تفك ترميزها، وتنفذها
          بوحدة الحساب والمنطق.
\end{itemize}

\begin{tashbih}
    في مطعم: \textbf{النادل} (\textenglish{BIU}) يأتي بالطلبات من الزبائن ويضعها في صف أمام الطباخ،
    بينما \textbf{الطباخ} (\textenglish{EU}) يحضّر الطلب الحالي. الطباخ لا ينتظر النادل، لأن الطلب القادم
    موجود دائمًا في الصف. هذه أول فكرة لخط الأنابيب!
\end{tashbih}

%================================================
\section{سجلات المعالج \textenglish{8086}}
%================================================

يملك المعالج \textenglish{8086} \textbf{14 سجلًا} من 16 بتًا.

\subsection{السجلات العامة}
أربعة سجلات للمعطيات، ويمكن تقسيم كل واحد منها إلى نصفين من 8 بتات:

\begin{center}
\begin{tikzpicture}[h/.style={draw=brandDark, thick, minimum width=2.2cm, minimum height=0.7cm, font=\codefont\bfseries}]
    \foreach \r/\n/\d [count=\i from 0] in {A/AX/{المركِّم: الحساب، الضرب، القسمة، الإدخال والإخراج},
                                         B/BX/{القاعدة: يحفظ عنوانًا في الذاكرة},
                                         C/CX/{العدّاد: عدد دورات الحلقات},
                                         D/DX/{المعطيات: مع \textenglish{AX} في الضرب والقسمة}} {
        \node[h, fill=brandTeal!25] (hi\i) at (-2.2,-0.9*\i) {\textenglish{\r H}};
        \node[h, fill=brandTeal!10] (lo\i) at (0,-0.9*\i) {\textenglish{\r L}};
        \node[font=\codefont\bfseries, lawn=tealDeep] at (1.6,-0.9*\i) {\textenglish{\n}};
        \node[font=\small, anchor=east] at (-3.6,-0.9*\i) {\d};
    }
    \node[font=\scriptsize] at (-2.2,0.6) {البتات 15 إلى 8};
    \node[font=\scriptsize] at (0,0.6) {البتات 7 إلى 0};
\end{tikzpicture}
\end{center}

\begin{mithal}
    إذا كان \textenglish{AX = \hex{1234}} فإن \textenglish{AH = \hex{12}} و\textenglish{AL = \hex{34}}.
    إذا غيّرنا \textenglish{AL} إلى \hex{FF} يصبح \textenglish{AX = \hex{12FF}}.
\end{mithal}

لماذا نقسم السجلات؟ لأننا نعمل كثيرًا بقيم صغيرة (حرف، رقم بين 0 و255). بدل أن نضيّع سجلًا كاملًا من
16 بتًا لحرف واحد، نستطيع وضع حرفين في سجل واحد: واحد في \textenglish{AH} وآخر في \textenglish{AL}.

\begin{intabih}
    هذه السجلات «عامة» لأننا نستطيع استعمالها في أي حساب، لكن بعض التعليمات تفرض سجلًا معينًا:
    \code{loop} تستعمل دائمًا \textenglish{CX}، و\code{mul} و\code{div} تستعملان دائمًا \textenglish{AX} و\textenglish{DX}،
    ووظائف \code{int 21h} تنتظر رقم الوظيفة في \textenglish{AH}. لهذا يجب أن نحفظ هذه الأدوار.
\end{intabih}

\subsection{سجلات المؤشرات والفهارس}
\begin{center}
\begin{tabular}{|c|l|l|}
    \hline
    \ras{السجل} & \ras{الاسم} & \ras{الدور} \\ \hline
    \textenglish{SI} & \textenglish{Source Index} & فهرس المصدر (للجداول والسلاسل) \\ \hline
    \textenglish{DI} & \textenglish{Destination Index} & فهرس الوجهة \\ \hline
    \textenglish{SP} & \textenglish{Stack Pointer} & يشير إلى قمة المكدس \\ \hline
    \textenglish{BP} & \textenglish{Base Pointer} & للوصول إلى معطيات داخل المكدس \\ \hline
\end{tabular}
\end{center}

\begin{mithal}
    لنمر على جدول من 10 عناصر، نضع عنوان بدايته في \textenglish{SI}، ثم نزيد \textenglish{SI} في كل دورة
    (\code{inc si}): في الدورة الأولى يشير \textenglish{SI} إلى العنصر الأول، وفي الثانية إلى الثاني، وهكذا.
    \textenglish{SI} هنا يلعب دور \textbf{الإصبع} الذي نضعه على العنصر الذي نقرؤه الآن.
\end{mithal}

\subsection{سجلات المقاطع}
\begin{center}
\begin{tabular}{|c|l|l|}
    \hline
    \ras{السجل} & \ras{الاسم} & \ras{يشير إلى بداية\dots} \\ \hline
    \textenglish{CS} & \textenglish{Code Segment} & مقطع البرنامج (التعليمات) \\ \hline
    \textenglish{DS} & \textenglish{Data Segment} & مقطع المعطيات (المتغيرات) \\ \hline
    \textenglish{SS} & \textenglish{Stack Segment} & مقطع المكدس \\ \hline
    \textenglish{ES} & \textenglish{Extra Segment} & مقطع إضافي (نسخ السلاسل مثلًا) \\ \hline
\end{tabular}
\end{center}

\subsection{مؤشر التعليمة وسجل الأعلام}
\begin{itemize}
    \item \textbf{\textenglish{IP}} \textenglish{(Instruction Pointer)}: عنوان التعليمة القادمة داخل مقطع البرنامج.
          هو \textenglish{CO} الذي رأيناه في الفصل الثاني. لا يمكن تغييره مباشرة بـ\code{mov}، بل بالقفز والاستدعاء.
    \item \textbf{\textenglish{FLAGS}}: 9 بتات مستعملة من أصل 16:
\end{itemize}

\begin{center}
\begin{tikzpicture}[b/.style={draw=brandDark, minimum width=0.75cm, minimum height=0.7cm, font=\codefont\scriptsize\bfseries}]
    \foreach \n [count=\i from 0] in {{}, {}, {}, {}, OF, DF, IF, TF, SF, ZF, {}, AF, {}, PF, {}, CF} {
        \pgfmathtruncatemacro\bit{15-\i}
        \node[b, fill=brandTeal!15] at (0.75*\i, 0) {\textenglish{\n}};
        \node[font=\codefont\tiny] at (0.75*\i, 0.55) {\textenglish{\bit}};
    }
\end{tikzpicture}
\end{center}

\begin{center}
\small
\begin{tabularx}{\linewidth}{|c|R||c|R|}
    \hline
    \ras{العلم} & \ras{المعنى} & \ras{العلم} & \ras{المعنى} \\ \hline
    \textenglish{CF} & حمل أو استلاف خارج من آخر بت & \textenglish{OF} & فيضان: نتيجة خاطئة للأعداد بإشارة \\ \hline
    \textenglish{PF} & عدد الآحاد في البايت الأول زوجي & \textenglish{DF} & اتجاه السلاسل (0: تصاعدي) \\ \hline
    \textenglish{AF} & حمل من البت 3 إلى البت 4 & \textenglish{IF} & السماح بالمقاطعات العتادية \\ \hline
    \textenglish{ZF} & النتيجة صفر & \textenglish{TF} & التنفيذ خطوة بخطوة (للتصحيح) \\ \hline
    \textenglish{SF} & النتيجة سالبة (البت الأخير 1) & & \\ \hline
\end{tabularx}
\end{center}

%================================================
\section{العنونة المقطعية}
%================================================

\subsection{المشكلة}
سجلات \textenglish{8086} من 16 بتًا، فهي تستطيع كتابة عناوين من 0 إلى $2^{16} - 1 = 65535$ فقط
(أي \textenglish{64 KB}). لكن حافلة العناوين من 20 بتًا (\textenglish{1 MB}). كيف نكتب عنوانًا من 20 بتًا
بسجلات من 16 بتًا؟

\subsection{الحل: المقطع والإزاحة}
نستعمل \textbf{سجلين}: سجل \textbf{المقطع} يحدد بداية منطقة من الذاكرة، و\textbf{الإزاحة}
\textenglish{(offset)} تحدد المسافة داخل هذه المنطقة. نكتب العنوان المنطقي هكذا:
\textenglish{segment:offset}.

\begin{tarif}
    \centering
    $\text{العنوان الفيزيائي} = \text{المقطع} \times 16 + \text{الإزاحة}$
    \par\small
    الضرب في 16 = إضافة صفر واحد على اليمين في الست عشري.
\end{tarif}

\begin{khatwa}[: العنوان الفيزيائي لـ \textenglish{CS = 1000H} و\textenglish{IP = 0234H}]
    \begin{enumerate}
        \item نضرب المقطع في 16 (نضيف صفرًا): \hex{1000} يصبح \hex{10000}.
        \item نضيف الإزاحة:
        \begin{center}
        \textenglish{\codefont
        \begin{tabular}{r}
              10000 \\
            + \ 0234 \\ \hline
              10234
        \end{tabular}}
        \end{center}
        \item العنوان الفيزيائي هو \hex{10234}، وهو عنوان من 20 بتًا (5 أرقام ست عشرية).
    \end{enumerate}
\end{khatwa}

\subsection{لماذا نضرب في 16؟}
لو استعملنا المقطع كما هو (دون ضرب)، لما تجاوزنا العنوان $\hex{FFFF} + \hex{FFFF} = \hex{1FFFE}$، أي
حوالي \textenglish{128 KB} فقط. بالضرب في 16، يصبح المقطع قادرًا على الإشارة إلى أي مكان في \textenglish{1 MB}:
$\hex{FFFF} \times 16 = \hex{FFFF0}$. ولهذا يبدأ كل مقطع دائمًا عند عنوان يقبل القسمة على 16، ويسمى
هذا الجزء من 16 بايتًا \mustalah{فقرة}{Paragraphe}{Paragraph}.

\begin{khatwa}[: ما المنطقة التي يغطيها المقطع \textenglish{DS = 2000H}؟]
    \begin{enumerate}
        \item أصغر إزاحة هي \hex{0000}: $\hex{20000} + \hex{0000} = \hex{20000}$.
        \item أكبر إزاحة هي \hex{FFFF}: $\hex{20000} + \hex{FFFF} = \hex{2FFFF}$.
        \item المقطع يغطي المنطقة من \hex{20000} إلى \hex{2FFFF}، أي $2^{16} = 65536$ بايتًا = \textenglish{64 KB}.
    \end{enumerate}
\end{khatwa}

\begin{tashbih}
    تخيّل مدينة فيها شوارع كثيرة. العنوان الكامل لبيت هو «\textbf{الشارع} رقم 12، \textbf{المنزل} رقم 34».
    سجل المقطع هو رقم الشارع (يحدد بداية منطقة)، والإزاحة هي رقم المنزل داخل الشارع. ساعي البريد يحتاج
    المعلومتين معًا ليجد البيت.
\end{tashbih}

\begin{center}
\begin{tikzpicture}
    \draw[thick, fill=brandLight] (0,0) rectangle (2.5,-7);
    \node[font=\small\bfseries] at (1.25,0.35) {الذاكرة \textenglish{(1 MB)}};
    \node[font=\codefont\scriptsize, anchor=west] at (2.6,-0.1) {\textenglish{00000}};
    \node[font=\codefont\scriptsize, anchor=west] at (2.6,-6.9) {\textenglish{FFFFF}};
    \fill[brandTeal!30] (0,-2) rectangle (2.5,-4.5);
    \draw[thick] (0,-2) rectangle (2.5,-4.5);
    \node[font=\small] at (1.25,-3.4) {المقطع \textenglish{(64 KB)}};
    \draw[sahmO] (-1.5,-2) node[left, font=\small] {بداية المقطع: \textenglish{10000}} -- (0,-2);
    \draw[sahmT] (-1.5,-2.8) node[left, font=\small] {العنوان \textenglish{10234}} -- (0,-2.8);
    \draw[decorate, decoration={brace, amplitude=5pt}, thick] (2.6,-2) -- (2.6,-2.8)
        node[midway, right=6pt, font=\small] {الإزاحة \textenglish{0234}};
\end{tikzpicture}
\end{center}

\begin{intabih}
    عنوان فيزيائي واحد يمكن كتابته بعدة طرق: \textenglish{1000:0234} و\textenglish{1023:0004} يعطيان
    العنوان نفسه \hex{10234}. تحقق بنفسك: $\hex{10230} + \hex{0004} = \hex{10234}$.
\end{intabih}

\subsection{الأزواج الافتراضية}
كل سجل إزاحة يُستعمل افتراضيًا مع سجل مقطع معين:

\begin{center}
\begin{tabular}{|c|c|l|}
    \hline
    \ras{الإزاحة} & \ras{المقطع الافتراضي} & \ras{الاستعمال} \\ \hline
    \textenglish{IP} & \textenglish{CS} & جلب التعليمات \\ \hline
    \textenglish{SP}، \textenglish{BP} & \textenglish{SS} & المكدس \\ \hline
    \textenglish{BX}، \textenglish{SI}، عنوان مباشر & \textenglish{DS} & المتغيرات \\ \hline
    \textenglish{DI} (في تعليمات السلاسل) & \textenglish{ES} & وجهة النسخ \\ \hline
\end{tabular}
\end{center}

\begin{hal}
    في برامج \textenglish{.COM} التي نكتبها في هذا الفصل، تكون المقاطع الأربعة \textenglish{CS = DS = SS = ES}
    متساوية: البرنامج كله داخل مقطع واحد من \textenglish{64 KB}، ويبدأ عند الإزاحة \hex{100}. هذا
    معنى السطر \code{org 100h} الذي سنجده في كل برامجنا.
\end{hal}

%================================================
\section{أنماط العنونة في \textenglish{8086}}
%================================================

\begin{tarif}
    \mustalah{العنوان الفعلي}{Adresse effective, EA}{Effective Address}: الإزاحة التي يحسبها
    المعالج لمعامل في الذاكرة. ثم يحسب العنوان الفيزيائي $= \text{DS} \times 16 + \text{EA}$.
\end{tarif}

\begin{center}
\small
\begin{tabularx}{\linewidth}{|>{\centering\arraybackslash}p{2.6cm}|>{\centering\arraybackslash}p{3.4cm}|R|}
    \hline
    \ras{النمط} & \ras{مثال} & \ras{العنوان الفعلي \textenglish{EA}} \\ \hline
    فوري & \code{mov ax, 5} & لا يوجد: القيمة في التعليمة \\ \hline
    سجلي & \code{mov ax, bx} & لا يوجد: القيمة في سجل \\ \hline
    مباشر & \code{mov ax, [0200h]} & \textenglish{EA} = \hex{0200} \\ \hline
    غير مباشر بسجل & \code{mov ax, [bx]} & \textenglish{EA} = \textenglish{BX} (أو \textenglish{SI}، \textenglish{DI}، \textenglish{BP}) \\ \hline
    قاعدي & \code{mov ax, [bx+4]} & \textenglish{EA} = \textenglish{BX} + 4 \\ \hline
    مفهرس & \code{mov ax, [si+4]} & \textenglish{EA} = \textenglish{SI} + 4 \\ \hline
    قاعدي مفهرس & \code{mov ax, [bx+si+4]} & \textenglish{EA} = \textenglish{BX} + \textenglish{SI} + 4 \\ \hline
\end{tabularx}
\end{center}

\begin{khatwa}[: أين يقرأ \code{mov ax, [bx+si+4]}؟]
    نفترض: \textenglish{DS = \hex{2000}}، \textenglish{BX = \hex{0100}}، \textenglish{SI = \hex{0020}}.
    \begin{enumerate}
        \item العنوان الفعلي: $\text{EA} = \hex{0100} + \hex{0020} + 4 = \hex{0124}$.
        \item العنوان الفيزيائي: $\hex{2000} \times 16 + \hex{0124} = \hex{20000} + \hex{0124} = \hex{20124}$.
        \item يقرأ المعالج كلمة (بايتين): البايت في \hex{20124} يذهب إلى \textenglish{AL}، والبايت في
              \hex{20125} يذهب إلى \textenglish{AH}.
    \end{enumerate}
\end{khatwa}

\begin{khatwa}[: \code{mov ax, [bp+2]} يستعمل المقطع \textenglish{SS}]
    نفترض \textenglish{SS = \hex{3000}} و\textenglish{BP = \hex{FFF0}}.
    \begin{enumerate}
        \item العنوان الفعلي: $\text{EA} = \hex{FFF0} + 2 = \hex{FFF2}$.
        \item بما أننا استعملنا \textenglish{BP}، فالمقطع الافتراضي هو \textenglish{SS} وليس \textenglish{DS}.
        \item العنوان الفيزيائي: $\hex{30000} + \hex{FFF2} = \hex{3FFF2}$.
    \end{enumerate}
    هذا مهم في الإجراءات: \textenglish{BP} يُستعمل لقراءة المعاملات الموجودة في المكدس.
\end{khatwa}

\begin{intabih}
    المعالج \textenglish{8086} يحفظ الكلمة في الذاكرة بترتيب \textbf{البايت الصغير أولًا}
    \textenglish{(little endian)}: الكلمة \hex{1234} تُخزن \textenglish{34} ثم \textenglish{12}. لا يُسمح أيضًا
    بمعاملين في الذاكرة في التعليمة نفسها: \code{mov [a], [b]} خطأ.
\end{intabih}

%================================================
\section{أول برنامج}
%================================================

\subsection{الأدوات}
نكتب البرنامج في \textbf{\textenglish{EMU8086}}: \textenglish{File $\rightarrow$ New $\rightarrow$}
\textenglish{COM template}، ثم نضغط \textenglish{Emulate} وننفذ تعليمة بتعليمة بالزر \textenglish{Single step}.

يمر كل برنامج بلغة التجميع بثلاث مراحل:
\begin{center}
\begin{tikzpicture}[node distance=1.1cm]
    \node[unit=brandTeal, minimum width=2.8cm] (src) {الملف المصدر\\\code{hello.asm}};
    \node[unit=brandOrange, minimum width=2.8cm, left=of src] (bin) {الملف التنفيذي\\\code{hello.com}};
    \node[unit=brandGreen, minimum width=2.8cm, left=of bin] (run) {التنفيذ\\في الذاكرة};
    \draw[sahmT] (src) -- node[above, font=\scriptsize] {المجمِّع} (bin);
    \draw[sahmT] (bin) -- node[above, font=\scriptsize] {التحميل} (run);
\end{tikzpicture}
\end{center}
\begin{enumerate}
    \item \textbf{الكتابة}: نكتب التعليمات في ملف نصي (المصدر).
    \item \textbf{التجميع}: المجمِّع يحوّل كل تعليمة إلى بايتات لغة الآلة، ويحسب عناوين المتغيرات والتسميات.
    \item \textbf{التحميل والتنفيذ}: النظام ينسخ البايتات في الذاكرة ابتداءً من الإزاحة \hex{100}، ويضع
          \textenglish{IP = \hex{100}}، فيبدأ المعالج التنفيذ.
\end{enumerate}

\subsection{البرنامج}
\asm{ch3/hello.asm}
\screen{ch3/hello.out}

\begin{sharh}
    \s{1}{تعليق: كل ما بعد \code{;} لا ينفذه المعالج، هو شرح للقارئ.}
    \s{2}{\code{org 100h}: البرنامج من نوع \textenglish{.COM} يبدأ عند الإزاحة \hex{100}.}
    \s{4}{نضع في \textenglish{DX} عنوان الرسالة \code{msg}.}
    \s{5}{نختار وظيفة نظام \textenglish{DOS} رقم \hex{09}: «اكتب السلسلة الموجودة في \textenglish{DX}».}
    \s{6}{\code{int 21h}: نطلب من \textenglish{DOS} تنفيذ الوظيفة المختارة في \textenglish{AH}.}
    \s{7}{\code{ret}: نهاية البرنامج والعودة إلى النظام.}
    \s{9}{نعرّف المتغير \code{msg} بـ\code{db} (\textenglish{define byte}). العلامة \code{\$} تعني نهاية النص.}
\end{sharh}

\begin{intabih}
    نضع المعطيات \textbf{بعد} \code{ret}. لو وضعناها قبل التعليمات، لحاول المعالج تنفيذ حروف
    الرسالة كأنها تعليمات!
\end{intabih}

\begin{khatwa}[: البرنامج في الذاكرة بعد التحميل]
    هذه هي البايتات الحقيقية التي أنتجها المجمِّع:
    \begin{center}
    \small
    \begin{tabular}{|c|l|l|}
        \hline \ras{الإزاحة} & \ras{البايتات} & \ras{المحتوى} \\ \hline
        \textenglish{0100} & \textenglish{\codefont 8D 16 09 01} & \code{lea dx, msg} (العنوان \hex{0109}) \\ \hline
        \textenglish{0104} & \textenglish{\codefont B4 09} & \code{mov ah, 09h} \\ \hline
        \textenglish{0106} & \textenglish{\codefont CD 21} & \code{int 21h} \\ \hline
        \textenglish{0108} & \textenglish{\codefont C3} & \code{ret} \\ \hline
        \textenglish{0109} & \textenglish{\codefont 53 61 6C 61 6D \dots} & الحروف \code{S a l a m \dots} \\ \hline
    \end{tabular}
    \end{center}
    البرنامج كله 9 بايتات من التعليمات، والرسالة تبدأ مباشرة بعدها في \hex{0109}. لاحظ أن المجمِّع حسب
    عنوان \code{msg} بنفسه ووضعه داخل تعليمة \code{lea}.
\end{khatwa}

\subsection{التعريف بالمتغيرات}
\begin{center}
\begin{tabular}{|l|l|}
    \hline
    \ras{الكتابة} & \ras{المعنى} \\ \hline
    \code{x db 5} & بايت واحد قيمته 5 \\ \hline
    \code{y dw 1234h} & كلمة (بايتان) قيمتها \hex{1234} \\ \hline
    \code{t db 14, 3, 27} & جدول من 3 بايتات \\ \hline
    \code{buf db 10 dup(0)} & 10 بايتات كلها 0 \\ \hline
    \code{s db 'Hi\$'} & سلسلة حروف \\ \hline
\end{tabular}
\end{center}

%================================================
\section{مراحل تنفيذ تعليمة داخل \textenglish{8086}}
%================================================

لنتابع التعليمة \code{mov cx, [si]} من البرنامج \code{transfer.asm} (القسم الآتي). في الذاكرة،
هذه التعليمة مكتوبة بالبايتين \textenglish{8B 0C} في الإزاحة \hex{010C}، و\textenglish{SI = \hex{0112}}.

\begin{center}
\small
\begin{tabularx}{\linewidth}{|>{\centering\arraybackslash}p{1.5cm}|>{\centering\arraybackslash}p{1.2cm}|R|}
    \hline
    \ras{المرحلة} & \ras{الوحدة} & \ras{ما يحدث} \\ \hline
    1 & \textenglish{BIU} & تحسب العنوان الفيزيائي $\text{CS} \times 16 + \text{IP}$، وتجلب البايتين \textenglish{8B 0C} إلى الطابور (غالبًا قبل أن نحتاجهما). \\ \hline
    2 & \textenglish{EU} & تأخذ \textenglish{8B} من الطابور: «نقل كلمة من الذاكرة إلى سجل». ثم \textenglish{0C}: «السجل \textenglish{CX}، والعنوان في \textenglish{SI}». \textenglish{IP} يصبح \hex{010E}. \\ \hline
    3 & \textenglish{EU} & تحسب العنوان الفعلي: \textenglish{EA = SI = \hex{0112}}. \\ \hline
    4 & \textenglish{BIU} & تحسب $\text{DS} \times 16 + \hex{0112}$، ترسله على حافلة العناوين، مع إشارة قراءة. \\ \hline
    5 & \textenglish{BIU} & تصل الكلمة \hex{ABCD} على حافلة المعطيات (\textenglish{CD} ثم \textenglish{AB}). \\ \hline
    6 & \textenglish{EU} & تكتب النتيجة في \textenglish{CX}: \textenglish{CX = \hex{ABCD}}. الأعلام لا تتغير مع \code{mov}. \\ \hline
\end{tabularx}
\end{center}

%================================================
\section{مجموعة التعليمات بالتفصيل}
%================================================

نستعمل هذه الرموز في الجداول: \textenglish{reg} = سجل، \textenglish{mem} = خانة ذاكرة، \textenglish{imm} = قيمة فورية.

\subsection{تعليمات النقل}
\begin{center}
\small
\begin{tabular}{|l|l|l|}
    \hline
    \ras{التعليمة} & \ras{المعنى} & \ras{مثال} \\ \hline
    \code{mov dest, src} & $\text{dest} \leftarrow \text{src}$ & \code{mov ax, bx} \\ \hline
    \code{xchg a, b} & تبادل القيمتين & \code{xchg ax, bx} \\ \hline
    \code{lea reg, mem} & \textenglish{reg} يأخذ عنوان المتغير & \code{lea si, tab} \\ \hline
    \code{push src} / \code{pop dest} & وضع في المكدس / أخذ منه & \code{push ax} \\ \hline
\end{tabular}
\end{center}

\asm{ch3/transfer.asm}

\begin{center}
\small
\begin{tabular}{|c|l|c|c|c|c|}
    \hline
    \ras{السطر} & \ras{التعليمة} & \ras{\textenglish{AX}} & \ras{\textenglish{BX}} & \ras{\textenglish{CX}} & \ras{\textenglish{SI}} \\ \hline
    4 & \code{mov ax, 1234h} & \cell{orangeDeep}{1234} & 0000 & 0000 & 0000 \\ \hline
    5 & \code{mov bl, ah} & 1234 & \cell{orangeDeep}{0012} & 0000 & 0000 \\ \hline
    6 & \code{mov bh, al} & 1234 & \cell{orangeDeep}{3412} & 0000 & 0000 \\ \hline
    7 & \code{xchg ax, bx} & \cell{orangeDeep}{3412} & \cell{orangeDeep}{1234} & 0000 & 0000 \\ \hline
    8 & \code{lea si, num} & 3412 & 1234 & 0000 & \cell{orangeDeep}{0112} \\ \hline
    9 & \code{mov cx, [si]} & 3412 & 1234 & \cell{orangeDeep}{ABCD} & 0112 \\ \hline
    10 & \code{mov [res], ax} & 3412 & 1234 & ABCD & 0112 \\ \hline
\end{tabular}
\end{center}
\begin{sharh}
    \s{5--6}{ننسخ نصفي \textenglish{AX} مقلوبين في \textenglish{BX}: \textenglish{BH = 34} و\textenglish{BL = 12}.}
    \s{8}{\code{lea} تعطي \textbf{العنوان} \hex{0112} وليس القيمة.}
    \s{9}{نقرأ الكلمة الموجودة في العنوان \hex{0112}: \hex{ABCD}.}
    \s{10}{السطر الوحيد الذي يكتب في الذاكرة: \code{res} تصبح \hex{3412}.}
\end{sharh}

\subsection{التعليمات الحسابية}
\begin{center}
\small
\begin{tabularx}{\linewidth}{|l|R|l|}
    \hline
    \ras{التعليمة} & \ras{المعنى} & \ras{الأعلام} \\ \hline
    \code{add d, s} & $d \leftarrow d + s$ & كلها \\ \hline
    \code{adc d, s} & $d \leftarrow d + s + \text{CF}$ (لجمع أعداد كبيرة) & كلها \\ \hline
    \code{sub d, s} & $d \leftarrow d - s$ & كلها \\ \hline
    \code{inc d} & $d \leftarrow d + 1$ & كلها إلا \textenglish{CF} \\ \hline
    \code{dec d} & $d \leftarrow d - 1$ & كلها إلا \textenglish{CF} \\ \hline
    \code{neg d} & $d \leftarrow -d$ (المتمم لاثنين) & كلها \\ \hline
    \code{cmp a, b} & يحسب $a - b$ \textbf{دون حفظ النتيجة} & كلها \\ \hline
    \code{mul s} & ضرب دون إشارة (\code{imul} بإشارة) & \textenglish{CF}، \textenglish{OF} \\ \hline
    \code{div s} & قسمة دون إشارة (\code{idiv} بإشارة) & غير محددة \\ \hline
\end{tabularx}
\end{center}

\asm{ch3/arith.asm}

\begin{center}
\small
\begin{tabular}{|c|l|c|c|c|}
    \hline
    \ras{السطر} & \ras{التعليمة} & \ras{النتيجة} & \ras{بالعشري} & \ras{\textenglish{CF}} \\ \hline
    4 & \code{mov al, 25} & \textenglish{AL = 19} & 25 & 0 \\ \hline
    5 & \code{add al, 17} & \textenglish{AL = 2A} & 42 & 0 \\ \hline
    6 & \code{sub al, 2} & \textenglish{AL = 28} & 40 & 0 \\ \hline
    7 & \code{inc al} & \textenglish{AL = 29} & 41 & 0 \\ \hline
    8 & \code{dec al} & \textenglish{AL = 28} & 40 & 0 \\ \hline
    9 & \code{neg al} & \textenglish{AL = D8} & $-40$ & 1 \\ \hline
    10 & \code{mov bx, 1000} & \textenglish{BX = 03E8} & 1000 & 1 \\ \hline
    11 & \code{add bx, 2500} & \textenglish{BX = 0DAC} & 3500 & 0 \\ \hline
\end{tabular}
\end{center}

\begin{khatwa}[: لماذا $-40$ يُكتب \hex{D8}؟ (المتمم لاثنين)]
    \begin{enumerate}
        \item نكتب 40 بالثنائي: \bin{0010\,1000}.
        \item نقلب كل البتات: \bin{1101\,0111}.
        \item نضيف 1: $\bin{1101\,1000} = \hex{D8}$.
    \end{enumerate}
    تحقق: $40 + (-40) = \hex{28} + \hex{D8} = \hex{100}$، وعلى 8 بتات يبقى \textenglish{00}. صحيح!
\end{khatwa}

\subsubsection{الضرب والقسمة}
المعالج يختار المكان حسب حجم المعامل:
\begin{center}
\small
\begin{tabular}{|c|c|c|c|}
    \hline
    \ras{التعليمة} & \ras{المعامل} & \ras{العملية} & \ras{النتيجة} \\ \hline
    \code{mul bl} & 8 بتات & $\text{AL} \times \text{BL}$ & \textenglish{AX} \\ \hline
    \code{mul cx} & 16 بتًا & $\text{AX} \times \text{CX}$ & \textenglish{DX:AX} (32 بتًا) \\ \hline
    \code{div bl} & 8 بتات & $\text{AX} \div \text{BL}$ & \textenglish{AL} = الحاصل، \textenglish{AH} = الباقي \\ \hline
    \code{div cx} & 16 بتًا & $\text{DX:AX} \div \text{CX}$ & \textenglish{AX} = الحاصل، \textenglish{DX} = الباقي \\ \hline
\end{tabular}
\end{center}

\asm{ch3/muldiv.asm}

\begin{sharh}
    \s{6}{$12 \times 11 = 132 = \hex{84}$: \textenglish{AX = \hex{0084}}.}
    \s{10}{$300 \times 500 = 150000$، وهو أكبر من 65535، فيتوزع على سجلين:
           \textenglish{DX = \hex{0002}} و\textenglish{AX = \hex{49F0}} (لأن $\hex{249F0} = 150000$).}
    \s{14}{$47 = 5 \times 9 + 2$: \textenglish{AL = 09} (الحاصل) و\textenglish{AH = 02} (الباقي)، أي \textenglish{AX = \hex{0209}}.}
\end{sharh}

\begin{intabih}
    قبل \code{div cx} يجب وضع \textenglish{DX = 0} إذا كان العدد المقسوم في \textenglish{AX} فقط، وإلا اعتبر
    المعالج \textenglish{DX} جزءًا من العدد. والقسمة على 0 توقف البرنامج (مقاطعة رقم 0).
\end{intabih}

\subsubsection{جمع أعداد أكبر من 16 بتًا: \code{adc}}
المعالج \textenglish{8086} يجمع 16 بتًا في المرة الواحدة. لجمع عددين من 32 بتًا، نقسم كل عدد إلى كلمتين
(عالية ومنخفضة)، ونجمع الكلمتين المنخفضتين بـ\code{add}، ثم العاليتين بـ\code{adc} التي تضيف الحمل
\textenglish{CF} القادم من الجمع الأول، تمامًا كما نحمل «واحدًا» إلى العمود التالي عند الجمع بالقلم.

\asm{ch3/add32.asm}

\begin{center}
\small
\begin{tabular}{|c|l|c|c|c|}
    \hline
    \ras{السطر} & \ras{التعليمة} & \ras{\textenglish{DX}} & \ras{\textenglish{AX}} & \ras{\textenglish{CF}} \\ \hline
    \textenglish{5--6} & تحميل العدد الأول & 0001 & FFFF & 0 \\ \hline
    7 & \code{add ax, 0001h} & 0001 & \cell{orangeDeep}{0000} & \cell{orangeDeep}{\textbf{1}} \\ \hline
    8 & \code{adc dx, 0000h} & \cell{orangeDeep}{0002} & 0000 & 0 \\ \hline
\end{tabular}
\end{center}
النتيجة \textenglish{DX:AX = 0002:0000} = \hex{00020000}. لو استعملنا \code{add} في السطر 8 لضاع الحمل،
ولحصلنا على \hex{00010000}، وهي نتيجة خاطئة.

\subsection{التعليمات المنطقية}
\begin{center}
\small
\begin{tabularx}{\linewidth}{|l|R|l|}
    \hline
    \ras{التعليمة} & \ras{الاستعمال المعتاد} & \ras{مثال} \\ \hline
    \code{and} & إطفاء بتات (القناع) & \code{and al, 0Fh} \\ \hline
    \code{or} & إشعال بتات & \code{or al, 80h} \\ \hline
    \code{xor} & قلب بتات، أو تصفير سجل & \code{xor cx, cx} \\ \hline
    \code{not} & قلب كل البتات & \code{not al} \\ \hline
    \code{test} & \code{and} دون حفظ النتيجة (فحص بت) & \code{test al, 1} \\ \hline
    \code{shl} / \code{shr} & إزاحة: ضرب أو قسمة على 2 & \code{shl bl, 1} \\ \hline
    \code{rol} / \code{ror} & تدوير: البت الخارج يعود من الجهة الأخرى & \code{rol al, 1} \\ \hline
\end{tabularx}
\end{center}

\asm{ch3/logic.asm}

\begin{center}
\small
\begin{tabular}{|c|l|c|c|}
    \hline
    \ras{السطر} & \ras{التعليمة} & \ras{\textenglish{AL} بالثنائي} & \ras{ست عشري} \\ \hline
    4 & \code{mov al, 10110110b} & \textenglish{\codefont 1011 0110} & \textenglish{B6} \\ \hline
    5 & \code{and al, 00001111b} & \textenglish{\codefont 0000 0110} & \textenglish{06} \\ \hline
    6 & \code{or  al, 10000000b} & \textenglish{\codefont 1000 0110} & \textenglish{86} \\ \hline
    7 & \code{xor al, 11111111b} & \textenglish{\codefont 0111 1001} & \textenglish{79} \\ \hline
    8 & \code{not al} & \textenglish{\codefont 1000 0110} & \textenglish{86} \\ \hline
\end{tabular}
\end{center}

\begin{sharh}
    \s{5}{\code{and} مع 1 يبقي البت، ومع 0 يطفئه: نحتفظ بالبتات الأربعة الأولى فقط.}
    \s{6}{\code{or} مع 1 يشعل البت: البت 7 أصبح 1.}
    \s{10--12}{\textenglish{BL}: $5 \rightarrow 10 \rightarrow 20 \rightarrow 10$. الإزاحة إلى اليسار تضرب في 2، وإلى اليمين تقسم على 2.}
    \s{13}{\code{xor cx, cx}: كل بت مع نفسه يعطي 0، إذن \textenglish{CX = 0}. أسرع وأقصر من \code{mov cx, 0}.}
\end{sharh}

\subsubsection{استعمالات عملية للأقنعة}
\begin{khatwa}[: هل العدد في \textenglish{AL} زوجي؟]
    العدد زوجي إذا كان البت 0 يساوي 0. نكتب \code{test al, 1}: هذه التعليمة تحسب $\text{AL} \cdot 1$ دون أن تغيّر
    \textenglish{AL}.
    \begin{itemize}
        \item إذا كان \textenglish{AL = 6} = \bin{0000\,0110}: البت 0 يساوي 0، النتيجة 0، و\textenglish{ZF = 1}: العدد زوجي.
        \item إذا كان \textenglish{AL = 7} = \bin{0000\,0111}: البت 0 يساوي 1، النتيجة 1، و\textenglish{ZF = 0}: العدد فردي.
    \end{itemize}
    بعدها نكتب \code{jz zawji} للقفز في حالة العدد الزوجي.
\end{khatwa}

\begin{khatwa}[: تحويل حرف صغير إلى كبير بالقناع]
    في جدول \textenglish{ASCII}: \textenglish{\code{'a'} = \hex{61} = \bin{0110\,0001}} و\textenglish{\code{'A'} = \hex{41} = \bin{0100\,0001}}.
    الفرق الوحيد هو البت 5. إذن \code{and al, 11011111b} (أي \code{0DFh}) يطفئ البت 5 فيحوّل أي حرف صغير إلى كبير،
    ويترك الحرف الكبير كما هو. هذه طريقة أخرى غير \code{sub al, 32} التي سنراها في قسم الإدخال.
\end{khatwa}

\subsection{التحكم في التدفق: القفزات}
عادة ينفذ المعالج التعليمات بالترتيب. \textbf{القفز} يغيّر \textenglish{IP} فيذهب التنفيذ إلى مكان آخر.

\begin{center}
\small
\begin{tabularx}{\linewidth}{|l|R|l|}
    \hline
    \ras{التعليمة} & \ras{تقفز إذا\dots} & \ras{الأعلام} \\ \hline
    \code{jmp L} & دائمًا (قفز غير مشروط) & -- \\ \hline
    \code{je} / \code{jz} & متساويان / النتيجة صفر & \textenglish{ZF = 1} \\ \hline
    \code{jne} / \code{jnz} & غير متساويين & \textenglish{ZF = 0} \\ \hline
    \code{jg} / \code{jge} & أكبر / أكبر أو يساوي (\textbf{بإشارة}) & \textenglish{ZF}، \textenglish{SF}، \textenglish{OF} \\ \hline
    \code{jl} / \code{jle} & أصغر / أصغر أو يساوي (\textbf{بإشارة}) & \textenglish{SF}، \textenglish{OF} \\ \hline
    \code{ja} / \code{jae} & أكبر / أكبر أو يساوي (\textbf{دون إشارة}) & \textenglish{CF}، \textenglish{ZF} \\ \hline
    \code{jb} / \code{jbe} & أصغر / أصغر أو يساوي (\textbf{دون إشارة}) & \textenglish{CF}، \textenglish{ZF} \\ \hline
\end{tabularx}
\end{center}

\begin{tashbih}
    \code{cmp} يشبه \textbf{وضع شيئين على ميزان}، والقفز المشروط يشبه قولك: «\textbf{إذا} مالت الكفة اليمنى
    فاذهب إلى اليمين». الميزان (\code{cmp}) يضبط الأعلام، والقفز ينظر إليها.
\end{tashbih}

\subsubsection{مثال: إذا \dots وإلا}
\asm{ch3/ifelse.asm}
\screen{ch3/ifelse.out}

\begin{center}
\begin{tikzpicture}[node distance=0.6cm]
    \node[unit] (c) {\code{cmp al, [b]}};
    \node[diamond, aspect=2.2, draw=brandOrange, fill=brandOrange!15, thick, below=of c, font=\small] (q) {$a > b$؟};
    \node[unit=brandGreen, below right=0.4cm and 0.6cm of q] (y) {\code{lea dx, msg1}};
    \node[unit=brandPurple, below left=0.4cm and 0.6cm of q] (n) {\code{lea dx, msg2}};
    \node[unit, below=1.8cm of q] (p) {\code{int 21h} (الطباعة)};
    \draw[sahm] (c) -- (q);
    \draw[sahm] (q) -| node[pos=0.2, above, font=\small] {نعم: \code{jg first}} (y);
    \draw[sahm] (q) -| node[pos=0.2, above, font=\small] {لا} (n);
    \draw[sahm] (y) |- (p);
    \draw[sahm] (n) |- node[pos=0.7, below, font=\small] {\code{jmp show}} (p);
\end{tikzpicture}
\end{center}

\begin{sharh}
    \s{5}{\code{cmp} يحسب $7 - 12 = -5$ ولا يحفظه، لكنه يضبط الأعلام (\textenglish{SF = 1}).}
    \s{6}{\code{jg}: هل $a > b$؟ لا، فلا نقفز ونكمل في السطر 7.}
    \s{7--8}{نختار الرسالة الثانية، ثم \code{jmp show} \textbf{ضروري} لكي لا ننفذ السطر 10.}
    \s{12--13}{الطباعة مشتركة بين الحالتين.}
\end{sharh}

\begin{intabih}
    ننسى كثيرًا \code{jmp show} في السطر 8! بدونه ينفذ البرنامج \code{lea dx, msg1} أيضًا،
    فيطبع دائمًا الرسالة الأولى.
\end{intabih}

\subsubsection{بإشارة أم دون إشارة؟ \code{jg} أو \code{ja}}
البتات نفسها يمكن قراءتها بطريقتين (رأينا ذلك في الفصل الثاني). \textenglish{AL = \hex{C8}} يساوي 200 دون إشارة،
ويساوي $-56$ بإشارة. لذلك يملك المعالج نوعين من القفزات: \code{ja}/\code{jb} للأعداد دون إشارة،
و\code{jg}/\code{jl} للأعداد بإشارة.

\asm{ch3/signed.asm}
\screen{ch3/signed.out}

\begin{sharh}
    \s{4--6}{\code{ja}: نقارن 200 مع 100 كأعداد موجبة. $200 > 100$، فنقفز.}
    \s{14--16}{\code{jg}: البتات نفسها \hex{C8}، لكن نقرؤها بإشارة: $-56 > 100$؟ لا، فلا نقفز.}
\end{sharh}

\begin{intabih}
    اختر القفزة حسب \textbf{نوع المعطيات}: الأعمار، العناوين، رموز الحروف، الأحجام: دون إشارة (\code{ja}، \code{jb}).
    درجات الحرارة، الأرصدة، الفروق: بإشارة (\code{jg}، \code{jl}).
\end{intabih}

\subsubsection{من الخوارزمية إلى لغة التجميع}
كل بنية تحكم تعلمناها في مقياس الخوارزميات لها «ترجمة» ثابتة في لغة التجميع. الفكرة دائمًا هي نفسها:
\textbf{نقارن} بـ\code{cmp}، ثم \textbf{نقفز} حسب النتيجة.

\newcommand{\stk}[1]{\begin{tabular}[t]{@{}l@{}}#1\end{tabular}}
\begin{center}
\small
\begin{tabular}{|l|l|l|}
    \hline
    \ras{البنية} & \ras{بلغة \textenglish{C}} & \ras{بلغة التجميع} \\ \hline
    إذا & \code{if (a == b) x = 1;} &
        \stk{\code{    cmp a, b}\\ \code{    jne end}\\ \code{    mov x, 1}\\ \code{end:}} \\ \hline
    طالما & \code{while (cx > 0) \{ ... \}} &
        \stk{\code{L:  cmp cx, 0}\\ \code{    jle end}\\ \code{    ...}\\ \code{    jmp L}\\ \code{end:}} \\ \hline
    من أجل & \code{for (cx = 10; cx > 0; cx--)} &
        \stk{\code{    mov cx, 10}\\ \code{L:  ...}\\ \code{    loop L}} \\ \hline
\end{tabular}
\end{center}

\begin{intabih}
    لاحظ في ترجمة \code{if}: نقفز بالشرط \textbf{المعاكس} (\code{jne} بدل \code{je}). نقول للمعالج: «إذا
    لم يكونا متساويين، تجاوز تعليمة الإسناد». هذه الطريقة تجعل البرنامج أقصر.
\end{intabih}

\subsection{الحلقات والتكرار}
\begin{tarif}
    \code{loop L} تساوي تعليمتين: \code{dec cx} ثم «إذا كان \textenglish{CX} $\neq 0$ فاقفز إلى \textenglish{L}».
    إذن \textenglish{CX} هو عدّاد الحلقة.
\end{tarif}

\subsubsection{مجموع الأعداد من 1 إلى 10}
\asm{ch3/loopsum.asm}

\begin{center}
\small
\begin{tabular}{|c|c|c|c|c|c|c|c|c|c|c|}
    \hline
    \ras{الدورة} & 1 & 2 & 3 & 4 & 5 & 6 & 7 & 8 & 9 & 10 \\ \hline
    \ras{\textenglish{CX} قبل \code{add}} & 10 & 9 & 8 & 7 & 6 & 5 & 4 & 3 & 2 & 1 \\ \hline
    \ras{\textenglish{AX} بعد \code{add}} & 10 & 19 & 27 & 34 & 40 & 45 & 49 & 52 & 54 & \cell{orangeDeep}{\textbf{55}} \\ \hline
\end{tabular}
\end{center}
بعد الدورة العاشرة يصبح \textenglish{CX = 0} فتتوقف الحلقة، و\textenglish{AX = \hex{0037} = 55}.
لاحظ أننا نجمع بالعكس: $10 + 9 + \dots + 1$، والنتيجة نفسها.

\begin{intabih}
    إذا دخلنا الحلقة و\textenglish{CX = 0}، فإن \code{loop} تنقص \textenglish{CX} أولًا فيصبح \hex{FFFF}، ثم تكرر
    65535 مرة أخرى! لتجنب ذلك نضع قبل الحلقة \code{jcxz fin}: «إذا كان \textenglish{CX = 0} فاقفز إلى النهاية».
\end{intabih}

\subsubsection{البحث عن أكبر قيمة في جدول}
\asm{ch3/maxarray.asm}

\begin{center}
\small
\begin{tabular}{|c|c|c|c|c|}
    \hline
    \ras{\textenglish{SI} يشير إلى} & \ras{العنصر} & \ras{العنصر $>$ \textenglish{AL}؟} & \ras{\textenglish{AL} بعدها} & \ras{\textenglish{CX} بعد \code{loop}} \\ \hline
    \textenglish{tab[0]} & 14 & لا (يساويه) & 14 & 5 \\ \hline
    \textenglish{tab[1]} & 3 & لا & 14 & 4 \\ \hline
    \textenglish{tab[2]} & 27 & نعم & 27 & 3 \\ \hline
    \textenglish{tab[3]} & 9 & لا & 27 & 2 \\ \hline
    \textenglish{tab[4]} & 31 & نعم & \cell{orangeDeep}{\textbf{31}} & 1 \\ \hline
    \textenglish{tab[5]} & 8 & لا & 31 & 0 \\ \hline
\end{tabular}
\end{center}
النتيجة: \textenglish{AL = \hex{1F} = 31}، وتُحفظ في \code{max}. استعملنا \code{jbe} لأن القيم
\textbf{دون إشارة}.

\subsubsection{التكرار بالبادئة \code{rep}}
تعليمات السلاسل تعمل على \textenglish{DS:SI} (المصدر) و\textenglish{ES:DI} (الوجهة)، والبادئة \code{rep}
تكررها \textenglish{CX} مرة.

\asm{ch3/repmovs.asm}
\screen{ch3/repmovs.out}

\begin{sharh}
    \s{7}{\code{cld}: نجعل \textenglish{DF = 0}، فيزيد \textenglish{SI} و\textenglish{DI} بعد كل نسخ.}
    \s{8}{\code{rep movsb} تنسخ بايتًا من \textenglish{[SI]} إلى \textenglish{[DI]}، وتكرر 6 مرات: الحروف الخمسة و\code{\$}.}
\end{sharh}

\subsection{توليد لغة الآلة}
المجمِّع يحوّل كل تعليمة إلى بايتات. هذه البايتات الحقيقية للبرنامج \code{machine.asm}:

\asm{ch3/machine.asm}

\begin{center}
\small
\begin{tabular}{|c|l|l|}
    \hline
    \ras{العنوان} & \ras{لغة الآلة (ست عشري)} & \ras{التعليمة} \\ \hline
    \textenglish{0100} & \textenglish{\codefont B8 34 12} & \code{mov ax, 1234h} \\ \hline
    \textenglish{0103} & \textenglish{\codefont 89 C3} & \code{mov bx, ax} \\ \hline
    \textenglish{0105} & \textenglish{\codefont 04 05} & \code{add al, 5} \\ \hline
    \textenglish{0107} & \textenglish{\codefont 88 07} & \code{mov [bx], al} \\ \hline
    \textenglish{0109} & \textenglish{\codefont 41} & \code{inc cx} \\ \hline
    \textenglish{010A} & \textenglish{\codefont EB 01} & \code{jmp fin} \\ \hline
    \textenglish{010C} & \textenglish{\codefont 90} & \code{nop} \\ \hline
    \textenglish{010D} & \textenglish{\codefont C3} & \code{ret} \\ \hline
\end{tabular}
\end{center}

\begin{khatwa}[: فك ترميز \textenglish{B8 34 12}]
    \begin{enumerate}
        \item \textenglish{B8} = \bin{1011\,1000}: الخمسة بتات الأولى \textenglish{10111} تعني «\code{mov} قيمة فورية في سجل 16 بتًا»،
              والبتات الثلاثة الأخيرة \textenglish{000} = رقم السجل \textenglish{AX}.
        \item \textenglish{34 12}: القيمة \hex{1234} مكتوبة بالبايت الصغير أولًا \textenglish{(little endian)}.
    \end{enumerate}
\end{khatwa}

\begin{khatwa}[: فك ترميز \textenglish{89 C3} (\code{mov bx, ax})]
    البايت الأول \textenglish{89} = \bin{100010\,0\,1}:
    \begin{itemize}
        \item \textenglish{100010}: الرمز العملي لـ\code{mov} بين سجل وسجل أو ذاكرة.
        \item \textenglish{d = 0}: السجل المذكور في الحقل \textenglish{reg} هو \textbf{المصدر}.
        \item \textenglish{w = 1}: نعمل على 16 بتًا (لو كان 0 لعملنا على 8 بتات).
    \end{itemize}
    البايت الثاني \textenglish{C3} = \bin{11\,000\,011}، ويسمى \textenglish{ModR/M}:
    \begin{center}
    \begin{tikzpicture}[f/.style={draw=brandDark, thick, minimum height=0.7cm, font=\codefont\bfseries}]
        \node[f, minimum width=1.6cm, fill=brandOrange!20] (m) at (0,0) {\textenglish{11}};
        \node[f, minimum width=2.2cm, fill=brandTeal!20, anchor=west] (r) at (m.east) {\textenglish{000}};
        \node[f, minimum width=2.2cm, fill=brandGreen!20, anchor=west] (rm) at (r.east) {\textenglish{011}};
        \node[below=0.1cm of m, font=\small] {\textenglish{mod}: سجل};
        \node[below=0.1cm of r, font=\small] {\textenglish{reg} = \textenglish{AX}};
        \node[below=0.1cm of rm, font=\small] {\textenglish{r/m} = \textenglish{BX}};
    \end{tikzpicture}
    \end{center}
    إذن: انقل \textenglish{AX} (المصدر) إلى \textenglish{BX}.
\end{khatwa}

\begin{center}
\small
\begin{tabular}{|c|c|c|c|c|c|c|c|c|}
    \hline
    \ras{الرقم} & 000 & 001 & 010 & 011 & 100 & 101 & 110 & 111 \\ \hline
    \ras{\textenglish{w = 1}} & \textenglish{AX} & \textenglish{CX} & \textenglish{DX} & \textenglish{BX} & \textenglish{SP} & \textenglish{BP} & \textenglish{SI} & \textenglish{DI} \\ \hline
    \ras{\textenglish{w = 0}} & \textenglish{AL} & \textenglish{CL} & \textenglish{DL} & \textenglish{BL} & \textenglish{AH} & \textenglish{CH} & \textenglish{DH} & \textenglish{BH} \\ \hline
\end{tabular}
\end{center}

\begin{mithal}[: \textenglish{88 07} و\textenglish{EB 01}]
    \textenglish{88} = \bin{100010\,0\,0}: \code{mov} على 8 بتات. \textenglish{07} = \bin{00\,000\,111}:
    \textenglish{mod = 00} (ذاكرة دون إزاحة)، \textenglish{reg = 000} (\textenglish{AL})، \textenglish{r/m = 111} (\textenglish{[BX]}).
    إذن \code{mov [bx], al}.\\[3pt]
    \textenglish{EB 01}: قفز قصير، والإزاحة \textbf{نسبية}: بعد قراءة \code{jmp} يكون \textenglish{IP = \hex{010C}}،
    ونضيف 1 فنصل إلى \hex{010D} (\code{fin}). هكذا قفزنا فوق \code{nop}.
\end{mithal}

%================================================
\section{إدارة الإدخال والإخراج}
%================================================

\subsection{المنافذ: \code{in} و\code{out}}
يتصل المعالج بالأجهزة عبر \mustalah{منافذ}{Ports}{I/O ports} مرقمة من 0 إلى \hex{FFFF}.
\begin{itemize}
    \item \code{in al, 60h}: اقرأ بايتًا من المنفذ \hex{60} (لوحة المفاتيح في الحاسوب الشخصي).
    \item \code{out 61h, al}: أرسل \textenglish{AL} إلى المنفذ \hex{61} (مكبر الصوت).
\end{itemize}
البرامج العادية نادرًا ما تستعمل المنافذ مباشرة، بل تطلب ذلك من نظام التشغيل عبر المقاطعات.

\subsection{خدمات \textenglish{DOS}: المقاطعة \code{int 21h}}
\begin{center}
\small
\begin{tabularx}{\linewidth}{|c|R|R|}
    \hline
    \ras{\textenglish{AH}} & \ras{الوظيفة} & \ras{المدخلات والمخرجات} \\ \hline
    \hex{01} & قراءة حرف من لوحة المفاتيح مع إظهاره & المخرج: الحرف في \textenglish{AL} \\ \hline
    \hex{02} & كتابة حرف & المدخل: الحرف في \textenglish{DL} \\ \hline
    \hex{09} & كتابة سلسلة تنتهي بـ\code{\$} & المدخل: عنوانها في \textenglish{DX} \\ \hline
    \hex{4C} & إنهاء البرنامج & \textenglish{AL} = رمز النهاية \\ \hline
\end{tabularx}
\end{center}

\asm{ch3/io.asm}
\screen{ch3/io.out}

\begin{sharh}
    \s{7--8}{ننتظر أن يضغط المستعمل حرفًا (هنا \code{m})، فيصل رمزه \textenglish{ASCII} في \textenglish{AL}: 109.}
    \s{9}{في جدول \textenglish{ASCII}، الحرف الكبير أصغر من الصغير بـ32: $109 - 32 = 77$، وهو رمز \code{'M'}.}
    \s{10}{نحفظ الحرف في \textenglish{BL}، لأن الطباعة في السطر 13 ستغيّر \textenglish{AL}.}
    \s{14--16}{الوظيفة \hex{02} تكتب الحرف الموجود في \textenglish{DL}.}
    \s{20}{\code{13, 10}: الرجوع إلى بداية السطر ثم سطر جديد.}
\end{sharh}

%================================================
\section{إدارة المكدس}
%================================================

\begin{tarif}
    \mustalah{المكدس}{Pile}{Stack}: منطقة من الذاكرة نضع فيها القيم ونأخذها بقاعدة
    \textbf{آخر من يدخل أول من يخرج} \textenglish{(LIFO: Last In, First Out)}. قمته محددة بـ\textenglish{SS:SP}.
\end{tarif}

\begin{tashbih}
    المكدس مثل \textbf{كومة صحون}: نضع الصحن الجديد فوق الكومة، ونأخذ دائمًا الصحن الذي في الأعلى.
\end{tashbih}

\begin{itemize}
    \item \code{push ax}: \nql{SP}{SP - 2}، ثم نكتب \textenglish{AX} في \textenglish{[SP]}.
    \item \code{pop ax}: نقرأ \textenglish{[SP]} في \textenglish{AX}، ثم \nql{SP}{SP + 2}.
\end{itemize}

\begin{intabih}
    في \textenglish{8086} المكدس \textbf{ينمو نحو العناوين الصغيرة}: كل \code{push} ينقص \textenglish{SP}.
    ونضع ونأخذ دائمًا \textbf{كلمات} (16 بتًا).
\end{intabih}

\asm{ch3/stack.asm}

\begin{center}
\begin{tikzpicture}[c/.style={draw=brandDark, fill=white, minimum width=1.6cm, minimum height=0.55cm, font=\codefont\small},
    a/.style={font=\codefont\scriptsize, lawn=tealDeep}]
    \foreach \t/\x/\v/\w/\sp in {{البداية}/0/{}/{}/FFFE, {بعد \code{push ax}}/3.3/1111/{}/FFFC,
                                {بعد \code{push bx}}/6.6/1111/2222/FFFA, {بعد \code{pop ax}}/9.9/1111/{}/FFFC} {
        \node[font=\small\bfseries] at (\x,1.1) {\t};
        \node[c] (u) at (\x,0.3) {\textenglish{\v}};
        \node[c] (d) at (\x,-0.25) {\textenglish{\w}};
        \node[a, left=0.05cm of u] {\textenglish{FFFC}};
        \node[a, left=0.05cm of d] {\textenglish{FFFA}};
        \node[font=\codefont\small] at (\x,-1) {\textenglish{SP = \sp}};
    }
\end{tikzpicture}
\end{center}

\begin{center}
\small
\begin{tabular}{|c|l|c|c|c|}
    \hline
    \ras{السطر} & \ras{التعليمة} & \ras{\textenglish{AX}} & \ras{\textenglish{BX}} & \ras{\textenglish{SP}} \\ \hline
    \textenglish{4--5} & \code{mov ax/bx} & 1111 & 2222 & FFFE \\ \hline
    6 & \code{push ax} & 1111 & 2222 & \cell{orangeDeep}{FFFC} \\ \hline
    7 & \code{push bx} & 1111 & 2222 & \cell{orangeDeep}{FFFA} \\ \hline
    8 & \code{pop ax} & \cell{orangeDeep}{2222} & 2222 & FFFC \\ \hline
    9 & \code{pop bx} & 2222 & \cell{orangeDeep}{1111} & FFFE \\ \hline
\end{tabular}
\end{center}
بما أن آخر قيمة دخلت (\hex{2222}) هي أول قيمة خرجت، فقد تبادل \textenglish{AX} و\textenglish{BX} قيمتيهما.

\subsection{لماذا نحتاج إلى المكدس؟}
\begin{enumerate}
    \item \textbf{حفظ عناوين العودة}: عندما نستدعي إجراء، يحفظ المعالج في المكدس أين يجب أن يعود.
    \item \textbf{حفظ السجلات مؤقتًا}: إجراء يحتاج إلى \textenglish{CX} وهو مستعمل في حلقة؟ نكتب \code{push cx}
          في بدايته و\code{pop cx} في نهايته، فلا يضيع عدّاد الحلقة.
    \item \textbf{تمرير المعاملات} إلى الإجراءات (سنراه بعد قليل).
    \item \textbf{المقاطعات}: المعالج يحفظ فيه حالته قبل أن يذهب إلى معالج المقاطعة.
\end{enumerate}

\begin{intabih}
    يجب أن يكون عدد \code{pop} مساويًا لعدد \code{push}، وبالترتيب المعاكس. \code{push ax} ثم \code{push bx}
    يجب أن يقابلهما \code{pop bx} ثم \code{pop ax}. إذا نسيت \code{pop} واحدًا داخل إجراء، فإن \code{ret} ستأخذ
    قيمة خاطئة كعنوان عودة، فيقفز البرنامج إلى مكان عشوائي.
\end{intabih}

%================================================
\section{الدوال والإجراءات}
%================================================

\begin{tarif}
    \mustalah{الإجراء}{Procédure, sous-programme}{Procedure}: مجموعة تعليمات لها اسم، نستدعيها بـ\code{call}
    ونعود منها بـ\code{ret}. نكتبها مرة واحدة ونستعملها عدة مرات.
\end{tarif}

\subsection{كيف يعرف \code{ret} إلى أين يعود؟}
\begin{itemize}
    \item \code{call P}: يضع في المكدس \textbf{عنوان التعليمة التي بعده} (\textenglish{IP})، ثم يقفز إلى \textenglish{P}.
    \item \code{ret}: يأخذ هذا العنوان من المكدس ويضعه في \textenglish{IP}، فيعود التنفيذ بعد \code{call}.
\end{itemize}

\begin{khatwa}[: الاستدعاءان في البرنامج الآتي (العناوين الحقيقية)]
    \begin{enumerate}
        \item التعليمة \code{call print\_num} الأولى في الإزاحة \hex{0103} وطولها 3 بايتات، إذن التعليمة التي بعدها في \hex{0106}.
              يضع المعالج \hex{0106} في المكدس، ثم يقفز إلى \code{print\_num}.
        \item في نهاية الإجراء، \code{ret} تأخذ \hex{0106} من المكدس: يعود التنفيذ إلى \code{mov ax, 7}.
        \item الاستدعاء الثاني في \hex{0109}، فيضع \hex{010C}، و\code{ret} تعود هذه المرة إلى \hex{010C}.
    \end{enumerate}
    الإجراء نفسه يعود إلى مكانين مختلفين، لأن عنوان العودة محفوظ في المكدس في كل مرة.
\end{khatwa}

\subsection{إجراء يطبع عددًا بالعشري}
المقاطعة \code{int 21h} تطبع حروفًا فقط. لكي نطبع العدد 2026 يجب أن نستخرج أرقامه واحدًا واحدًا:

\begin{khatwa}[: استخراج أرقام 2026]
    \begin{center}
    \begin{tabular}{|c|c|c|c|}
        \hline
        \ras{العملية} & \ras{الحاصل (\textenglish{AX})} & \ras{الباقي (\textenglish{DX})} & \ras{المكدس} \\ \hline
        $2026 \div 10$ & 202 & 6 & 6 \\ \hline
        $202 \div 10$ & 20 & 2 & 6، 2 \\ \hline
        $20 \div 10$ & 2 & 0 & 6، 2، 0 \\ \hline
        $2 \div 10$ & 0 (نتوقف) & 2 & 6، 2، 0، 2 \\ \hline
    \end{tabular}
    \end{center}
    الأرقام تخرج بالعكس (6، 2، 0، 2). نضعها في المكدس، ثم نأخذها: فتخرج بالترتيب الصحيح 2، 0، 2، 6.
\end{khatwa}

\asm{ch3/procnum.asm}
\screen{ch3/procnum.out}

\begin{sharh}
    \s{4--5}{نضع العدد في \textenglish{AX} (المعامل)، ثم نستدعي الإجراء.}
    \s{10}{\code{proc} و\code{endp} يحددان بداية الإجراء ونهايته.}
    \s{14}{\textenglish{DX = 0} ضروري قبل \code{div bx} (القسمة 16 بتًا تقسم \textenglish{DX:AX}).}
    \s{16--17}{نضع الرقم في المكدس، ونعدّ الأرقام في \textenglish{CX}.}
    \s{21--25}{\textenglish{CX} مرة: نأخذ رقمًا، نحوله إلى حرف بإضافة \code{'0'} (48)، ونطبعه.}
    \s{28}{\code{ret} تعود إلى السطر 6 في المرة الأولى، وإلى السطر 8 في المرة الثانية.}
\end{sharh}

\subsection{تمرير المعاملات عبر المكدس}
\asm{ch3/procstack.asm}

بعد \code{mov bp, sp} في السطر 14، يبدو المكدس هكذا (القيم من التنفيذ الحقيقي):

\begin{center}
\begin{tikzpicture}[c/.style={draw=brandDark, fill=white, minimum width=3.2cm, minimum height=0.6cm, font=\small},
    a/.style={font=\codefont\small, lawn=tealDeep}]
    \foreach \ad/\v/\l [count=\i from 0] in {FFFC/{30 (المعامل الثاني)}/{BP+6}, FFFA/{12 (المعامل الأول)}/{BP+4},
                                           FFF8/{عنوان العودة \textenglish{010B}}/{BP+2}, FFF6/{\textenglish{BP} القديم}/{BP}} {
        \node[c] (n\i) at (0,-0.6*\i) {\v};
        \node[a, right=0.15cm of n\i] {\textenglish{\ad}};
        \node[font=\codefont\small\bfseries, left=0.15cm of n\i] {\textenglish{\l}};
    }
    \draw[sahmO] (-4.6,-1.8) node[left, font=\small] {\textenglish{SP = BP = FFF6}} -- (-3.0,-1.8);
\end{tikzpicture}
\end{center}

\begin{sharh}
    \s{4--7}{نضع المعاملين في المكدس: \textenglish{SP} ينتقل من \hex{FFFE} إلى \hex{FFFA}.}
    \s{8}{\code{call} يضع عنوان العودة \hex{010B}: \textenglish{SP = \hex{FFF8}}.}
    \s{13--14}{نحفظ \textenglish{BP}، ثم نجعل \textenglish{BP} ثابتًا على قمة المكدس (\hex{FFF6}).}
    \s{15--16}{المعامل الأول في \textenglish{[BP+4]}، والثاني في \textenglish{[BP+6]}: $12 + 30 = 42$.}
    \s{18}{\code{ret 4}: نعود، ثم نزيل 4 بايتات (المعاملين) من المكدس. \textenglish{SP} يعود إلى \hex{FFFE}.}
\end{sharh}

%================================================
\section{المقاطعات}
%================================================

\begin{tarif}
    \mustalah{المقاطعة}{Interruption}{Interrupt}: حدث يوقف البرنامج الحالي مؤقتًا، لينفذ المعالج
    برنامجًا خاصًا (\textbf{معالج المقاطعة})، ثم يعود ليكمل من حيث توقف.
\end{tarif}

\begin{tashbih}
    أنت تقرأ كتابًا، فيرن \textbf{جرس الباب}. تضع علامة في الصفحة، تفتح الباب، ثم تعود إلى القراءة
    من الصفحة نفسها. الجرس هو المقاطعة، والعلامة في الكتاب هي ما يحفظه المعالج في المكدس.
\end{tashbih}

\subsection{أنواع المقاطعات}
\begin{center}
\small
\begin{tabularx}{\linewidth}{|>{\centering\arraybackslash}p{3cm}|R|R|}
    \hline
    \ras{النوع} & \ras{من يسببها؟} & \ras{أمثلة} \\ \hline
    عتادية & جهاز خارجي يرسل إشارة على المنفذ \textenglish{INTR} أو \textenglish{NMI} & ضغط مفتاح، الساعة، الطابعة جاهزة \\ \hline
    برمجية & التعليمة \code{int n} في البرنامج & \code{int 21h} لخدمات \textenglish{DOS} \\ \hline
    استثناءات & المعالج نفسه عند خطأ & القسمة على صفر (مقاطعة 0) \\ \hline
\end{tabularx}
\end{center}

\subsection{جدول أشعة المقاطعات}
في بداية الذاكرة (من \hex{00000} إلى \hex{003FF}) يوجد جدول من 256 خانة، لكل مقاطعة 4 بايتات:
الإزاحة ثم المقطع لمعالج المقاطعة.

\begin{tarif}
    \centering
    شعاع المقاطعة رقم $n$ يوجد في العنوان: \quad $4 \times n$
\end{tarif}

\begin{mithal}
    شعاع \code{int 21h}: $\hex{21} \times 4 = 33 \times 4 = 132 = \hex{84}$. إذن عنوان معالج \textenglish{DOS}
    محفوظ في \hex{00084} إلى \hex{00087}.
\end{mithal}

\subsection{ماذا يفعل المعالج عند \code{int n}؟}
\begin{enumerate}
    \item يضع \textenglish{FLAGS} في المكدس، ثم يطفئ \textenglish{IF} و\textenglish{TF}.
    \item يضع \textenglish{CS} ثم \textenglish{IP} في المكدس (العلامة في الكتاب).
    \item يقرأ \textenglish{IP} و\textenglish{CS} الجديدين من الخانة $n \times 4$، فيقفز إلى معالج المقاطعة.
    \item في نهاية المعالج، \code{iret} تأخذ \textenglish{IP} و\textenglish{CS} و\textenglish{FLAGS} من المكدس، فيعود البرنامج.
\end{enumerate}

\subsection{المقاطعة العتادية خطوة بخطوة}
\begin{khatwa}[: تضغط على مفتاح أثناء تشغيل برنامج]
    \begin{enumerate}
        \item لوحة المفاتيح ترسل إشارة إلى \mustalah{متحكم المقاطعات}{Contrôleur d'interruptions}{PIC 8259}.
        \item المتحكم يرسل إشارة على المدخل \textenglish{INTR} للمعالج.
        \item إذا كان \textenglish{IF = 1}، يُكمل المعالج \textbf{التعليمة الحالية} أولًا، ثم يقبل المقاطعة.
        \item المتحكم يعطي المعالج رقم المقاطعة (\hex{09} للوحة المفاتيح).
        \item المعالج يحفظ \textenglish{FLAGS} و\textenglish{CS} و\textenglish{IP}، ثم يقفز إلى العنوان المكتوب في $4 \times 9 = \hex{24}$.
        \item معالج المقاطعة يقرأ المفتاح من المنفذ \hex{60}، يحفظه، ثم ينفذ \code{iret}.
        \item البرنامج يكمل من حيث توقف، دون أن «يشعر» بشيء.
    \end{enumerate}
\end{khatwa}

\begin{tarif}
    \begin{itemize}
        \item \code{cli} \textenglish{(clear IF)}: تمنع المقاطعات العتادية مؤقتًا، و\code{sti} \textenglish{(set IF)} تعيد السماح بها.
              نستعملهما عندما يجب ألّا يُقاطَع جزء حساس من البرنامج.
        \item المدخل \textenglish{NMI} \textenglish{(Non Maskable Interrupt)}: مقاطعة \textbf{لا يمكن منعها}، تُستعمل للأخطاء الخطيرة
              مثل عطل في الذاكرة أو انقطاع التيار.
    \end{itemize}
\end{tarif}

\subsection{برنامج: مقاطعتنا الخاصة}
\asm{ch3/myint.asm}
\screen{ch3/myint.out}

\begin{sharh}
    \s{4--5}{\textenglish{ES = 0}: نريد الكتابة في جدول الأشعة الموجود في المقطع 0.}
    \s{6--8}{نكتب عنوان \code{handler} في الخانة $\hex{60} \times 4 = \hex{180}$: الإزاحة ثم المقطع.}
    \s{10--11}{كل \code{int 60h} ينفذ \code{handler}: المرة الأولى يطبع \code{A}، والثانية \code{B}.}
    \s{19}{\code{iret} وليس \code{ret}، لأن المقاطعة وضعت في المكدس \textenglish{FLAGS} أيضًا.}
\end{sharh}

%================================================
\section{تمارين}
%================================================

\begin{tamrin}
    احسب العنوان الفيزيائي في كل حالة: (أ) \textenglish{DS = 3000H}، \textenglish{BX = 0050H} مع \code{mov al, [bx]}.
    (ب) \textenglish{SS = 4000H}، \textenglish{SP = FFFEH}.
\end{tamrin}
\begin{hall}
    (أ) $\hex{30000} + \hex{0050} = \hex{30050}$. \quad (ب) $\hex{40000} + \hex{FFFE} = \hex{4FFFE}$.
\end{hall}

\begin{tamrin}
    ما قيمة \textenglish{AX} بعد هذه التعليمات؟ \code{mov ax, 0F0Fh} ثم \code{and ax, 00FFh} ثم \code{shl ax, 1}.
\end{tamrin}
\begin{hall}
    بعد \code{and}: \hex{000F}. بعد الإزاحة (ضرب في 2): \textenglish{\hex{001E} = 30}.
\end{hall}

\begin{tamrin}
    اكتب برنامجًا يحسب $5! = 1 \times 2 \times 3 \times 4 \times 5$ بالتعليمة \code{loop}، ويضع النتيجة في \textenglish{AX}.
\end{tamrin}
\begin{hall}
    \asm{ch3/ex_fact.asm}
    النتيجة: \textenglish{AX = 0078} = 120 (جربناه في المحاكي).
\end{hall}

\begin{tamrin}
    بعد \code{mov sp, 0100h} ثم \code{push ax} ثم \code{push bx} ثم \code{pop cx}: ما قيمة \textenglish{SP}، وأي قيمة في \textenglish{CX}؟
\end{tamrin}
\begin{hall}
    بعد الدفعتين \textenglish{SP = \hex{00FC}}، وبعد \code{pop} \textenglish{SP = \hex{00FE}}.
    \textenglish{CX} يأخذ آخر قيمة دخلت: قيمة \textenglish{BX}.
\end{hall}

\begin{tamrin}
    ما عنوان شعاع المقاطعة رقم \hex{10} (خدمات الشاشة في \textenglish{BIOS})؟
\end{tamrin}
\begin{hall}
    $\hex{10} \times 4 = 16 \times 4 = 64 = \hex{40}$: الخانات من \hex{00040} إلى \hex{00043}.
\end{hall}

\begin{tamrin}
    فك ترميز التعليمة \textenglish{89 D8}.
\end{tamrin}
\begin{hall}
    \textenglish{89}: \code{mov} 16 بتًا، \textenglish{d = 0}. \textenglish{D8} = \bin{11\,011\,000}: \textenglish{mod = 11} (سجل)،
    \textenglish{reg = 011} (\textenglish{BX}، المصدر)، \textenglish{r/m = 000} (\textenglish{AX}). إذن: \code{mov ax, bx}.
\end{hall}

\chapter{دراسة حالة: المعالجات ذات 32 و64 بتًا}

\begin{khulasa}[: ماذا سنتعلم في هذا الفصل؟]
\begin{itemize}
    \item لماذا انتقلنا من 16 بتًا إلى 32 ثم 64 بتًا، وما الذي تغيّر.
    \item تطور المعمارية من \textenglish{IA-32} إلى \textenglish{AMD64}.
    \item سجلات المعالجات الحديثة، والعنونة المسطحة.
    \item تعليمات جديدة: \textenglish{SSE} و\textenglish{AVX} التي تحسب عدة أعداد في وقت واحد.
    \item خط الأنابيب: كيف ينفذ المعالج عدة تعليمات في الوقت نفسه.
\end{itemize}
\end{khulasa}

\begin{intabih}
    برامج هذا الفصل مكتوبة بصيغة \textbf{\textenglish{NASM}} لنظام \textbf{\textenglish{Linux}}، وقد نفذناها فعلًا على معالج
    \textenglish{Intel Core i7-8665U} (4 أنوية). للتجميع والتنفيذ:
    \code{nasm -f elf64 prog.asm} ثم \code{ld prog.o -o prog} ثم \code{./prog}.
\end{intabih}

%================================================
\section{من 16 إلى 32 ثم 64 بتًا}
%================================================

\begin{tarif}
    نقول إن المعالج من \textbf{$n$ بتًا} عندما تكون \textbf{سجلاته العامة} من $n$ بتًا: هذا هو حجم العدد الذي
    يحسبه في تعليمة واحدة، وحجم العنوان الذي يستطيع وضعه في سجل.
\end{tarif}

\subsection{الفروق الأساسية}
\begin{center}
\small
\begin{tabularx}{\linewidth}{|>{\centering\arraybackslash}p{3.2cm}|C|C|C|}
    \hline
    \ras{} & \ras{16 بتًا (\textenglish{8086})} & \ras{32 بتًا (\textenglish{IA-32})} & \ras{64 بتًا (\textenglish{x86-64})} \\ \hline
    أكبر عدد دون إشارة & $2^{16}-1 = 65\,535$ & $2^{32}-1 \approx 4.3$ مليار & $2^{64}-1 \approx 1.8 \times 10^{19}$ \\ \hline
    الذاكرة القصوى & \textenglish{1 MB} (بالمقاطع) & \textenglish{4 GB} & \textenglish{256 TB} (48 بتًا عمليًا) \\ \hline
    عدد السجلات العامة & 8 & 8 & \textbf{16} \\ \hline
    أسماء السجلات & \textenglish{AX, BX, \dots} & \textenglish{EAX, EBX, \dots} & \textenglish{RAX, RBX, \dots, R15} \\ \hline
    نمط العنونة & مقطعي & مسطح (عمليًا) & مسطح \\ \hline
\end{tabularx}
\end{center}

\subsection{لماذا نحتاج إلى أكثر من 32 بتًا؟}
\begin{enumerate}
    \item \textbf{الذاكرة}: بعنوان من 32 بتًا لا نستطيع استعمال أكثر من \textenglish{4 GB}. حواسيب اليوم فيها
          \textenglish{16} أو \textenglish{32 GB}، والخوادم فيها مئات الغيغابايتات.
    \item \textbf{الأداء}: سجل من 64 بتًا يحسب عددًا كبيرًا في تعليمة واحدة، بينما يحتاج معالج 32 بتًا إلى
          تعليمتين على الأقل (كما رأينا \code{add} و\code{adc} في الفصل الثالث). ومضاعفة عدد السجلات (16 بدل 8)
          تقلل الذهاب إلى الذاكرة.
    \item \textbf{التوافق}: المعالج من 64 بتًا ينفذ أيضًا البرامج القديمة ذات 32 بتًا دون تغيير.
\end{enumerate}

\begin{khatwa}[: ما معنى «\textenglish{4 GB}» بالضبط؟]
    \begin{enumerate}
        \item عنوان من 32 بتًا يأخذ $2^{32}$ قيمة مختلفة.
        \item $2^{32} = 2^{2} \times 2^{30} = 4 \times 2^{30}$.
        \item و$2^{30}$ بايت $= 1$ \textenglish{GB}، إذن السعة القصوى $= 4$ \textenglish{GB}.
    \end{enumerate}
    بـ48 بتًا: $2^{48} = 2^{8} \times 2^{40} = 256$ \textenglish{TB}، أي أكبر بـ65536 مرة!
\end{khatwa}

%================================================
\section{تطور المعمارية: من \textenglish{IA-32} إلى \textenglish{AMD64}}
%================================================

\begin{center}
\begin{tikzpicture}[e/.style={unit=#1, minimum width=2.9cm, minimum height=1.3cm, font=\small}]
    \node[e=brandMuted] (a) at (0,0) {\textenglish{8086} (1978)\\16 بتًا، مقاطع};
    \node[e=brandPurple] (b) at (-3.6,0) {\textenglish{80286} (1982)\\النمط المحمي};
    \node[e=brandTeal] (c) at (-7.2,0) {\textenglish{80386} (1985)\\\textbf{\textenglish{IA-32}}};
    \node[e=brandOrange] (d) at (-10.8,0) {\textenglish{AMD64} (2003)\\\textbf{64 بتًا}};
    \draw[sahmT] (a) -- (b); \draw[sahmT] (b) -- (c); \draw[sahmT] (c) -- (d);
\end{tikzpicture}
\end{center}

\begin{itemize}
    \item \textbf{\textenglish{IA-32}} \textenglish{(Intel Architecture 32 bits)}: بدأت مع \textenglish{80386}. السجلات أصبحت من 32 بتًا
          (\textenglish{EAX}، الحرف \textenglish{E} يعني \textenglish{Extended})، وظهر \textbf{النمط المحمي} الذي يمنع البرامج من
          الكتابة في ذاكرة بعضها البعض.
    \item \textbf{\textenglish{AMD64}}: المفاجأة أن شركة \textenglish{AMD} (وليس \textenglish{Intel}) هي التي صممت الامتداد إلى 64 بتًا سنة 2003،
          مع المحافظة على التوافق. تبنّته \textenglish{Intel} بعد ذلك باسم \textenglish{Intel 64}. نسميه اليوم \textenglish{x86-64}.
\end{itemize}

\subsection{أنماط عمل المعالج}
\begin{center}
\small
\begin{tabularx}{\linewidth}{|>{\centering\arraybackslash}p{3.3cm}|R|R|}
    \hline
    \ras{النمط} & \ras{الوصف} & \ras{متى؟} \\ \hline
    الحقيقي \textenglish{(real mode)} & يتصرف كمعالج \textenglish{8086}: 16 بتًا و\textenglish{1 MB} & عند تشغيل الحاسوب، قبل نظام التشغيل \\ \hline
    المحمي \textenglish{(protected mode)} & 32 بتًا، حماية الذاكرة & أنظمة 32 بتًا \\ \hline
    الطويل \textenglish{(long mode)} & 64 بتًا، ومعه نمط فرعي للتوافق مع برامج 32 بتًا & كل الأنظمة الحديثة \\ \hline
\end{tabularx}
\end{center}

\begin{hal}
    كل حاسوب حديث يستيقظ وهو «يعتقد» أنه \textenglish{8086}! يبدأ في النمط الحقيقي، ثم ينقله نظام التشغيل إلى
    النمط المحمي ثم الطويل. هكذا تبقى المعمارية متوافقة منذ 1978.
\end{hal}

%================================================
\section{السجلات في معالجات 32 و64 بتًا}
%================================================

\subsection{السجلات المتداخلة}
كل سجل جديد \textbf{يحتوي} السجل القديم في نصفه الأيمن (البتات الصغيرة):

\begin{center}
\begin{tikzpicture}[b/.style={draw=brandDark, thick, minimum height=0.75cm, font=\codefont\bfseries}]
    \node[b, minimum width=12cm, fill=brandOrange!15, anchor=east] (r) at (6,1.6) {\textenglish{RAX: 64 bits}};
    \node[b, minimum width=6cm, fill=brandTeal!15, anchor=east] (e) at (6,0.8) {\textenglish{EAX: 32 bits}};
    \node[b, minimum width=3cm, fill=brandGreen!20, anchor=east] (x) at (6,0) {\textenglish{AX: 16}};
    \node[b, minimum width=1.5cm, fill=brandPurple!20, anchor=east] (h) at (4.5,-0.8) {\textenglish{AH}};
    \node[b, minimum width=1.5cm, fill=brandPurple!35, anchor=east] (l) at (6,-0.8) {\textenglish{AL}};
    \foreach \x/\t in {6/0, 4.5/8, 3/16, 0/32, -6/64} \node[font=\codefont\scriptsize] at (\x,-1.45) {\textenglish{\t}};
    \draw[brandMuted, dashed] (0,-1.25) -- (0,1.2); \draw[brandMuted, dashed] (3,-1.25) -- (3,0.4);
\end{tikzpicture}
\end{center}

\begin{center}
\small
\begin{tabular}{|c|c|c|c|l|}
    \hline
    \ras{64 بتًا} & \ras{32 بتًا} & \ras{16 بتًا} & \ras{8 بتات} & \ras{الاستعمال المعتاد} \\ \hline
    \textenglish{RAX} & \textenglish{EAX} & \textenglish{AX} & \textenglish{AL} & النتيجة المعادة من دالة \\ \hline
    \textenglish{RBX} & \textenglish{EBX} & \textenglish{BX} & \textenglish{BL} & عام \\ \hline
    \textenglish{RCX} & \textenglish{ECX} & \textenglish{CX} & \textenglish{CL} & العدّاد، المعامل الرابع \\ \hline
    \textenglish{RDX} & \textenglish{EDX} & \textenglish{DX} & \textenglish{DL} & المعامل الثالث \\ \hline
    \textenglish{RSI} & \textenglish{ESI} & \textenglish{SI} & \textenglish{SIL} & المعامل الثاني \\ \hline
    \textenglish{RDI} & \textenglish{EDI} & \textenglish{DI} & \textenglish{DIL} & المعامل الأول \\ \hline
    \textenglish{RSP} & \textenglish{ESP} & \textenglish{SP} & \textenglish{SPL} & قمة المكدس \\ \hline
    \textenglish{RBP} & \textenglish{EBP} & \textenglish{BP} & \textenglish{BPL} & قاعدة إطار المكدس \\ \hline
    \textenglish{R8 \dots R15} & \textenglish{R8D \dots R15D} & \textenglish{R8W \dots} & \textenglish{R8B \dots} & \textbf{جديدة في 64 بتًا} \\ \hline
\end{tabular}
\end{center}

\begin{itemize}
    \item \textbf{\textenglish{RIP}}: مؤشر التعليمة (كان \textenglish{IP} ثم \textenglish{EIP}).
    \item \textbf{\textenglish{RFLAGS}}: سجل الأعلام، والأعلام \textenglish{ZF}، \textenglish{CF}، \textenglish{SF}، \textenglish{OF} في المكان نفسه كما في \textenglish{8086}.
    \item \textbf{سجلات المقاطع}: ما زالت موجودة (\textenglish{CS}، \textenglish{DS}\dots) لكن بدايتها كلها 0 في 64 بتًا، ما عدا
          \textenglish{FS} و\textenglish{GS} التي يستعملها النظام لمعطيات خاصة بكل خيط تنفيذ.
\end{itemize}

\subsection{تجربة: ماذا يحدث للبتات الأخرى؟}
\asm{ch4/regs64.asm}
\screen{ch4/regs64.out}

\begin{sharh}
    \s{5--6}{نضع عددًا من 64 بتًا في \textenglish{RAX} ونعرضه: \textenglish{1122334455667788}.}
    \s{9}{الكتابة في \textenglish{EAX} (32 بتًا) \textbf{تصفّر} النصف الأعلى: \textenglish{00000000AABBCCDD}.}
    \s{13}{الكتابة في \textenglish{AX} (16 بتًا) تترك باقي البتات كما هي: \textenglish{\dots 5566}\textbf{\textenglish{1234}}.}
    \s{17}{الكتابة في \textenglish{AL} تغيّر البايت الأخير فقط: \textenglish{\dots 77}\textbf{\textenglish{FF}}.}
    \s{24--42}{إجراء يعرض \textenglish{RAX} بالست عشري: في كل دورة يدوّر 4 بتات، يأخذها بالقناع \code{0xF}، ويحوّلها إلى حرف.}
\end{sharh}

\begin{intabih}
    هذه قاعدة غريبة لكنها مهمة: في 64 بتًا، \textbf{كل كتابة في سجل من 32 بتًا تصفّر البتات العليا}، أما الكتابة
    في 8 أو 16 بتًا فلا تفعل ذلك. صممت \textenglish{AMD} هذه القاعدة لتسريع المعالج.
\end{intabih}

%================================================
\section{أنماط العنونة المتقدمة}
%================================================

\subsection{من العنونة المقطعية إلى العنونة المسطحة}
\begin{center}
\begin{tikzpicture}
    \draw[thick, fill=brandLight] (0,0) rectangle (2,-4);
    \foreach \y/\c in {0/brandTeal, -1.3/brandOrange, -2.6/brandGreen} \fill[\c!30] (0,\y) rectangle (2,\y-1.2);
    \foreach \y/\t in {-0.6/\textenglish{CS}, -1.9/\textenglish{DS}, -3.2/\textenglish{SS}} \node[font=\codefont] at (1,\y) {\t};
    \node[font=\small\bfseries, align=center] at (1,0.6) {مقطعي \textenglish{(8086)}};
    \node[font=\scriptsize, text width=4cm, align=center] at (1,-4.6) {مناطق من \textenglish{64 KB}، والعنوان = مقطع:إزاحة};
    \draw[thick, fill=brandTeal!15] (-6,0) rectangle (-4,-4);
    \node[font=\small, align=center] at (-5,-2) {مساحة واحدة\\متصلة};
    \node[font=\small\bfseries, align=center] at (-5,0.6) {مسطح \textenglish{(x86-64)}};
    \node[font=\scriptsize, text width=4cm, align=center] at (-5,-4.6) {العنوان = عدد واحد من 0 إلى النهاية};
    \draw[sahmO] (-0.4,-2) -- (-3.6,-2);
\end{tikzpicture}
\end{center}

في \mustalah{النموذج المسطح}{Modèle mémoire plat}{Flat memory model} يرى البرنامج الذاكرة \textbf{كشريط واحد طويل} من
البايتات، مرقمة من 0. لا نحتاج إلى حساب «مقطع × 16 + إزاحة»: العنوان هو مجرد عدد في سجل من 64 بتًا.

\begin{tashbih}
    العنونة المقطعية تشبه أن تقول: «العمارة 12، الطابق 3». العنونة المسطحة تشبه أن تعطي لكل شقة في المدينة
    رقمًا واحدًا فريدًا: «الشقة رقم 73\,502». أبسط بكثير، بشرط أن نملك أعدادًا كبيرة بما يكفي.
\end{tashbih}

\subsection{الشكل العام للعنوان}
\begin{tarif}
    \centering
    $\text{العنوان} = \text{القاعدة} + \text{الفهرس} \times \text{السلّم} + \text{الإزاحة}$
    \par\small
    القاعدة والفهرس: أي سجلين. السلّم: 1 أو 2 أو 4 أو 8. الإزاحة: عدد ثابت.
\end{tarif}

\begin{center}
\small
\begin{tabular}{|l|l|}
    \hline
    \ras{الكتابة (\textenglish{NASM})} & \ras{الاستعمال} \\ \hline
    \code{[rbx]} & القاعدة فقط: مؤشر \\ \hline
    \code{[rbx + 16]} & قاعدة + إزاحة: حقل في هيكل \textenglish{(struct)} \\ \hline
    \code{[tab + rcx*8]} & فهرس × سلّم + إزاحة: العنصر \textenglish{i} من جدول أعداد من 8 بايتات \\ \hline
    \code{[rbx + rcx*4 + 8]} & كل العناصر معًا \\ \hline
    \code{[rel msg]} & نسبي إلى \textenglish{RIP}: خاص بـ64 بتًا \\ \hline
\end{tabular}
\end{center}

\begin{khatwa}[: أين يقرأ \code{[tab + rcx*8]} عندما \textenglish{RCX = 3}؟]
    \begin{enumerate}
        \item كل عنصر من \code{tab} عدد من 64 بتًا = 8 بايتات (\code{dq}).
        \item العنصر رقم 3 (الرابع، لأننا نبدأ من 0) على بعد $3 \times 8 = 24$ بايتًا من البداية.
        \item إذا كان \code{tab} في العنوان \hex{402000}، فالعنوان هو $\hex{402000} + 24 = \hex{402018}$.
    \end{enumerate}
    السلّم يُغنينا عن ضرب الفهرس بأنفسنا: المعالج يفعل ذلك داخل التعليمة نفسها.
\end{khatwa}

\asm{ch4/sum64.asm}
\screen{ch4/sum64.out}

\begin{sharh}
    \s{5--6}{\code{xor} سجل مع نفسه = 0: المجموع والفهرس يبدآن من 0.}
    \s{8}{العنونة \textbf{بالسلّم}: نقرأ العنصر \textenglish{i} من الجدول، ونجمعه في \textenglish{RAX}.}
    \s{9--11}{\code{i++} ثم «طالما \textenglish{i < 5}» نعود. هذه حلقة \code{while} كما رأينا في الفصل الثالث.}
    \s{18--36}{نفس طريقة الفصل الثالث لعرض عدد بالعشري (القسمة على 10)، لكن بسجلات 64 بتًا.}
    \s{39}{\code{dq} \textenglish{(define quad-word)}: أعداد من 64 بتًا. العدد الأول ($10^{12}$) لا يدخل في 32 بتًا.}
\end{sharh}

\subsection{العنونة النسبية إلى \textenglish{RIP}}
في 64 بتًا نكتب كثيرًا \code{lea rsi, [rel buf]}: العنوان يُحسب \textbf{بالنسبة إلى التعليمة الحالية}
($\text{RIP} + \text{مسافة}$). الفائدة: يعمل البرنامج في أي مكان يضعه فيه النظام في الذاكرة، وهذا مهم للأمان
(النظام يغيّر مكان البرامج في كل تشغيل لكي يصعب على المهاجمين معرفة العناوين).

%================================================
\section{مجموعة التعليمات}
%================================================

\subsection{التعليمة نفسها على 16 و32 و64 بتًا}
جمّعنا البرنامج الآتي، وهذه بايتات لغة الآلة الحقيقية:

\asm{ch4/same.asm}

\begin{center}
\small
\begin{tabular}{|l|l|l|}
    \hline
    \ras{التعليمة} & \ras{لغة الآلة} & \ras{الملاحظة} \\ \hline
    \code{add ax, bx} & \textenglish{\codefont \textcolor{orangeDeep}{66} 01 D8} & البادئة \textenglish{66}: «انتبه، 16 بتًا» \\ \hline
    \code{add eax, ebx} & \textenglish{\codefont 01 D8} & 32 بتًا هو الحجم الافتراضي \\ \hline
    \code{add rax, rbx} & \textenglish{\codefont \textcolor{orangeDeep}{48} 01 D8} & البادئة \textenglish{REX.W = 48}: «64 بتًا» \\ \hline
    \code{add r8, r9} & \textenglish{\codefont \textcolor{orangeDeep}{4D} 01 C8} & \textenglish{REX = 4D}: 64 بتًا + سجلات جديدة \\ \hline
\end{tabular}
\end{center}

\begin{khatwa}[: لماذا \textenglish{01 D8} موجودة في الأسطر الثلاثة الأولى؟]
    \begin{enumerate}
        \item \textenglish{01} هو الرمز العملي «جمع سجل في سجل أو ذاكرة»، وهو نفسه منذ \textenglish{8086}.
        \item \textenglish{D8} = \bin{11\,011\,000} هو بايت \textenglish{ModR/M} الذي رأيناه في الفصل الثالث: \textenglish{mod = 11} (سجل)،
              \textenglish{reg = 011} (\textenglish{BX})، \textenglish{r/m = 000} (\textenglish{AX}).
        \item الحجم (16 أو 32 أو 64) لا يُكتب في التعليمة نفسها، بل في \textbf{بادئة} قبلها. هكذا بقيت التعليمات القديمة
              صالحة، وأضيفت الأحجام الجديدة «من الخارج».
        \item في \textenglish{4D} = \bin{0100\,1101}: \textenglish{0100} تعني «بادئة \textenglish{REX}»، ثم بت \textenglish{W = 1} (64 بتًا)، وبتّان
              يضيفان 8 إلى رقم السجلين، فيصبح \textenglish{AX} هو \textenglish{R8} و\textenglish{BX} هو \textenglish{R9}.
    \end{enumerate}
\end{khatwa}

\subsection{مقارنة: برنامج 32 بتًا وبرنامج 64 بتًا}
\textbf{بـ32 بتًا:}
\asm{ch4/hello32.asm}
\screen{ch4/hello32.out}
\textbf{بـ64 بتًا:}
\asm{ch4/hello64.asm}
\screen{ch4/hello64.out}

\begin{center}
\small
\begin{tabularx}{\linewidth}{|>{\centering\arraybackslash}p{3.6cm}|C|C|}
    \hline
    \ras{} & \ras{32 بتًا} & \ras{64 بتًا} \\ \hline
    طلب خدمة من النظام & \code{int 0x80} (مقاطعة) & \code{syscall} (تعليمة خاصة، أسرع) \\ \hline
    رقم الخدمة \textenglish{write} & 4 في \textenglish{EAX} & 1 في \textenglish{RAX} \\ \hline
    المعاملات & \textenglish{EBX}، \textenglish{ECX}، \textenglish{EDX} & \textenglish{RDI}، \textenglish{RSI}، \textenglish{RDX} \\ \hline
    حجم العنوان & 4 بايتات & 8 بايتات \\ \hline
\end{tabularx}
\end{center}

\begin{hal}
    نفذنا البرنامجين على الحاسوب نفسه ذي 64 بتًا. البرنامج ذو 32 بتًا عمل دون أي تعديل: هذا هو النمط الفرعي
    للتوافق داخل النمط الطويل.
\end{hal}

\subsection{ماذا يكتب المترجم؟ من \textenglish{C} إلى لغة التجميع}
كتبنا دالة بلغة \textenglish{C}، وطلبنا من المترجم \textenglish{gcc} أن يعطينا لغة التجميع التي يولّدها
(\code{gcc -O1 -S -masm=intel sum.c}):

\begin{center}
\begin{minipage}[t]{0.47\linewidth}
    \cfile{ch4/sum.c}
\end{minipage}\hfill
\begin{minipage}[t]{0.49\linewidth}
    \asm{ch4/sum_gcc.s}
\end{minipage}
\end{center}

\begin{sharh}
    \s{1}{حسب قواعد \textenglish{Linux} لـ64 بتًا: المعامل الأول (\code{t}) يصل في \textenglish{RDI}، والثاني (\code{n}) في \textenglish{RSI}.}
    \s{2--3}{إذا كان $n \leq 0$ نقفز إلى \code{.L4}: النتيجة 0.}
    \s{5}{\code{lea}: يحسب عنوان نهاية الجدول $t + n \times 8$ \textbf{دون أن يقرأ الذاكرة}.}
    \s{8--11}{الحلقة: نجمع العنصر، نتقدم بـ8 بايتات، ونتوقف عند النهاية. المترجم استعمل مؤشرًا بدل الفهرس \code{i}.}
    \s{13--14}{النتيجة تُعاد دائمًا في \textenglish{RAX}.}
\end{sharh}

\subsection{\textenglish{SSE} و\textenglish{AVX}: حساب عدة أعداد مرة واحدة}
\begin{tarif}
    \mustalah{\textenglish{SIMD}}{Une instruction, données multiples}{Single Instruction, Multiple Data}: تعليمة واحدة تطبق
    العملية نفسها على \textbf{عدة أعداد في الوقت نفسه}.
\end{tarif}

\begin{center}
\begin{tikzpicture}[c/.style={draw=brandDark, thick, minimum width=1.3cm, minimum height=0.65cm, font=\codefont}]
    \foreach \v [count=\i from 0] in {1.5, 2.5, 3.5, 4.5} \node[c, fill=brandTeal!15] at (-1.3*\i,1.6) {\textenglish{\v}};
    \foreach \v [count=\i from 0] in {10, 20, 30, 40} \node[c, fill=brandOrange!15] at (-1.3*\i,0.6) {\textenglish{\v}};
    \foreach \v [count=\i from 0] in {11.5, 22.5, 33.5, 44.5} \node[c, fill=brandGreen!20] at (-1.3*\i,-0.9) {\textenglish{\v}};
    \draw[very thick] (-4.6,0.1) -- (0.7,0.1);
    \node[font=\codefont\bfseries] at (1.4,0.6) {\textenglish{+}};
    \node[font=\small, anchor=west] at (1.2,1.6) {\textenglish{XMM0} (128 بتًا = 4 أعداد)};
    \node[font=\small, anchor=west] at (1.2,-0.9) {النتيجة: \textbf{تعليمة واحدة} \code{addps}};
\end{tikzpicture}
\end{center}

\begin{center}
\small
\begin{tabular}{|c|c|c|c|c|}
    \hline
    \ras{الامتداد} & \ras{السنة} & \ras{السجلات} & \ras{الحجم} & \ras{أعداد \textenglish{float} في المرة} \\ \hline
    \textenglish{SSE} & 1999 & \textenglish{XMM0--XMM15} & 128 بتًا & 4 \\ \hline
    \textenglish{AVX} & 2011 & \textenglish{YMM0--YMM15} & 256 بتًا & 8 \\ \hline
    \textenglish{AVX-512} & 2016 & \textenglish{ZMM0--ZMM31} & 512 بتًا & 16 \\ \hline
\end{tabular}
\end{center}

\cfile{ch4/simd.c}
\screen{ch4/simd.out}

عندما نفتح البرنامج المترجم، نجد فعلًا التعليمة الحقيقية:
\begin{codebox}\begin{english}\codefont\small
    addps \ xmm0, XMMWORD PTR [rbp-0x50] \quad \textcolor{red!75!black}{; 4 additions at once}
\end{english}\end{codebox}

ونسخة \textenglish{AVX} تجمع 8 أعداد بالتعليمة \code{vaddps ymm0, ymm0, [\dots]}:
\cfile{ch4/avx.c}
\screen{ch4/avx.out}

\begin{tashbih}
    التعليمة العادية تشبه أستاذًا يصحح ورقة امتحان واحدة في كل مرة. \textenglish{SIMD} تشبه أستاذًا يضع 8 أوراق جنبًا
    إلى جنب، ويصحح السؤال الأول فيها كلها بحركة واحدة. لهذا تُستعمل في معالجة الصور، الفيديو، الألعاب، والذكاء الاصطناعي،
    حيث نطبق العملية نفسها على ملايين الأعداد.
\end{tashbih}

%================================================
\section{دورة تنفيذ التعليمة وخط الأنابيب}
%================================================

\subsection{المراحل الخمس}
في المعالجات الحديثة، تنقسم كل تعليمة إلى مراحل، ولكل مرحلة دارة خاصة بها:
\begin{center}
\small
\begin{tabular}{|c|c|l|}
    \hline
    \ras{المرحلة} & \ras{الاسم} & \ras{العمل} \\ \hline
    \textenglish{IF} & الجلب & قراءة التعليمة من الذاكرة (بعنوان \textenglish{RIP}) \\ \hline
    \textenglish{ID} & فك الترميز & فهم التعليمة وقراءة السجلات \\ \hline
    \textenglish{EX} & التنفيذ & حساب النتيجة في \textenglish{UAL} أو حساب العنوان \\ \hline
    \textenglish{MEM} & الذاكرة & القراءة أو الكتابة في الذاكرة إن احتجنا \\ \hline
    \textenglish{WB} & الكتابة & حفظ النتيجة في السجل \\ \hline
\end{tabular}
\end{center}

\subsection{دون خط أنابيب ومع خط أنابيب}
\begin{tashbih}
    في مغسلة ملابس: الغسل ساعة، التجفيف ساعة، الكي ساعة. لغسل 3 أكياس واحدًا بعد الآخر نحتاج 9 ساعات.
    لكن عندما يخرج الكيس الأول من الغسالة إلى المجفف، نضع الكيس الثاني في الغسالة فورًا. الآلات الثلاث تعمل
    معًا، فننتهي في 5 ساعات فقط. هذا هو خط الأنابيب.
\end{tashbih}

\begin{center}
\begin{tikzpicture}[s/.style={draw=brandDark, minimum width=0.95cm, minimum height=0.55cm, font=\codefont\scriptsize\bfseries}]
    \node[font=\small\bfseries] at (-0.6,0.9) {مع خط الأنابيب: 4 تعليمات في 8 دورات};
    \foreach \c in {1,...,8} \node[font=\codefont\scriptsize] at (-0.95*\c+0.95*0.5-0.95*0, 0.35) {\textenglish{\c}};
    \foreach \i/\col in {1/brandTeal, 2/brandOrange, 3/brandGreen, 4/brandPurple} {
        \node[font=\small] at (0.9,-0.6*\i+0.1) {\textenglish{I\i}};
        \foreach \st [count=\k from 0] in {IF, ID, EX, MEM, WB} {
            \pgfmathsetmacro\x{-0.95*(\i+\k)+0.475}
            \node[s, fill=\col!25] at (\x, -0.6*\i+0.1) {\textenglish{\st}};
        }
    }
\end{tikzpicture}
\end{center}

\begin{tarif}
    لـ$n$ تعليمة و$k$ مرحلة (كل مرحلة دورة واحدة):
    \begin{itemize}
        \item دون خط أنابيب: $n \times k$ دورة.
        \item مع خط أنابيب: $k + (n - 1)$ دورة: الأولى تحتاج $k$ دورات، ثم تنتهي تعليمة في \textbf{كل} دورة.
    \end{itemize}
\end{tarif}

\begin{khatwa}[: 100 تعليمة في خط أنابيب من 5 مراحل]
    \begin{enumerate}
        \item دون خط أنابيب: $100 \times 5 = 500$ دورة.
        \item مع خط أنابيب: $5 + (100 - 1) = 104$ دورات.
        \item التسريع: $\dfrac{500}{104} \approx 4.8$ مرات، أي قريب من عدد المراحل (5).
    \end{enumerate}
\end{khatwa}

\subsection{المخاطر: متى يتوقف خط الأنابيب؟}
\begin{enumerate}
    \item \textbf{خطر المعطيات}: \code{add rax, rbx} ثم \code{add rcx, rax}. التعليمة الثانية تحتاج \textenglish{RAX} الجديد،
          لكن الأولى لم تكتبه بعد. الحل: المعالج «يمرر» النتيجة مباشرة من \textenglish{EX} إلى التعليمة الثانية، أو ينتظر دورة.
    \item \textbf{خطر التحكم}: بعد \code{jmp} أو \code{jne}، التعليمات التي جُلبت بعدها قد تكون خاطئة. الحل:
          \mustalah{التنبؤ بالقفز}{Prédiction de branchement}{Branch prediction}: المعالج يخمّن هل سيقفز، وإن أخطأ
          يرمي التعليمات الخاطئة ويبدأ من جديد.
    \item \textbf{خطر البنية}: تعليمتان تحتاجان الدارة نفسها في الدورة نفسها (مثل الذاكرة).
\end{enumerate}

\subsection{مثال: \code{mov} و\code{add} و\code{jmp} في 64 بتًا}
نتابع ثلاث تعليمات من البرنامج \code{same.asm} وبرامج هذا الفصل، ببايتاتها الحقيقية:

\begin{center}
\small
\begin{tabularx}{\linewidth}{|>{\centering\arraybackslash}p{1.1cm}|>{\centering\arraybackslash}p{2.7cm}|R|R|}
    \hline
    \ras{المرحلة} & \ras{\code{mov rax, 60}} & \ras{\code{add rax, rbx}} & \ras{\code{jmp fin}} \\ \hline
    \textenglish{IF} & جلب 7 بايتات \textenglish{48 C7 C0 3C 00 00 00} & جلب 3 بايتات \textenglish{48 01 D8} & جلب بايتين \textenglish{EB 00} \\ \hline
    \textenglish{ID} & \textenglish{REX.W}: 64 بتًا، الوجهة \textenglish{RAX}، والقيمة 60 داخل التعليمة & 64 بتًا، المصدر \textenglish{RBX}، الوجهة \textenglish{RAX}، وقراءة السجلين & قفز قصير، الإزاحة 0 \\ \hline
    \textenglish{EX} & لا حساب (القيمة جاهزة) & \textenglish{UAL}: $\text{RAX} + \text{RBX}$ وحساب الأعلام & حساب الهدف: $\text{RIP} + 0$ \\ \hline
    \textenglish{MEM} & لا شيء & لا شيء & لا شيء \\ \hline
    \textenglish{WB} & \nql{RAX}{60} & \textenglish{RAX} يأخذ النتيجة، وتحديث \textenglish{RFLAGS} & \textenglish{RIP} يأخذ الهدف، ورمي ما جُلب خطأً \\ \hline
\end{tabularx}
\end{center}

\begin{hal}
    معالج \textenglish{Core i7} الذي جربنا عليه لا يكتفي بخط أنابيب من 5 مراحل: فيه أكثر من 14 مرحلة،
    ويستطيع إطلاق \textbf{4 تعليمات في كل دورة}، بل وينفذها \textbf{بغير ترتيبها} إذا كانت مستقلة عن بعضها،
    ثم يرتب النتائج. كل هذا يحدث دون أن يرى المبرمج أي فرق.
\end{hal}

%================================================
\section{تمارين}
%================================================

\begin{tamrin}
    كم بايتًا يستطيع معالج أن يعنون بعنوان من 36 بتًا؟ (هذا ما سمح به الامتداد \textenglish{PAE} في معالجات 32 بتًا).
\end{tamrin}
\begin{hall}
    $2^{36} = 2^{6} \times 2^{30} = 64$ \textenglish{GB}.
\end{hall}

\begin{tamrin}
    \textenglish{RAX = \hex{FFFFFFFFFFFFFFFF}}. ما قيمة \textenglish{RAX} بعد: (أ) \code{mov eax, 1} (ب) \code{mov ax, 1} (ج) \code{mov al, 1}؟
\end{tamrin}
\begin{hall}
    (أ) \hex{0000000000000001} (الكتابة في 32 بتًا تصفّر النصف الأعلى).\\
    (ب) \hex{FFFFFFFFFFFF0001}. \quad (ج) \hex{FFFFFFFFFFFFFF01}.
\end{hall}

\begin{tamrin}
    جدول من أعداد \textenglish{int} (4 بايتات) يبدأ في \hex{601000}. ما عنوان العنصر رقم 10؟ واكتب العنونة بصيغة \textenglish{NASM}.
\end{tamrin}
\begin{hall}
    $\hex{601000} + 10 \times 4 = \hex{601000} + 40 = \hex{601028}$. الصيغة: \code{[tab + rcx*4]} مع \textenglish{RCX = 10}.
\end{hall}

\begin{tamrin}
    خط أنابيب من 6 مراحل، و1000 تعليمة. احسب عدد الدورات دون خط أنابيب ومعه، والتسريع.
\end{tamrin}
\begin{hall}
    دون: $1000 \times 6 = 6000$. مع: $6 + 999 = 1005$. التسريع $= \dfrac{6000}{1005} \approx 5.97$.
\end{hall}

\begin{tamrin}
    في البرنامج الآتي، أين يوجد خطر المعطيات؟ \code{mov rax, [x]} ثم \code{add rax, 5} ثم \code{mov [y], rax}.
\end{tamrin}
\begin{hall}
    الثانية تحتاج \textenglish{RAX} الذي تقرؤه الأولى من الذاكرة، والثالثة تحتاج \textenglish{RAX} الذي تحسبه الثانية: كل
    تعليمة تعتمد على التي قبلها مباشرة.
\end{hall}

\chapter{بنية المعالجات الحديثة}

\begin{khulasa}[: ماذا سنتعلم في هذا الفصل؟]
\begin{itemize}
    \item كيف تطورت المعالجات، ولماذا توقف ارتفاع التردد.
    \item المعالجات متعددة الأنوية والمعالجة المتوازية: الفوائد والصعوبات، وقانون \textenglish{Amdahl}.
    \item الذاكرة المخبئية وتسلسل الذاكرة، ولماذا يكون ترتيب قراءة المعطيات مهمًا جدًا.
\end{itemize}
\end{khulasa}

\begin{intabih}
    كل الأرقام في هذا الفصل \textbf{قسناها فعلًا} على حاسوب محمول بمعالج \textenglish{Intel Core i7-8665U}
    (4 أنوية، 8 خيوط، ذاكرة مخبئية \textenglish{L1 = 32 KB} و\textenglish{L2 = 256 KB} لكل نواة، و\textenglish{L3 = 8 MB} مشتركة).
    إذا أعدت التجارب على حاسوبك ستحصل على أرقام مختلفة قليلًا، لكن الظواهر نفسها.
\end{intabih}

%================================================
\section{المعالجات الحديثة: الاتجاهات والتطور}
%================================================

\subsection{خمسون سنة من التقدم}
\begin{center}
\small
\begin{tabular}{|c|c|c|c|c|}
    \hline
    \ras{المعالج} & \ras{السنة} & \ras{الترانزستورات} & \ras{التردد} & \ras{الأنوية} \\ \hline
    \textenglish{Intel 4004} & 1971 & \textenglish{2\,300} & \textenglish{0.74 MHz} & 1 \\ \hline
    \textenglish{Intel 8086} & 1978 & \textenglish{29\,000} & \textenglish{5--10 MHz} & 1 \\ \hline
    \textenglish{Intel 80386} & 1985 & \textenglish{275\,000} & \textenglish{16--33 MHz} & 1 \\ \hline
    \textenglish{Pentium 4} & 2000 & 42 مليون & \textenglish{1.3--3.8 GHz} & 1 \\ \hline
    \textenglish{Core i7-8665U} & 2019 & حوالي ملياري & \textenglish{1.9--4.8 GHz} & 4 \\ \hline
    معالجات 2024 & 2024 & عشرات المليارات & حوالي \textenglish{5 GHz} & من 8 إلى أكثر من 100 \\ \hline
\end{tabular}
\end{center}

\subsection{جدار التردد}
حتى سنة 2004 تقريبًا، كانت كل سنة تأتي بمعالج أسرع لأن \textbf{التردد} يرتفع. ثم توقف التردد عند حوالي
\textenglish{4} إلى \textenglish{5 GHz}. السبب الرئيسي هو \textbf{الحرارة}:
\begin{itemize}
    \item كل نبضة ساعة تستهلك طاقة. كلما ارتفع التردد ارتفعت الطاقة المستهلكة، والطاقة تتحول إلى حرارة.
    \item عند \textenglish{4 GHz} تقريبًا يسخن المعالج إلى درجة لا نستطيع تبريدها في حاسوب عادي.
\end{itemize}

\begin{tashbih}
    عدّاء يركض أسرع وأسرع، حتى يصل إلى سرعة لا يستطيع أن يحافظ عليها دون أن ينهار. بدل أن نطلب منه
    الركض أسرع، نضيف عدّائين آخرين يتقاسمون معه المسافة. هذا بالضبط ما فعله صانعو المعالجات.
\end{tashbih}

\subsection{الاتجاهات الحالية}
\begin{enumerate}
    \item \textbf{أنوية أكثر} بدل تردد أعلى (القسم الآتي).
    \item \textbf{ذاكرة مخبئية أكبر وأذكى}، لأن الذاكرة أصبحت أبطأ بكثير من المعالج (القسم الأخير).
    \item \textbf{تعليمات متخصصة}: \textenglish{SIMD} (الفصل الرابع)، وتسريع الذكاء الاصطناعي والتشفير.
    \item \textbf{أنوية مختلفة في الشريحة نفسها}: أنوية قوية للعمل الثقيل وأنوية اقتصادية للعمل الخفيف، لتوفير البطارية.
\end{enumerate}

%================================================
\section{المعالجات متعددة الأنوية والتوازي}
%================================================

\subsection{المفاهيم}
\begin{tarif}
    \begin{itemize}
        \item \mustalah{النواة}{Cœur}{Core}: معالج كامل (وحدة تحكم، وحدة حساب، سجلات، ذاكرة مخبئية \textenglish{L1}).
        \item \mustalah{المعالج متعدد الأنوية}{Processeur multi-cœur}{Multi-core processor}: عدة أنوية في شريحة واحدة.
        \item \mustalah{الخيط}{Thread}{Thread}: سلسلة من التعليمات تُنفذ بشكل مستقل. البرنامج قد يملك خيطًا واحدًا أو عدة خيوط.
        \item \mustalah{المعالجة المتوازية}{Traitement parallèle}{Parallel processing}: تنفيذ عدة أجزاء من العمل
              \textbf{في اللحظة نفسها} على أنوية مختلفة.
    \end{itemize}
\end{tarif}

\begin{center}
\begin{tikzpicture}[core/.style={unit, minimum width=2.4cm, minimum height=1.5cm}]
    \draw[brandDark, very thick, rounded corners=8pt, fill=brandLight] (-6.1,-2.6) rectangle (6.1,2.1);
    \node[font=\bfseries] at (0,1.75) {شريحة معالج من 4 أنوية};
    \foreach \x/\i in {4.5/1, 1.5/2, -1.5/3, -4.5/4} {
        \node[core] at (\x,0.6) {النواة \i\\{\scriptsize \textenglish{L1}، \textenglish{L2}}};
    }
    \node[unit=brandOrange, minimum width=11.2cm] at (0,-1.1) {الذاكرة المخبئية المشتركة \textenglish{L3}};
    \node[unit=brandPurple, minimum width=4cm] (ram) at (0,-3.6) {الذاكرة الرئيسية \textenglish{RAM}};
    \draw[sahmT, <->] (0,-1.6) -- (ram);
\end{tikzpicture}
\end{center}

\begin{tashbih}
    مطبخ فيه طباخ واحد (نواة واحدة) يحضّر الطبق تلو الآخر. مطبخ فيه 4 طباخين (4 أنوية) يحضّر 4 أطباق
    في الوقت نفسه. لكنهم يتقاسمون الثلاجة نفسها (\textenglish{L3} و\textenglish{RAM})، وقد ينتظر أحدهم دوره
    أمامها. وطبق واحد معقد لا يصبح أسرع إلا إذا استطعنا تقسيمه إلى مهام مستقلة.
\end{tashbih}

\subsection{\textenglish{Hyper-Threading}: نواة واحدة لخيطين}
تقنية \mustalah{تعدد الخيوط المتزامن}{Multithreading simultané}{SMT, Hyper-Threading} تجعل النواة الواحدة تبدو
كأنها نواتان: عندما ينتظر خيط معطيات من الذاكرة، يستعمل الخيط الثاني وحدات الحساب الفارغة. لهذا يُظهر
معالجنا 8 «معالجات منطقية» رغم أنه يملك 4 أنوية فقط. لكن إذا كان الخيطان يحتاجان وحدة الحساب نفسها
دائمًا، فلا نربح شيئًا تقريبًا، كما سنرى في التجربة.

\subsection{تجربة: العمل نفسه على 1 و2 و4 و8 خيوط}
البرنامج الآتي يحسب مجموع الجذور التربيعية لـ400 مليون عدد، ويقسم العمل بين $k$ خيوط:

\cfile{ch5/threads.c}
\screen{ch5/threads.out}

\begin{sharh}
    \s{18--26}{كل خيط يأخذ الأعداد $id$، $id + k$، $id + 2k$\dots فيحسب جزءًا من المجموع لا يتداخل مع الآخرين.}
    \s{34--35}{\code{pthread\_create}: نطلق $k$ خيوط، يعمل كل واحد منها على نواة.}
    \s{36--39}{\code{pthread\_join}: ننتظر انتهاء كل خيط، ثم نجمع النتائج الجزئية.}
    \s{40--42}{التسريع $=$ الزمن بخيط واحد $\div$ الزمن بـ$k$ خيوط.}
\end{sharh}

\begin{center}
\begin{tikzpicture}[x=1.1cm, y=1.1cm]
    \draw[->, thick] (0,0) -- (8.6,0) node[right, font=\small] {الخيوط};
    \draw[->, thick] (0,0) -- (0,4.4) node[above, font=\small] {التسريع};
    \foreach \y in {1,2,3,4} \draw (0.08,\y) -- (-0.08,\y) node[left, font=\scriptsize] {\textenglish{\y}};
    \foreach \x in {1,2,4,8} \draw (\x,0.08) -- (\x,-0.08) node[below, font=\scriptsize] {\textenglish{\x}};
    \draw[brandMuted, dashed, thick] (1,1) -- (4,4) node[above, font=\scriptsize, lawn=brandMuted] {التسريع المثالي};
    \draw[orangeDeep, very thick] (1,1) -- (2,1.98) -- (4,3.02) -- (8,3.05);
    \foreach \x/\y in {1/1, 2/1.98, 4/3.02, 8/3.05} {
        \fill[orangeDeep] (\x,\y) circle (2.5pt);
        \node[font=\codefont\scriptsize, above left] at (\x,\y) {\textenglish{\y}};
    }
    \draw[brandTeal, dashed] (4,0) -- (4,4.2);
    \node[font=\scriptsize, lawn=tealDeep, align=center] at (6,3.9) {4 أنوية فقط:\\الخيوط الإضافية لا تساعد هنا};
\end{tikzpicture}
\end{center}

\begin{khatwa}[: كيف نقرأ هذه النتائج؟]
    \begin{enumerate}
        \item \textbf{خيطان}: التسريع \textenglish{1.98}، أي قريب جدًا من 2. العمل انقسم إلى نصفين مستقلين تمامًا.
        \item \textbf{4 خيوط}: التسريع \textenglish{3.02} وليس 4. عندما تعمل الأنوية الأربعة معًا يسخن المعالج،
              فيخفض تردده تلقائيًا. وجزء من البرنامج (إطلاق الخيوط وجمع النتائج) لا يتوازى.
        \item \textbf{8 خيوط}: التسريع \textenglish{3.05}، تقريبًا مثل 4 خيوط. معالجنا فيه 4 أنوية حقيقية فقط،
              والخيوط الإضافية تتقاسم وحدات الحساب نفسها، فلا يتقدم العمل أسرع.
    \end{enumerate}
\end{khatwa}

\subsection{قانون \textenglish{Amdahl}: حدود التوازي}
معظم البرامج فيها جزء \textbf{لا يمكن توزيعه} (قراءة الملف، عرض النتيجة، انتظار المستعمل\dots).

\begin{tarif}
    إذا كان الجزء القابل للتوازي من البرنامج هو $p$ (بين 0 و1)، ولدينا $n$ نواة، فإن التسريع الأقصى هو:
    \[ S(n) = \frac{1}{(1 - p) + \dfrac{p}{n}} \]
\end{tarif}

\begin{khatwa}[: برنامج 90\% منه قابل للتوازي]
    $p = 0.9$، والجزء التسلسلي $1 - p = 0.1$.
    \begin{enumerate}
        \item بـ4 أنوية: $S = \dfrac{1}{0.1 + \frac{0.9}{4}} = \dfrac{1}{0.1 + 0.225} = \dfrac{1}{0.325} \approx 3.08$.
        \item بـ16 نواة: $S = \dfrac{1}{0.1 + 0.05625} \approx 6.4$.
        \item بعدد لا نهائي من الأنوية: $\dfrac{p}{n} \rightarrow 0$، فيصبح $S = \dfrac{1}{0.1} = 10$.
    \end{enumerate}
    مهما أضفنا أنوية، لن يكون البرنامج أسرع من 10 مرات، لأن الـ10\% التسلسلية تبقى كما هي.
\end{khatwa}

\begin{hal}
    لاحظ أن قانون \textenglish{Amdahl} مع $p = 0.9$ و4 أنوية يعطي \textenglish{3.08}، وهو قريب جدًا من التسريع
    \textenglish{3.02} الذي قسناه! القانون البسيط يشرح الواقع جيدًا.
\end{hal}

\subsection{الفوائد والتحديات}
\begin{center}
\small
\begin{tabularx}{\linewidth}{|R|R|}
    \hline
    \ras{الفوائد} & \ras{التحديات} \\ \hline
    أداء أكبر دون رفع التردد & يجب أن يُكتب البرنامج بخيوط متعددة، وهذا أصعب \\ \hline
    استهلاك طاقة أقل لنفس العمل & الجزء التسلسلي يحد من التسريع (قانون \textenglish{Amdahl}) \\ \hline
    تشغيل عدة برامج في الوقت نفسه بسلاسة & \textbf{التزامن}: خيطان يكتبان في المتغير نفسه قد يفسدان النتيجة \\ \hline
    إذا تعطلت مهمة، تبقى الأنوية الأخرى تعمل & \textbf{تماسك الذاكرة المخبئية}: يجب أن ترى كل الأنوية القيمة نفسها للمتغير \\ \hline
\end{tabularx}
\end{center}

\begin{mithal}[: مشكلة التزامن]
    خيطان ينفذان \code{x = x + 1} على المتغير نفسه، و\textenglish{x = 5} في البداية:
    \begin{enumerate}
        \item الخيط 1 يقرأ \textenglish{x = 5}.
        \item الخيط 2 يقرأ \textenglish{x = 5} أيضًا (قبل أن يكتب الأول).
        \item الخيط 1 يحسب $5 + 1$ ويكتب 6.
        \item الخيط 2 يحسب $5 + 1$ ويكتب 6 أيضًا.
    \end{enumerate}
    النتيجة 6 بدل 7: ضاعت زيادة واحدة! لهذا نستعمل \textbf{الأقفال} \textenglish{(mutex)} أو تعليمات \textbf{ذرّية}
    (\code{lock add} في \textenglish{x86}) تمنع خيطين من تعديل المتغير في الوقت نفسه. في برنامجنا تجنبنا المشكلة:
    كل خيط يكتب في خانته الخاصة \code{part[id]}.
\end{mithal}

%================================================
\section{الذاكرة المخبئية وتسلسل الذاكرة}
%================================================

\subsection{المشكلة: الذاكرة بطيئة جدًا}
المعالج قادر على تنفيذ عدة تعليمات في كل نانوثانية، لكن قراءة الذاكرة الرئيسية تحتاج حوالي 100 نانوثانية.
لو انتظر المعالج الذاكرة في كل مرة، لبقي معظم الوقت دون عمل.

\begin{tarif}
    \mustalah{الذاكرة المخبئية}{Mémoire cache, antémémoire}{Cache memory}: ذاكرة صغيرة وسريعة جدًا (من نوع
    \textenglish{SRAM}) قريبة من المعالج، تحفظ \textbf{نسخة} من المعطيات التي استُعملت مؤخرًا أو التي سنحتاجها قريبًا.
\end{tarif}

\begin{tashbih}
    عندما تراجع للامتحان، لا تضع كل كتب المكتبة على مكتبك. تضع على المكتب الكتب التي تحتاجها الآن (\textenglish{L1})،
    وفي الحقيبة كتب الأسبوع (\textenglish{L2})، وفي خزانة البيت كتب السنة (\textenglish{L3})، والباقي في المكتبة
    البعيدة (\textenglish{RAM}). كلما كان المكان أقرب، كان أصغر وأسرع.
\end{tashbih}

\subsection{تسلسل الذاكرة}
\begin{center}
\begin{tikzpicture}
    \foreach \y/\w/\t/\c/\d in {
        3.6/1.6/{السجلات}/brandOrange/{أقل من \textenglish{1 KB}، أقل من \textenglish{1 ns}},
        2.7/2.8/{\textenglish{L1}}/brandOrange!70!brandTeal/{\textenglish{32 KB}، حوالي \textenglish{1 ns}},
        1.8/4.0/{\textenglish{L2}}/brandTeal/{\textenglish{256 KB}، حوالي \textenglish{4 ns}},
        0.9/5.2/{\textenglish{L3}}/brandTeal!70!brandPurple/{\textenglish{8 MB}، حوالي \textenglish{10 ns}},
        0.0/6.4/{الذاكرة الرئيسية \textenglish{RAM}}/brandPurple/{\textenglish{8 GB}، حوالي \textenglish{100 ns}},
        -0.9/7.6/{القرص \textenglish{SSD}}/brandMuted/{\textenglish{256 GB}، حوالي \textenglish{100 µs}}} {
        \fill[\c!35, draw=\c, thick] (-\w/2,\y) rectangle (\w/2,\y+0.8);
        \node[font=\small\bfseries] at (0,\y+0.4) {\t};
        \node[font=\scriptsize, anchor=east] at (-4.1,\y+0.4) {\d};
    }
    \draw[sahmO] (4.3,-0.9) -- (4.3,4.4) node[midway, right, font=\small, align=left] {أسرع\\أصغر\\أغلى};
\end{tikzpicture}
\end{center}

\subsection{تجربة: نرى طبقات الذاكرة بأعيننا}
البرنامج الآتي يقرأ جدولًا بترتيب عشوائي (كل قراءة تعطي عنوان القراءة التالية، فلا يستطيع المعالج التخمين)،
ويقيس متوسط زمن القراءة الواحدة، لأحجام مختلفة من الجدول:

\cfile{ch5/latency.c}
\screen{ch5/latency.out}

\begin{center}
\begin{tikzpicture}[x=1.45cm, y=0.042cm]
    \draw[->, thick] (0,0) -- (7.4,0) node[right, font=\small] {حجم المعطيات};
    \draw[->, thick] (0,0) -- (0,125) node[above, font=\small] {الزمن \textenglish{(ns)}};
    \foreach \y in {0,25,50,75,100} \draw (0.06,\y) -- (-0.06,\y) node[left, font=\scriptsize] {\textenglish{\y}};
    \foreach \x/\t in {1/16 KB, 2/64 KB, 3/256 KB, 4/1 MB, 5/4 MB, 6/16 MB, 7/64 MB}
        \node[font=\codefont\scriptsize, below] at (\x,-2) {\textenglish{\t}};
    \fill[brandOrange!15] (0.5,0) rectangle (1.6,125); \node[font=\scriptsize] at (1.05,118) {\textenglish{L1}};
    \fill[brandTeal!15] (1.6,0) rectangle (3.4,125); \node[font=\scriptsize] at (2.5,118) {\textenglish{L2}};
    \fill[brandPurple!12] (3.4,0) rectangle (5.6,125); \node[font=\scriptsize] at (4.5,118) {\textenglish{L3}};
    \fill[brandMuted!12] (5.6,0) rectangle (7.3,125); \node[font=\scriptsize] at (6.45,118) {\textenglish{RAM}};
    \draw[orangeDeep, very thick] (1,1.7) -- (2,2.2) -- (3,4.1) -- (4,9.2) -- (5,12.8) -- (6,83.5) -- (7,112.7);
    \foreach \x/\y in {1/1.7, 2/2.2, 3/4.1, 4/9.2, 5/12.8, 6/83.5, 7/112.7} {
        \fill[orangeDeep] (\x,\y) circle (2.5pt);
        \node[font=\codefont\scriptsize, above] at (\x,\y+2) {\textenglish{\y}};
    }
\end{tikzpicture}
\end{center}

\begin{khatwa}[: ماذا نرى في المنحنى؟]
    \begin{enumerate}
        \item \textbf{حتى \textenglish{32 KB}}: الجدول كله في \textenglish{L1}، والقراءة تستغرق \textenglish{1.7 ns} فقط.
        \item \textbf{\textenglish{64 KB} و\textenglish{256 KB}}: الجدول لم يعد يدخل في \textenglish{L1}، فنقرأ من \textenglish{L2}: الزمن يرتفع إلى \textenglish{2} ثم \textenglish{4 ns}.
        \item \textbf{\textenglish{1 MB} و\textenglish{4 MB}}: نقرأ من \textenglish{L3} المشتركة: حوالي \textenglish{9} إلى \textenglish{13 ns}.
        \item \textbf{\textenglish{16 MB} فأكثر}: الجدول أكبر من \textenglish{L3} (\textenglish{8 MB})، فنذهب إلى \textenglish{RAM}:
              الزمن يقفز إلى \textenglish{84} ثم \textenglish{113 ns}، أي أبطأ بحوالي \textbf{65 مرة} من \textenglish{L1}!
    \end{enumerate}
    الدرجات في المنحنى تقع بالضبط عند أحجام الذاكرات المخبئية لمعالجنا.
\end{khatwa}

\subsection{مبدأ المحلية: لماذا تنجح الذاكرة المخبئية؟}
الذاكرة المخبئية صغيرة، لكنها تحتوي في الغالب على ما نحتاجه، بفضل سلوكين طبيعيين في البرامج:
\begin{itemize}
    \item \mustalah{المحلية الزمنية}{Localité temporelle}{Temporal locality}: المعطاة التي استعملناها الآن سنستعملها
          غالبًا \textbf{مرة أخرى قريبًا} (عدّاد الحلقة، المجموع\dots).
    \item \mustalah{المحلية المكانية}{Localité spatiale}{Spatial locality}: بعد قراءة عنوان، نقرأ غالبًا
          \textbf{العناوين المجاورة} (عناصر جدول، تعليمات متتالية).
\end{itemize}

\begin{tarif}
    \mustalah{سطر الذاكرة المخبئية}{Ligne de cache}{Cache line}: الذاكرة المخبئية لا تنقل بايتًا واحدًا، بل
    كتلة كاملة من \textbf{64 بايتًا} مجاورة. عندما نقرأ \code{t[0]}، تصل معه \code{t[1]} إلى \code{t[15]} (للأعداد من 4 بايتات) مجانًا.
\end{tarif}

\subsection{تجربة: ترتيب القراءة يغيّر كل شيء}
\cfile{ch5/locality.c}
\screen{ch5/locality.out}

\begin{center}
\begin{tikzpicture}[c/.style={draw=brandDark, minimum size=0.55cm}]
    \foreach \i in {0,...,3} \foreach \j in {0,...,3} \node[c] at (-\j*0.55-4.2, -\i*0.55) {};
    \foreach \i in {0,...,3} \draw[sahmT] (-6.0, -\i*0.55) -- (-4.0, -\i*0.55);
    \node[font=\small, align=center] at (-5.0,0.75) {سطرًا بسطر: كل قراءة\\بجوار التي قبلها};
    \foreach \i in {0,...,3} \foreach \j in {0,...,3} \node[c] at (\j*0.55+3.0, -\i*0.55) {};
    \foreach \j in {0,...,3} \draw[sahmO] (\j*0.55+3.0, 0.2) -- (\j*0.55+3.0, -1.85);
    \node[font=\small, align=center] at (3.8,0.75) {عمودًا بعمود: كل قراءة\\تقفز \textenglish{16 KB}};
\end{tikzpicture}
\end{center}

\begin{khatwa}[: لماذا الأعمدة أبطأ بـ33 مرة؟]
    \begin{enumerate}
        \item في لغة \textenglish{C}، الجدول \code{m[4096][4096]} محفوظ \textbf{سطرًا بعد سطر}: \code{m[0][0]} ثم \code{m[0][1]} ثم\dots
        \item \textbf{سطرًا بسطر}: كل سطر ذاكرة مخبئية (64 بايتًا) يحتوي على 16 عددًا متجاورًا. نقرأ من \textenglish{RAM} مرة، ثم نستعمل
              16 عددًا من الذاكرة المخبئية. والمعالج يلاحظ أننا نتقدم بانتظام فيجلب الأسطر القادمة مسبقًا.
        \item \textbf{عمودًا بعمود}: \code{m[0][j]} و\code{m[1][j]} بينهما $4096 \times 4 = 16\,384$ بايتًا. كل قراءة
              في سطر ذاكرة مخبئية جديد، فنستعمل عددًا واحدًا فقط من كل 16 عددًا نقلناها.
        \item النتيجة الحقيقية: \textenglish{0.009 s} مقابل \textenglish{0.301 s}، أي حوالي \textbf{33 مرة} أبطأ،
              لنفس العمليات تمامًا!
    \end{enumerate}
\end{khatwa}

\subsection{الإصابة والإخفاق وزمن الوصول المتوسط}
\begin{tarif}
    \begin{itemize}
        \item \mustalah{الإصابة}{Succès}{Hit}: المعطاة موجودة في الذاكرة المخبئية.
        \item \mustalah{الإخفاق}{Défaut}{Miss}: المعطاة غير موجودة، فنذهب إلى المستوى التالي (الأبطأ).
        \item \textbf{معدل الإصابة} $h$: نسبة القراءات التي وجدت معطياتها في الذاكرة المخبئية.
    \end{itemize}
    \[ \text{الزمن المتوسط} = h \times t_{cache} + (1 - h) \times t_{RAM} \]
\end{tarif}

\begin{khatwa}[: \textenglish{$t_{cache} = 1$ ns}، \textenglish{$t_{RAM} = 100$ ns}]
    \begin{enumerate}
        \item معدل إصابة 90\%: $0.9 \times 1 + 0.1 \times 100 = 0.9 + 10 = 10.9$ \textenglish{ns}.
        \item معدل إصابة 99\%: $0.99 \times 1 + 0.01 \times 100 = 0.99 + 1 = 1.99$ \textenglish{ns}.
    \end{enumerate}
    الانتقال من 90\% إلى 99\% يجعل الذاكرة \textbf{أسرع بأكثر من 5 مرات}. لهذا تهتم المعالجات الحديثة جدًا
    بتحسين معدل الإصابة.
\end{khatwa}

\subsection{أين نضع المعطاة في الذاكرة المخبئية؟ (التطابق)}
الذاكرة المخبئية مقسمة إلى أسطر. علينا أن نقرر في أي سطر نضع كتلة قادمة من الذاكرة:

\begin{center}
\small
\begin{tabularx}{\linewidth}{|>{\centering\arraybackslash}p{3.4cm}|R|R|}
    \hline
    \ras{النوع} & \ras{القاعدة} & \ras{الميزة والعيب} \\ \hline
    التطابق المباشر \textenglish{(direct mapped)} & كل كتلة لها سطر واحد ممكن: رقم الكتلة \textenglish{mod} عدد الأسطر & بسيط وسريع، لكن كتلتان قد تتصارعان على السطر نفسه \\ \hline
    التجميعي الكامل \textenglish{(fully associative)} & أي كتلة في أي سطر & أقل صراعات، لكن البحث في كل الأسطر مكلف \\ \hline
    التجميعي بمجموعات \textenglish{(set associative)} & الكتلة تذهب إلى مجموعة محددة، وفيها إلى أي سطر من $k$ أسطر & الحل الوسط المستعمل في كل المعالجات الحديثة \\ \hline
\end{tabularx}
\end{center}

\begin{khatwa}[: تقسيم عنوان في ذاكرة مخبئية بتطابق مباشر]
    ذاكرة مخبئية سعتها \textenglish{1 KB}، وأسطرها من 16 بايتًا، والعناوين من 16 بتًا.
    \begin{enumerate}
        \item عدد الأسطر: $1024 \div 16 = 64 = 2^6$ سطرًا، إذن نحتاج \textbf{6 بتات} لرقم السطر.
        \item داخل السطر 16 بايتًا $= 2^4$، إذن \textbf{4 بتات} للإزاحة.
        \item الباقي: $16 - 6 - 4 = 6$ بتات، وتسمى \textbf{الوسم} \textenglish{(tag)}: تحدد أي كتلة من الذاكرة موجودة الآن في السطر.
    \end{enumerate}
    \begin{center}
    \begin{tikzpicture}[f/.style={draw=brandDark, thick, minimum height=0.7cm, font=\small\bfseries}]
        % drawn left to right, like the binary address: tag first
        \node[f, minimum width=3cm, fill=brandOrange!20] (t) at (0,0) {الوسم: 6 بتات};
        \node[f, minimum width=3cm, fill=brandTeal!20, anchor=west] (l) at (t.east) {السطر: 6 بتات};
        \node[f, minimum width=2cm, fill=brandGreen!20, anchor=west] (o) at (l.east) {الإزاحة: 4 بتات};
        \node[font=\codefont\scriptsize, below=1pt of t.south west, anchor=north west] {\textenglish{bit 15}};
        \node[font=\codefont\scriptsize, below=1pt of o.south east, anchor=north east] {\textenglish{bit 0}};
    \end{tikzpicture}
    \end{center}
    مثال: $\hex{1A2C} = \bin{000110\,100010\,1100}$. الوسم $\bin{000110} = 6$، والسطر $\bin{100010} = 34$، والإزاحة $\bin{1100} = 12$.
\end{khatwa}

\begin{intabih}
    العنوان الثنائي يُكتب دائمًا من اليسار إلى اليمين، حتى في كتاب بالعربية: البتات العليا (الوسم) على اليسار
    والبت 0 على اليمين، تمامًا كما رسمناه في الشكل.
\end{intabih}

\subsection{الاستبدال والكتابة}
\begin{itemize}
    \item \textbf{سياسة الاستبدال}: عندما تمتلئ المجموعة، أي سطر نطرد؟ غالبًا \textenglish{LRU} \textenglish{(Least Recently Used)}:
          السطر الذي لم نستعمله منذ أطول مدة.
    \item \textbf{سياسة الكتابة}:
    \begin{itemize}
        \item \textenglish{Write-through}: نكتب في الذاكرة المخبئية وفي \textenglish{RAM} معًا. بسيط لكنه بطيء.
        \item \textenglish{Write-back}: نكتب في الذاكرة المخبئية فقط، ولا نكتب في \textenglish{RAM} إلا عندما نطرد السطر. أسرع، وهو المستعمل اليوم.
    \end{itemize}
\end{itemize}

\begin{hal}
    تذكر نموذج هارفارد من الفصل الأول: ذاكرة \textenglish{L1} مقسمة إلى جزأين، \textenglish{L1i} للتعليمات
    و\textenglish{L1d} للمعطيات (في معالجنا \textenglish{32 KB} لكل واحدة)، فيستطيع المعالج جلب تعليمة وقراءة معطاة
    في الوقت نفسه. أما \textenglish{L2} و\textenglish{L3} و\textenglish{RAM} فمشتركة بينهما، كما في نموذج فون نيومان.
\end{hal}

%================================================
\section{تمارين}
%================================================

\begin{tamrin}
    برنامج 75\% منه قابل للتوازي. احسب التسريع بـ2 و4 و8 أنوية، والتسريع الأقصى.
\end{tamrin}
\begin{hall}
    $1 - p = 0.25$. \quad $S(2) = \dfrac{1}{0.25 + 0.375} = 1.6$، \quad $S(4) = \dfrac{1}{0.25 + 0.1875} \approx 2.29$،
    \quad $S(8) = \dfrac{1}{0.25 + 0.09375} \approx 2.91$. الأقصى $= \dfrac{1}{0.25} = 4$.
\end{hall}

\begin{tamrin}
    ذاكرة مخبئية زمنها \textenglish{2 ns}، والذاكرة الرئيسية \textenglish{80 ns}، ومعدل الإصابة 95\%. ما الزمن المتوسط؟
\end{tamrin}
\begin{hall}
    $0.95 \times 2 + 0.05 \times 80 = 1.9 + 4 = 5.9$ \textenglish{ns}.
\end{hall}

\begin{tamrin}
    ذاكرة مخبئية بتطابق مباشر: \textenglish{4 KB}، أسطر من 32 بايتًا، عناوين من 32 بتًا. كم بتًا للإزاحة وللسطر وللوسم؟
\end{tamrin}
\begin{hall}
    الإزاحة: $32 = 2^5$، إذن 5 بتات. الأسطر: $4096 \div 32 = 128 = 2^7$، إذن 7 بتات. الوسم: $32 - 7 - 5 = 20$ بتًا.
\end{hall}

\begin{tamrin}
    اشرح لماذا يكون هذا الكود بطيئًا، وكيف نصلحه:
    \begin{codebox}\begin{english}\codefont\small
    for (j = 0; j < N; j++) \\
    \hspace*{2em}for (i = 0; i < N; i++) \\
    \hspace*{4em}s += m[i][j];
    \end{english}\end{codebox}
\end{tamrin}
\begin{hall}
    يقرأ الجدول عمودًا بعمود، فكل قراءة في سطر ذاكرة مخبئية مختلف (لا محلية مكانية). نصلحه بتبديل الحلقتين:
    \code{for (i ...) for (j ...)}، فنقرأ سطرًا بسطر.
\end{hall}

\begin{tamrin}
    لماذا لم يتحسن البرنامج \code{threads.c} عند الانتقال من 4 إلى 8 خيوط في معالجنا؟
\end{tamrin}
\begin{hall}
    لأن المعالج يملك 4 أنوية حقيقية فقط. الخيوط 5 إلى 8 تعمل بتقنية \textenglish{Hyper-Threading} على الأنوية نفسها،
    والبرنامج يستعمل وحدة الحساب طوال الوقت، فلا توجد وحدات فارغة يستفيد منها الخيط الثاني.
\end{hall}

\appendix

%================================================
\chapter{ملخص تعليمات \textenglish{8086}}
%================================================

هذا الملحق يجمع في مكان واحد كل التعليمات التي درسناها في الفصل الثالث، لتعود إليها بسرعة أثناء كتابة
برامجك. الرموز: \textenglish{d} الوجهة، \textenglish{s} المصدر، \textenglish{r} سجل، \textenglish{m} خانة ذاكرة.

\section*{وظائف \textenglish{DOS} الأكثر استعمالًا (\code{int 21h})}
\begin{center}
\small
\begin{tabular}{|c|l|l|}
    \hline
    \ras{\textenglish{AH}} & \ras{الوظيفة} & \ras{المدخلات / المخرجات} \\ \hline
    \hex{01} & قراءة حرف مع إظهاره & المخرج في \textenglish{AL} \\ \hline
    \hex{02} & كتابة حرف & الحرف في \textenglish{DL} \\ \hline
    \hex{09} & كتابة سلسلة تنتهي بـ\code{\$} & العنوان في \textenglish{DX} \\ \hline
    \hex{0A} & قراءة سطر في ذاكرة مؤقتة & العنوان في \textenglish{DX} \\ \hline
    \hex{4C} & إنهاء البرنامج & الرمز في \textenglish{AL} \\ \hline
\end{tabular}
\end{center}

\section*{التعليمات}

\begin{center}
\small\renewcommand{\arraystretch}{1.12}
\begin{tabularx}{\linewidth}{|l|>{\raggedleft\arraybackslash}X|l|}
    \hline
    \ras{التعليمة} & \ras{المعنى} & \ras{الأعلام المتأثرة} \\ \hline
    \multicolumn{3}{|c|}{\cellcolor{brandOrange!20}\textbf{النقل}} \\ \hline
    \code{mov d, s} & \nql{d}{s} & لا شيء \\ \hline
    \code{xchg a, b} & تبادل \textenglish{a} و\textenglish{b} & لا شيء \\ \hline
    \code{lea r, m} & \textenglish{r} يأخذ عنوان \textenglish{m} & لا شيء \\ \hline
    \code{push s} & \nql{SP}{SP - 2}، ثم \nql{[SP]}{s} & لا شيء \\ \hline
    \code{pop d} & \nql{d}{[SP]}، ثم \nql{SP}{SP + 2} & لا شيء \\ \hline
    \code{in al, port} / \code{out port, al} & قراءة من منفذ / كتابة في منفذ & لا شيء \\ \hline
    \multicolumn{3}{|c|}{\cellcolor{brandOrange!20}\textbf{الحساب}} \\ \hline
    \code{add d, s} & \nql{d}{d + s} & كلها \\ \hline
    \code{adc d, s} & \nql{d}{d + s + CF} & كلها \\ \hline
    \code{sub d, s} & \nql{d}{d - s} & كلها \\ \hline
    \code{sbb d, s} & \nql{d}{d - s - CF} & كلها \\ \hline
    \code{inc d} & \nql{d}{d + 1} & كلها إلا \textenglish{CF} \\ \hline
    \code{dec d} & \nql{d}{d - 1} & كلها إلا \textenglish{CF} \\ \hline
    \code{neg d} & \nql{d}{-d} & كلها \\ \hline
    \code{cmp a, b} & يحسب \textenglish{a - b} دون حفظ & كلها \\ \hline
    \code{mul s} (8 بتات) & \nql{AX}{AL * s} & \textenglish{CF}، \textenglish{OF} \\ \hline
    \code{mul s} (16 بتًا) & \nql{DX:AX}{AX * s} & \textenglish{CF}، \textenglish{OF} \\ \hline
    \code{div s} (8 بتات) & \nql{AL}{AX / s}، \nql{AH}{AX mod s} & غير محددة \\ \hline
    \code{div s} (16 بتًا) & \nql{AX}{DX:AX / s}، \nql{DX}{DX:AX mod s} & غير محددة \\ \hline
    \multicolumn{3}{|c|}{\cellcolor{brandOrange!20}\textbf{المنطق والإزاحة}} \\ \hline
    \code{and}، \code{or}، \code{xor} & عمليات بتًا ببت & \textenglish{CF = OF = 0}، \textenglish{ZF}، \textenglish{SF} \\ \hline
    \code{not d} & قلب كل البتات & لا شيء \\ \hline
    \code{test a, b} & \code{and} دون حفظ & مثل \code{and} \\ \hline
    \code{shl d, n} / \code{shr d, n} & إزاحة: ضرب أو قسمة على $2^n$ & \textenglish{CF} = آخر بت خرج \\ \hline
    \code{rol d, n} / \code{ror d, n} & تدوير & \textenglish{CF} \\ \hline
    \multicolumn{3}{|c|}{\cellcolor{brandOrange!20}\textbf{التحكم في التدفق}} \\ \hline
    \code{jmp L} & قفز دائمًا & لا شيء \\ \hline
    \code{je}، \code{jne}، \code{jz}، \code{jnz} & حسب \textenglish{ZF} & -- \\ \hline
    \code{jg}، \code{jge}، \code{jl}، \code{jle} & مقارنة \textbf{بإشارة} & -- \\ \hline
    \code{ja}، \code{jae}، \code{jb}، \code{jbe} & مقارنة \textbf{دون إشارة} & -- \\ \hline
    \code{loop L} & \nql{CX}{CX - 1}، ثم القفز إذا كان \textenglish{CX} $\neq 0$ & لا شيء \\ \hline
    \code{jcxz L} & القفز إذا كان \textenglish{CX = 0} & -- \\ \hline
    \code{call P} / \code{ret} & استدعاء إجراء / العودة منه & لا شيء \\ \hline
    \code{int n} / \code{iret} & مقاطعة / العودة منها & \textenglish{IF = TF = 0} / تُستعاد \\ \hline
    \multicolumn{3}{|c|}{\cellcolor{brandOrange!20}\textbf{السلاسل}} \\ \hline
    \code{movsb} / \code{movsw} & \nql{ES:[DI]}{DS:[SI]}، ثم تقدم \textenglish{SI} و\textenglish{DI} & لا شيء \\ \hline
    \code{rep} & كرر التعليمة \textenglish{CX} مرة & -- \\ \hline
    \code{cld} / \code{std} & \textenglish{DF = 0} (تصاعدي) / \textenglish{DF = 1} & \textenglish{DF} \\ \hline
\end{tabularx}
\end{center}

%================================================
\chapter{\textenglish{EMU8086} و\textenglish{NASM}: الفروق في الكتابة}
%================================================

برامج الفصل الثالث مكتوبة لـ\textenglish{EMU8086} (صيغة قريبة من \textenglish{MASM})، وبرامج الفصل الرابع لـ\textenglish{NASM}.
المعالج هو نفسه، لكن المجمِّعين لا يكتبان بعض الأشياء بالطريقة نفسها:

\begin{center}
\small
\begin{tabularx}{\linewidth}{|>{\raggedleft\arraybackslash}X|l|l|}
    \hline
    \ras{ما نريد} & \ras{\textenglish{EMU8086 / MASM}} & \ras{\textenglish{NASM}} \\ \hline
    عنوان متغير في سجل & \code{lea dx, msg} أو \code{mov dx, offset msg} & \code{mov dx, msg} \\ \hline
    قيمة متغير & \code{mov al, x} أو \code{mov al, [x]} & \code{mov al, [x]} (الأقواس إجبارية) \\ \hline
    تحديد الحجم & \code{mov byte ptr [bx], 5} & \code{mov byte [bx], 5} \\ \hline
    عدد ست عشري & \code{0FFh} & \code{0xFF} أو \code{0FFh} \\ \hline
    إجراء & \code{p proc} \dots{} \code{p endp} & تسمية بسيطة \code{p:} \\ \hline
    ثابت & \code{N equ 10} & \code{N equ 10} \\ \hline
    جدول من 10 أصفار & \code{t db 10 dup(0)} & \code{t times 10 db 0} \\ \hline
    التجميع & زر \textenglish{Emulate} & \code{nasm -f bin prog.asm -o prog.com} \\ \hline
\end{tabularx}
\end{center}

\begin{intabih}
    عادة جيدة: اكتب دائمًا الأقواس \code{[ ]} عندما تقصد محتوى الذاكرة، حتى في \textenglish{EMU8086}.
    هكذا يبقى برنامجك واضحًا، ويعمل مع المجمِّعين بتعديلات قليلة. هذا ما فعلناه في كل برامج هذا الدرس.
\end{intabih}

%================================================
\chapter{قاموس المصطلحات}
%================================================

\begin{center}
\small
\begin{tabular}{|l|l|l|}
    \hline
    \ras{بالعربية} & \ras{\textenglish{Français}} & \ras{\textenglish{English}} \\ \hline
    الحاسوب & \textenglish{Ordinateur} & \textenglish{Computer} \\ \hline
    المعالج & \textenglish{Processeur, UCT} & \textenglish{Processor, CPU} \\ \hline
    الوحدة المركزية & \textenglish{Unité centrale} & \textenglish{Central unit} \\ \hline
    الذاكرة المركزية & \textenglish{Mémoire centrale} & \textenglish{Main memory} \\ \hline
    ذاكرة متطايرة & \textenglish{Mémoire volatile} & \textenglish{Volatile memory} \\ \hline
    الحافلة & \textenglish{Bus} & \textenglish{Bus} \\ \hline
    حافلة العناوين & \textenglish{Bus d'adresses} & \textenglish{Address bus} \\ \hline
    حافلة المعطيات & \textenglish{Bus de données} & \textenglish{Data bus} \\ \hline
    حافلة التحكم & \textenglish{Bus de contrôle} & \textenglish{Control bus} \\ \hline
    السجل & \textenglish{Registre} & \textenglish{Register} \\ \hline
    العدّاد البرنامجي & \textenglish{Compteur ordinal (CO)} & \textenglish{Program counter (PC)} \\ \hline
    سجل التعليمة & \textenglish{Registre d'instruction (RI)} & \textenglish{Instruction register (IR)} \\ \hline
    سجل عنوان الذاكرة & \textenglish{Registre d'adresse mémoire (RAM)} & \textenglish{Memory address register (MAR)} \\ \hline
    سجل معطيات الذاكرة & \textenglish{Registre d'information mémoire (RIM)} & \textenglish{Memory data register (MDR)} \\ \hline
    المركِّم & \textenglish{Accumulateur} & \textenglish{Accumulator} \\ \hline
    العلم & \textenglish{Indicateur, drapeau} & \textenglish{Flag} \\ \hline
    وحدة الحساب والمنطق & \textenglish{Unité arithmétique et logique (UAL)} & \textenglish{Arithmetic logic unit (ALU)} \\ \hline
    وحدة التحكم & \textenglish{Unité de commande (UC)} & \textenglish{Control unit} \\ \hline
    المفكك & \textenglish{Décodeur} & \textenglish{Decoder} \\ \hline
    المتتالي & \textenglish{Séquenceur} & \textenglish{Sequencer} \\ \hline
    نصف المجمّع & \textenglish{Demi-additionneur} & \textenglish{Half adder} \\ \hline
    المجمّع الكامل & \textenglish{Additionneur complet} & \textenglish{Full adder} \\ \hline
    المتمم لاثنين & \textenglish{Complément à deux} & \textenglish{Two's complement} \\ \hline
    مجموعة التعليمات & \textenglish{Jeu d'instructions} & \textenglish{Instruction set} \\ \hline
    الرمز العملي & \textenglish{Code opération} & \textenglish{Opcode} \\ \hline
    نمط العنونة & \textenglish{Mode d'adressage} & \textenglish{Addressing mode} \\ \hline
    العنوان الفعلي & \textenglish{Adresse effective} & \textenglish{Effective address} \\ \hline
    العنوان الفيزيائي & \textenglish{Adresse physique} & \textenglish{Physical address} \\ \hline
    المقطع والإزاحة & \textenglish{Segment, déplacement} & \textenglish{Segment, offset} \\ \hline
    دورة التعليمة & \textenglish{Cycle d'instruction} & \textenglish{Instruction cycle} \\ \hline
\end{tabular}
\end{center}

\begin{center}
\small
\begin{tabular}{|l|l|l|}
    \hline
    \ras{بالعربية} & \ras{\textenglish{Français}} & \ras{\textenglish{English}} \\ \hline
    الجلب، فك الترميز، التنفيذ & \textenglish{Recherche, décodage, exécution} & \textenglish{Fetch, decode, execute} \\ \hline
    لغة التجميع & \textenglish{Langage d'assemblage} & \textenglish{Assembly language} \\ \hline
    المجمِّع & \textenglish{Assembleur} & \textenglish{Assembler} \\ \hline
    المكدس & \textenglish{Pile} & \textenglish{Stack} \\ \hline
    الإجراء & \textenglish{Procédure, sous-programme} & \textenglish{Procedure, subroutine} \\ \hline
    المقاطعة & \textenglish{Interruption} & \textenglish{Interrupt} \\ \hline
    جدول أشعة المقاطعات & \textenglish{Table des vecteurs d'interruption} & \textenglish{Interrupt vector table} \\ \hline
    المنفذ & \textenglish{Port d'E/S} & \textenglish{I/O port} \\ \hline
    النمط المحمي & \textenglish{Mode protégé} & \textenglish{Protected mode} \\ \hline
    النموذج المسطح & \textenglish{Modèle mémoire plat} & \textenglish{Flat memory model} \\ \hline
    خط الأنابيب & \textenglish{Pipeline} & \textenglish{Pipeline} \\ \hline
    التنبؤ بالقفز & \textenglish{Prédiction de branchement} & \textenglish{Branch prediction} \\ \hline
    النواة & \textenglish{Cœur} & \textenglish{Core} \\ \hline
    الخيط & \textenglish{Thread, fil d'exécution} & \textenglish{Thread} \\ \hline
    الذاكرة المخبئية & \textenglish{Mémoire cache, antémémoire} & \textenglish{Cache memory} \\ \hline
    سطر الذاكرة المخبئية & \textenglish{Ligne de cache} & \textenglish{Cache line} \\ \hline
    الإصابة / الإخفاق & \textenglish{Succès / Défaut} & \textenglish{Hit / Miss} \\ \hline
    المحلية الزمنية والمكانية & \textenglish{Localité temporelle et spatiale} & \textenglish{Temporal and spatial locality} \\ \hline
\end{tabular}
\end{center}

\section*{المراجع}
\begin{enumerate}
    \item \textenglish{A. S. Tanenbaum, \emph{Structured Computer Organization}, Pearson, 2012.}
    \item \textenglish{D. A. Patterson, J. L. Hennessy, \emph{Computer Organization and Design:}}\\ \textenglish{\emph{The Hardware/Software Interface}, Morgan Kaufmann.}
    \item \textenglish{P. Zanella, Y. Ligier, E. Lazard, \emph{Architecture et technologie des ordinateurs}, 6e édition, Dunod, 2018.}
    \item \textenglish{B. B. Brey, \emph{The Intel Microprocessors: 8086/8088 \dots{} Core2 with 64-bit Extensions}, Pearson.}
    \item \textenglish{EMU8086, Microprocessor Emulator.}
    \item \textenglish{NASM, The Netwide Assembler: https://www.nasm.us/}
    \item \textenglish{MASM, Microsoft Macro Assembler reference:}\\ \textenglish{\small https://learn.microsoft.com/en-us/cpp/assembler/masm/}
    \item \textenglish{GAS, The GNU Assembler: https://sourceware.org/binutils/docs/as/}
\end{enumerate}


\end{document}
